% concatenated from Latex_Pretty_good_plus_state_transfer_in_cycles.tex
\documentclass[12pt,a4paper]{article}


%%%%%%%%%%% Latex Packages %%%%%%%%%%%
\RequirePackage{etex}
\usepackage[utf8]{inputenc}
%\usepackage[margin=1in, footskip=0.25in]{geometry}
\usepackage{amsmath, amssymb, amsthm, amsfonts, mathrsfs}
\usepackage{bbm, commath}
\usepackage{multicol}
\usepackage{blindtext}
\usepackage{graphicx}
\usepackage{tikz-network}
\usetikzlibrary {positioning, patterns, calc}
\usetikzlibrary{graphs, graphs.standard, quotes}
\usepackage{authblk}
\usepackage{enumitem}
\usepackage[colorlinks=true,linkcolor=black,citecolor=black]{hyperref}
\usepackage{tabularx}
\usepackage{array}
\usepackage{subcaption}

%%Line numbers
%\usepackage{lineno}
%\linenumbers

% \usepackage{blindtext}
\usepackage{natbib}
\setlength{\bibsep}{1pt}
 \usepackage{setspace}
 \let\oldbibitem\bibitem
 \renewcommand{\bibitem}{\setstretch{1.2}\oldbibitem}
\renewcommand{\bibfont}{\small}


\textheight=24cm \textwidth=16cm \oddsidemargin=0.2in
\evensidemargin=0.2in \topmargin=-0.25in
\renewcommand{\baselinestretch}{1.5}

%%%%%%%%%%% Theorem, Definition, Lemma %%%%%%%%%%%


\newtheorem{thm}{Theorem}
\newtheorem{defn}{Definition}
\newtheorem{prop}{Proposition}
\newtheorem{lem}{Lemma}
\newtheorem{cor}{Corollary}
\newtheorem{exm}{Example}
\newtheorem{conj}{Conjecture}
\newtheorem{rem}{Remark}
%%%%%%%%%%% Re-define Commands %%%%%%%%%%%

\def \Fl {{\mathbbm F}}
\def \Nl {{\mathbbm N}}
\def \Zl {{\mathbbm Z}}
\def \Ql {{\mathbbm Q}}
\def \Rl {{\mathbbm R}}
\def \Cl {{\mathbbm C}}
\def \e {{\bf e}}
\def \u {{\bf u}}
\def \v {{\bf v}}
\def \w {{\bf w}} 
\def \x {{\bf x}}
\def \J {{\mathbb J}}
\def \l {{\mathbbm 1}}
\def \o {{\bf 0}}
\def \e {{\mathbf e}}
\def \ra {\rightarrow}
\def\mathbi #1{\text{\em #1}}

\newcommand{\ob}[1]{\left(#1\right)}
\newcommand{\cb}[1]{\left\lbrace #1\right\rbrace}
\newcommand{\tb}[1]{\left[#1\right]}
\newcommand{\mb}[1]{\left|#1\right|}
\newcommand{\up}[1]{\overset{#1}{\uplus}~}
\newcommand{\rk}[1]{\text{Rank}\left\lbrace#1\right\rbrace}
\newcommand{\eig}[1]{\sigma\left(#1\right)}
%\renewcommand{\figurename}{Figure}


%%%%%%%%%%% Title Inputs %%%%%%%%%%%
\title{\it Pretty good plus state transfer in cycles}
\author[1]{Sarojini Mohapatra}
\author[2]{Hiranmoy Pal}
\affil[1,2]{National Institute of Technology Rourkela, India-769008 palh@nitrkl.ac.in}
%\affil[2]{National Institute of Technology Rourkela, India-769008, 521ma6013@nitrkl.ac.in}
\date{\today}



%%%%%%%%%%% Begin Document %%%%%%%%%%%

\begin{document}
	
	\maketitle
	
	
	%%%%%%%%%%% Abstract %%%%%%%%%%%
	
		\begin{abstract}
        We investigate fractional revival in graphs with respect to the adjacency, Laplacian, and signless Laplacian matrices. We observe that, under certain conditions, fractional revival is preserved under graph complementation. 
       Then we establish a connection between fractional revival in a graph and in its double cover, and obtain a complete characterization of pretty good plus state transfer in cycles and their complements. This leads to characterizations of pretty good vertex state transfer in weighted paths with potential. 
\\\\
   
     \noindent{\it Keywords:}  Spectra of graphs, Circulant graph, Continuous-time quantum walk, Fractional revival, Pretty good plus state transfer.  
\\\\{\it MSC: 15A16, 05C50, 81P45.}
	\end{abstract}


   % \tableofcontents	
	
	%%%%%%%%%%% Introduction %%%%%%%%%%%
	\newpage

\section{Introduction} 
The transfer of quantum states in quantum spin networks is a fundamental problem in quantum information theory. 
\emph{Continuous-time quantum walks}, introduced by Farhi and Gutmann \cite{farhi}, provide a useful framework for analyzing such quantum transport phenomena.
% The use of quantum walk on a graph to analyze such information transfer was initiated by Bose \cite{bose}. 
A continuous-time quantum walk on a graph $G$ is governed  by the transition matrix 
\[U_{M(G)}(t)=\exp{(itM(G))}=\sum_{k\geq 0}\frac{(it)^k}{k!}(M(G))^k,\quad \text{where $t\in\mathbb{R}, $}\]
and M(G) denotes the Hamiltonian associated with $G.$ In this context, we consider $M(G)$ to be the adjacency matrix, the Laplacian matrix, or the signless Laplacian matrix of $G.$ Whenever the context is clear, we write 
$U(t)$ instead of 
$U_{M(G)}(t).$ Let $G (V(G), E(G), w)$ be a weighted undirected graph with vertex set $V(G)$ and edge set $E(G),$ where the weight function $w:E(G)\to \mathbb{R}$ assigns a real number to each edge of $G.$ The \textit{adjacency matrix} $A(G)$ is defined by $A(G)_{j,k}=w(j,k)$ for $(j,k)\in E(G),$ and $0$ otherwise. A \textit{potential on $G$} is a diagonal matrix $\Delta,$ where $\Delta_{j,j}$  is the potential at the vertex $j.$ For a graph $G$ with potential, the Hamiltonian $M(G)$ is considered as $A(G)+\Delta.$
% Let $G (V(G), E(G), w, \eta)$ be a weighted undirected graph with vertex set $V(G)$ and edge set $E(G),$ with the weight function $w:E(G)\to \mathbb{R}^+$ assigns a positive real weight to each edge of $G$ and potential $\eta: V(G)\to \mathbb{R}^+\cup\{0\}$ assigns a potential to each vertex of $G.$ The adjacency matrix $A(G)$ is defined by $A(G)_{j,k}=w(j,k)$ for $(j,k)\in E(G),$ and $A(G)_{j,j}=\eta(j)$, if there is a potential at vertex $j,$  and $0$ otherwise.
Unless otherwise specified, we assume that $G$ is finite, simple, and undirected. In the case of a simple graph, each edge has weight $1,$ and the potential $\Delta$ is the zero matrix. The \textit{Laplacian} and \textit{signless Laplacian} matrices of a simple graph are defined by $L(G)=D(G)-A(G),$ and $Q(G)=D(G)+A(G),$ respectively, where $D(G)$ is the degree matrix of $G.$  Let $\lambda_1,\lambda_2,\ldots,\lambda_d$ be the distinct eigenvalues of $M(G)$ with corresponding eigenprojection matrices $E_{\lambda_1},E_{\lambda_2},\ldots,E_{\lambda_d}.$ Then the spectral decomposition of $U_{M(G)}(t)$ is given by
   \[U_{M(G)}(t)=\sum_{j=1}^{d}\exp{(it\lambda_j)}E_{\lambda_j}.\]
   % The expression shows that the eigenvalue and eigenvectors of $M(G)$ completely determine the behavior of the quantum walk.
% If the graph is regular, then the transition matrices governed by the adjacency matrix, the Laplacian matrix, or the signless Laplacian matrix differ only by a global phase; consequently, the state transfer
% properties remain invariant under all such choices of $M(G).$ 

A \textit{real pure state} is represented by a unit vector in $\mathbb{R}^n$ \cite{god7}. Let $a$ and $b$ be two vertices in a graph $G.$ The characteristic vector $\e_a$ is called the \textit{vertex state} associated with $a.$ For a non-zero real number $s,$ a real pure state of the form $\frac{1}{\sqrt{1+s^2}}\ob{\e_a+s\e_b}$ is called an \textit{$s$-pair state} \cite{kim}. In particular, when $s=-1,$ the state is called a \textit{pair state} and for $s=1,$ it is called a \textit{plus state}. 
The \textit{eigenvalue support} of a state $\u$ relative to $M(G)$ is the set 
$\{\lambda_j:E_{\lambda_j}\u\neq 0\}.$
A state $\u$ is called a fixed state if and only if $\u$ is an eigenvector of $M(G)$ associated with the lone eigenvalue in the eigenvalue support of $\u$ \cite{god7}.
A graph $G$ is said to exhibit \emph{perfect state transfer (PST)} between two linearly independent real pure states $\u$ and $\v$ if there exists a
time $\tau$ and a phase factor $\gamma\in\mathbb{C}$ with $|\gamma|=1$ such that
\begin{equation}\label{5eq1}
  U_{M(G)}(\tau)\u=\gamma\v.  
\end{equation} 
If $\u=\v,$ then the state 
$\u$ is said to be periodic in 
$G.$
When $\u$ and $\v$ in \eqref{5eq1} are both vertex states, pair states, or plus states, then the corresponding PST is called \textit{vertex PST, pair PST}, or \textit{plus PST}, respectively.  PST in quantum spin networks was first introduced by Bose \cite{bose} and has attracted considerable attention over the past two decades, with extensive research devoted to characterizing vertex PST.
% Following the work of Bose \cite{bose}, the study of state transfer in graphs has attracted considerable attention, and over the past two decades, extensive research has been devoted to characterizing vertex PST. 
% It has been characterized for several families of graphs, including paths $P_n$ \cite{chr2, god1}, circulant graphs \cite {ange1, bas1, mil4, bas3, bas2}, Cayley graphs \cite{ber, che, pal2, tan}, trees \cite{cou7, cou0}, join and product of graphs \cite{cou, kirk3}, corona graphs \cite{ack, ack1}, blow-up graphs \cite{pal9, mon3}. 
A fundamental result due to Godsil \cite{god2} establishes that for any integer $k>0,$ there is a finite number of connected graphs of maximum degree $k$ with vertex PST. This result has prompted the investigation of more general notions of state transfer beyond vertex states, 
% Motivated by this result there is an emerging literature on extension of vertex-PST,
such as pair (plus) PST \cite{che1,  jia, ojha1, pal10}, $s$-pair PST \cite{kim}, PST between real pure states \cite{god8, god7,pal11}.  
% As vertex state transfer, plus PST is also a rare phenomenon as discussed in \cite[Theorem 3.8]{god7}.

A relaxation to vertex PST known as \emph{pretty good state transfer (PGST)} was introduced in \cite{god1, vin}. Extension of this notion to more general states, including pair states and real pure states, have also been investigated \cite{god8, pal11, ojha1, pal10}. A graph  $G$ is said to exhibit \emph{PGST} between two linearly independent real pure states $\u$ and $\v$ if there exists a sequence of real number $t_k\in\mathbb{R},$ and a complex number $\gamma$ of unit modulus such that 
\begin{equation}\label{5eq2}
       \lim_{k \to \infty}U_{M(G)}\ob{t_k}\u=\gamma \v.
   \end{equation}
   As in vertex PST, \cite{god7} shows that plus PST with respect to the adjacency, Laplacian, and signless Laplacian matrices is a rare phenomenon. This motivates us to investigate plus PGST in graphs. In particular, we study plus PGST in cycles and their complements. 
   
   
   % Furthermore, there are relatively few results on \textit{pretty good plus state transfer (plus PGST)}, which motivates us to investigate plus PGST. 
   % motivates us
   % As vertex state transfer, plus PST is also a rare phenomenon as discussed in \cite[Theorem 3.8]{god7}.
   % In this article, we mainly focus on pretty good plus state transfer. 
%    A graph  $G$ is said to exhibit \textit{plus PGST} between two linearly independent states $\frac{1}{\sqrt{2}}\ob{\e_a+\e_b}$ and $\frac{1}{\sqrt{2}}\ob{\e_c+\e_d}$ if there exists a sequence of real number $t_k\in\mathbb{R},$ and a complex number $\gamma$ of unit modulus such that 
% \begin{equation}\label{5eq2}
%        \lim_{k \to \infty}U_{M(G)}\ob{t_k}\ob{\e_a+\e_b}=\gamma\ob{\e_c+\e_d}.
%    \end{equation}
   % Some recent work investigating pair PGST includes characterization of pair PGST in paths and cycles \cite{}, corona graphs \cite{}. There is not much literature on the study of plus PGST. This motivates us to investigate on plus PGST. 

\emph{Fractional revival (FR)} is another generalization of PST, which is relevant for quantum entanglement generation. FR occurs whenever a continuous-time quantum walk maps
the characteristic vector of a vertex to a superposition of the characteristic vectors of a
subset of vertices containing the initial vertex.
% A main focus will be on the case when the
% subset has two vertices.
% In quantum information, it is known that entanglement is a useful and important resource. 
% A relevant quantum transport
% phenomenon for entanglement generation is fractional revival. 
 % We say fractional revival occurs at a and b at time τ if
% e−iτH|a⟩ = α|a⟩ + β|b⟩ (7)
% for complex scalars α and β ̸= 0 with |α|2 + |β|2 = 1.A graph $G$ is said to have fractional revival from a vertex $a$ whenever the transition matrix of $G$ maps the state associated with $a$ to a linear combination of the states of a collection of vertices containing the initial one \cite{bern1, cha2, gan1, mon1}. 
% A graph $G$ is said to exhibit \textit{FR} at $a$ and $b$ at time $t$ if 
%  \[U_{M(G)}(t)\e_a=\alpha\e_a+\beta\e_b,\]
%  for some complex scalars $\alpha$ and $\beta\neq 0$ with $|\alpha|^2+|\beta|^2=1.$ 
 In a path $P_n$ on $n$ vertices, FR occurs between two vertices with respect to the adjacency matrix if and only if $n\in\{2,3,4\},$ while among cycles only $C_4$ and $C_6$ admit FR \cite{chan2}. Furthermore, no tree admits Laplacian FR except for the paths on two
and three vertices \cite{chan3}.
In this work, we consider the initial state to be a real pure state. Let 
$\u$ and $\v$ be two linearly independent states in a graph $G.$ The graph $G$ is said to exhibit \emph{FR} from $\u$ to $\v$
 at time $\tau$ if there exist complex scalars $\alpha,\beta$ with $\beta\neq 0$
such that
\begin{equation}\label{5e1}
U_{M(G)}(\tau)\u=\alpha\u+\beta\v.
\end{equation}
% This is equivalent to the notion of generalized fractional revival.
In this case, we also say that $(\alpha,\beta)$-FR occurs from $\u$ to $\v$ at $\tau.$ 
In particular, if $\u$ and $\v$ both represent pair (plus) states, then the graph is said to have \textit{pair (plus) FR}. If $\u$ and $\v$ in \eqref{5e1} are orthogonal states, then $|\alpha|^2+|\beta|^2=1.$
% The investigation of fractional revival in graphs with the initial state as a vertex state includes characterizations of paths and cycles \cite{chan2}, Cayley graphs \cite{}, trees \cite{chan3}, threshold graphs \cite{}, etc. 
% For the path graph $P_n$, FR occurs between two vertices with respect to the adjacency matrix if and only if $n\in\{2,3,4\},$ while among cycles only $C_4$ and $C_6$ admit FR \cite{chan2}. Furthermore, no
% tree admits Laplacian FR except for the paths on two
% and three vertices \cite{chan3}.
If $\alpha=0$ in \eqref{5e1}, then $G$ admits PST between $\u$ and $\v.$  A non-trivial relation between FR from a vertex state and a pair state is established in \cite{pal10}.
A relaxation to FR called as \textit{pretty good fractional revival (PGFR)} between vertex states was introduced by Chan et al. in \cite{cha3}. 
% A graph $G$ is said to exhibit \textit{PGFR} from vertex $a$ to $b$ if there is a sequence of time $t_k\in\mathbb{R}$ such that 
%    \[
%        \lim_{k \to \infty}U_{M(G)}\ob{t_k}\e_a=\alpha\e_a+\beta\e_b,
%   \]
%   for some complex scalars $\alpha$ and $\beta\neq 0$ with $|\alpha|^2+|\beta|^2=1.$
  We consider PGFR from a real pure state $\u$ to $\v,$ where $\u$ and $\v$ are linearly independent.
   A graph $G$ is said to exhibit PGFR from $\u$ to $\v$ if there is a sequence of time $t_k\in\mathbb{R}$ such that for some $\alpha,\beta\in\mathbb{C}$ with $\beta\neq 0,$
   \[
       \lim_{k \to \infty}U_{M(G)}\ob{t_k}\u=\alpha\u+\beta\v.
  \]
% In particular, if $\u$ and $\v$ represent pair states $\frac{1}{\sqrt{2}}\ob{\e_a-\e_b}$ and $\frac{1}{\sqrt{2}}\ob{\e_c-\e_d},$ respectively, then we have fractional revival between $\frac{1}{\sqrt{2}}\ob{\e_a-\e_b}$ and $\frac{1}{\sqrt{2}}\ob{\e_c-\e_d}$ at $\tau.$ 

 
% In this article we consider fractional revival from a real pure state.  
% A graph $G$ is said to exhibit $(\alpha,\beta)$- pretty good fractional revival (PGFR) from $\u$ to $\v$ if there is a sequence of time $t_k\in\mathbb{R}$ such that 
%    \begin{equation}\label{1e1}
%        \lim_{k \to \infty}U_G\ob{t_k}\u=\alpha\u+\beta\v.
%    \end{equation}



\begin{rem}\label{5r1}
    If $G$ is a regular graph, then the transition matrices governed by the adjacency matrix, the Laplacian matrix, or the signless Laplacian matrix differ only by a global phase. Consequently, the state transfer
properties are identical under all such choices of $M(G).$ 
\end{rem}

   

% We begin by outlining a few preliminary results that will be useful in the subsequent sections. Let $\lambda_1,\lambda_2,\ldots,\lambda_d$ be the distinct eigenvalues of $M(G)$ with corresponding eigenprojection matrices $E_{\lambda_1},E_{\lambda_2},\ldots,E_{\lambda_d}.$ The spectral decomposition of $U_{M(G)}(t)$ is given by
%    \[U_{M(G)}(t)=\sum_{j=1}^{d}\exp{(it\lambda_j)}E_{\lambda_j}.\]
%    Consequently, the eigenvalue and eigenvectors of $M(G)$ completely determine the behavior of the quantum walk. The \textit{eigenvalue support} of a state $\u$ relative to $M(G),$ denoted by $\sigma_\u(M(G)),$ is the set 
%    \[\sigma_\u(M(G))=\{\lambda_j:E_{\lambda_j}\u\neq 0\}.\]

% A vector $\u$ is called a fixed state if and only if $\u$ is an eigenvector for $M(G)$ associated with the lone eigenvalue in $\sigma_{\u}(M(G))$ \cite{god7}. 
% The structure of a graph allows us to simplify the analysis of quantum walks using equitable partitions. A partition $\Pi$ of the vertex set of a wighted graph $G$ with disjoint cells $V_1, V_2, \ldots, V_d$ is said to be \textit{equitable} if
% \[\sum_{b \in V_k} w(a,b)=c_{jk}~~~\text{for each}~ a\in V_j,\]
% where $c_{jk}\in \Rl$ is a constant depending only on the cells $V_j$ and $V_k$. The \textit{characteristic matrix} $\mathcal{C}$ associated to $\Pi$ is an $n \times d$ matrix where the columns represent the normalized characteristic vectors of $V_1, V_2,\dots, V_d.$ The graph with the $d$ cells of $\Pi$ as its vertices having adjacency matrix $A(G/\Pi),$ where $A(G/\Pi)_{j,k}=\sqrt{c_{jk}c_{kj}},$ is called the \textit{symmetrized quotient graph} of $G$ over $\Pi,$ and is denoted by $G/\Pi.$ The transition matrices of $G$ and $G/\Pi$ are related as follows.

% \begin{prop}\cite{bac, god0}\label{5p1}
%  Let $G$ be a weighted graph with adjacency matrix $A(G).$ If $\Pi$ is an equitable partition of $G$ having characteristic matrix $\mathcal{C}$ then $A(G)\mathcal{C}=\mathcal{C}A(G/\Pi),$ where $A(G/\Pi)$ is the adjacency matrix of the symmetrized quotient graph $G/\Pi.$ Moreover, $U_{A(G)}(t)\mathcal{C}=\mathcal{C}U_{A(G/\Pi)}(t),$ for all $t\in\mathbb{R}.$
% \end{prop}
Throughout the paper, $\mathbf{1}$ denotes the all-ones vector, $I$ the identity matrix, $J$ the all-ones square matrix, and $\mathbf{0}$ the zero matrix of appropriate order. 
% The spectrum of $M(G)$ is denoted by 
% $\cb{\lambda_1^{(m_1)},\lambda_2^{(m_2)},\ldots,\lambda_d^{(m_d)}},$
% where $(m_j)$ indicates the multiplicity of the eigenvalue $\lambda_j.$ 
The article is organized as follows. In Section \ref{5s2}, we include the effect of certain graph automorphisms on the existence of PGST from plus and vertex states. In Section \ref{5sec2}, we study FR from a real pure state in the complement of a graph with respect to the adjacency, Laplacian, and signless Laplacian matrices. Section \ref{5sec3} establishes a relation between FR in a graph and its double cover. In Section \ref{5sec5}, we provide a complete characterization of plus PGST in cycles and their complements. A few observations on vertex PGST in certain weighted paths with potential is presented in Section \ref{5sec6}. 
% For a regular graph, the evolution of a state 
% $\u$ satisfying $\mathbf{1}^T\u=0$ is preserved under graph complementation for 
% $M\in\{A,Q\}$ as shown in \cite[Lemma 3]{pal11}. We further observe that the graph regularity is not necessary for the preservation of FR from a state $\u$ relative to the adjacency and signless Laplacian matrices when $\u$ is orthogonal to $\mathbf{1}.$ 







%  A graph $G$ is said to have fractional revival from a vertex $a$ whenever the transition matrix of $G$ maps the state associated with $a$ to a linear combination of the states of a collection of vertices containing the initial one \cite{bern1, cha2, gan1, mon1}. A graph $G$ exhibits fractional revival at $a$ and $b$ at time $\tau$ if 
%  \[U_G(\tau)\e_a=\alpha\e_a+\beta\e_b,\]
%  for some complex scalars $\alpha$ and $\beta\neq 0$ with $|\alpha|^2+|\beta|^2=1.$ We say that fractional revival occurs from a unit vector $\u\in\mathbb{C}^{|V(G)|}$
%  at time $\tau$ if there exist complex scalars $\alpha,\beta$ with $\beta\neq 0$
% and a unit vector $\v\in\mathbb{C}^{|V(G)|}$
% such that
% \begin{equation}\label{5e1}
% U_G(\tau)\u=\alpha\u+\beta\v.
% \end{equation}
% In particular, if $\u$ and $\v$ represent pair states $\frac{1}{\sqrt{2}}\ob{\e_a-\e_b}$ and $\frac{1}{\sqrt{2}}\ob{\e_c-\e_d},$ respectively, then we have fractional revival between $\frac{1}{\sqrt{2}}\ob{\e_a-\e_b}$ and $\frac{1}{\sqrt{2}}\ob{\e_c-\e_d}$ at $\tau.$ If $\alpha$ becomes zero, $G$ is said to exhibit perfect state transfer between real pure states $\u$ and $\v$ at $\tau.$ 
% A graph $G$ is said to exhibit $(\alpha,\beta)$- pretty good fractional revival (PGFR) between $\frac{1}{\sqrt{2}}\ob{\e_a-\e_b}$ and $\frac{1}{\sqrt{2}}\ob{\e_c-\e_d}$ if there is a sequence of time $t_k\in\mathbb{R}$ such that 
%    \begin{equation}\label{1e1}
%        \lim_{k \to \infty}U_G\ob{t_k}\ob{\e_a-\e_b}=\alpha\ob{\e_a-\e_b}+\beta\ob{\e_c-\e_d}.
%    \end{equation}

   % Let $M$ be a Hamiltonian of a graph $G.$ Since $M$ is real symmetric, we may write $M$ in its spectral decomposition 
   % \[M=\sum_{j=1}^{d}\lambda_jE_{\lambda_j},\]
   % where $\lambda_1,\lambda_2,\ldots,\lambda_d$ are the distinct eigenvalues of $M$ with corresponding eigenprojection matrices $E_1,E_2,\ldots,E_d.$ The spectral decomposition of $U_G(t)$ given by
   % \[U_G(t)=\sum_{j=1}^{d}\exp{(it\lambda_j)}E_{\lambda_j}.\]
   % Thus, the eigenvalue and eigenvectors of $M$ completely determine the behavior of the quantum walk. Let $x\in\mathbb{R}^n\backslash\{0\}.$ The eigenvalue support of $x$ relative to $M,$ denoted $\sigma_x(M),$ is the set 
   % \[\sigma_x(M)=\{\lambda_j:E_{\lambda_j}x\neq 0\}.\]
% \section{Preliminaries}
%  Let $J$ denote the all-ones matrix and $I$ be the identity matrix. define  $\mathbf{1}_n.$ Let $M(G)$ denote the adjacency, Laplacian, or signless Laplacian matrix of $G.$ Let $M(G)$ have $d$ distinct eigenvalues. We denote the spectrum of $M(G)$ by 
% \[\text{spec}(M(G))=\{\lambda_1^{[m_1]},\lambda_2^{[m_2]},\ldots,\lambda_d^{[m_d]}\},\] 
% where $[m_j]$ indicates the multiplicity of the eigenvalue $\lambda_j.$ 
% \begin{prop}\cite{god7}
%     A vector $x\in\mathbb{R}^n\backslash\{0\}$ is a fixed state if and only if $x$ is an eigenvector for $M$ associated with the lone eigenvalue in $\sigma_x(M).$
% \end{prop}
% The above proposition demonstrates that a fixed state is not involved in perfect state transfer.

\section{Algebraic properties}\label{5s2}
An \emph{automorphism} $f$ of a graph $G$ is a bijection on the vertex set $V(G)$ such that vertices 
$a$ and $b$ are adjacent in $G$ if and only if  $f(a)$ and $f(b)$ are adjacent in $G.$
 If $P$ is the permutation matrix corresponding to the automorphism $f,$ then $P$ commutes with $M(G).$  Since the transition matrix $U_{M(G)}(t)$ is a polynomoal in $M(G),$ the matrix $P$ commutes with $U_{M(G)}(t)$ as well. The automorphism $f$ is said to \emph{fix} a vertex $a$ in $G$ if $P\e_a=\e_a.$ 
 In the following, we observe that the existence of certain automorphisms in a graph $G$ with plus PGST guarantees the existence of pair PGST in $G.$
\begin{thm}\label{5l7}
 Let a graph $G$ admit pretty good plus state transfer between $\frac{1}{\sqrt{2}}\ob{\e_a+\e_b}$ and $\frac{1}{\sqrt{2}}\ob{\e_c+\e_d}.$ Then the graph $G$ admits pretty good pair state transfer if there exists an automorphism of $G$ with  permutation matrix 
 $P$ satisfying one of the following conditions.
 \begin{enumerate}
    \item $P\e_a=\e_a$ and $P\e_b\neq \e_b,$ 
    \item  $P\e_a=\e_b$ and $P\e_b\neq \e_a.$ 
\end{enumerate}
 \end{thm}
 \begin{proof}
  Suppose the graph $G$ admits plus PGST between $\frac{1}{\sqrt{2}}\ob{\e_a+\e_b}$ and $\frac{1}{\sqrt{2}}\ob{\e_c+\e_d}.$ Then by \eqref{5eq2}, we have 
  \begin{equation}\label{5eqn7}
         \lim_{k \to \infty}U_{M(G)}\ob{t_k}\ob{\e_a+\e_b}=\gamma\ob{\e_c+\e_d}.
     \end{equation}
    Suppose condition-1 holds. Then
     multiplying $P$ on both sides of \eqref{5eqn7} and subtracting the resulting equation from \eqref{5eqn7}, yields
 \begin{equation}\label{5eqn8}
        \lim_{k \to \infty}U_{M(G)}\ob{t_k}\ob{\e_b-P\e_b}=\gamma\ob{\e_c+\e_d-P\e_c-P\e_d}.
   \end{equation}
    Since $U_{M(G)}(t)$ is unitary, \eqref{5eqn8} implies that $\left\|\frac{1}{\sqrt{2}}\ob{\e_b-P\e_b}\right\|=\left\|\frac{1}{\sqrt{2}}\ob{\e_c+\e_d-P\e_c-P\e_d}\right\|.$ This equality is possible only if $\frac{1}{\sqrt{2}}\ob{\e_c+\e_d-P\e_c-P\e_d}$ represents a pair state. Hence, $G$ admits pair PGST.  
    Similarly, if condition-2 holds, an analogous argument shows that  
    $G$ admits pair PGST.
 \end{proof}
 Next, we observe that certain graph automorphisms prevent the existence of PGST between a vertex state and a plus state.
 \begin{thm}\label{5th1}
 Let $a, b$ and $c$ be vertices of a graph $G.$ If there exists an automorphism of $G$ that fixes $b$ but not $a,$ then there is no pretty good state transfer between 
$\e_a$ and $\frac{1}{\sqrt{2}}\ob{\e_b+\e_c}.$ 
\end{thm} 
%  If there exists an automorphism of 
% $G$ that fixes 
% $b$
% but not 
% $a,$ then there is no pretty good state transfer between 
% $\e_a$ and $\frac{1}{\sqrt{2}}\ob{\e_b+\e_c}.$
\begin{proof}
    Suppose $G$ admits PGST between $\e_a$ and $\frac{1}{\sqrt{2}}\ob{\e_b+\e_c}.$  Then by \eqref{5eq2}, we have
    \begin{equation}\label{5e29}
        \lim_{k \to \infty}U_{M(G)}\ob{t_k}\e_a=\frac{\gamma}{\sqrt{2}}\ob{\e_b+\e_c}.
    \end{equation}
     Let $P$ be the permutation matrix of the automorphism of $G$ that fixes $b$ but not $a.$  Multiplying $P$ on both sides of \eqref{5e29} and subtracting the resulting equation from \eqref{5e29}, gives
 \[
        \lim_{k \to \infty}U_{M(G)}\ob{t_k}(\e_a-P\e_a)=\frac{\gamma}{\sqrt{2}}\ob{\e_c-P\e_c},
    \]
 which is a contradiction as $\left\|\e_a-P\e_a\right\|\neq \left\|\frac{1}{\sqrt{2}}\ob{\e_c-P\e_c}\right\|.$ 
Hence, the result follows.
\end{proof}

\section{State transfer in graph complement}\label{5sec2}
The \emph{complement} of a graph $G,$ denoted by $\overline{G},$ is the graph on the vertex set of $G,$ where two distinct vertices are adjacent in $\overline {G}$ if and only if they are not adjacent in $G.$ 
The matrix $M(\overline{G})$ corresponding to the adjacency, Laplacian, and signless Laplacian matrices of the complement $\overline{G}$ on $n$ vertices is given by
\[M(\overline{G})=\delta J+\zeta I-M(G),\]
   where $\delta=\left\{ \begin{array}{rcl}
 -1, & \mbox{if} & M =L \\ 
 1, & \mbox{if} & M \in\{A,Q\},
 \end{array}\right.$
\quad and \quad
$\zeta=\left\{ \begin{array}{rcl}
 -1, & \mbox{if} & M=A, \\ 
n, & \mbox{if} & M=L, \\
n-2, & \mbox{if} & M=Q. 
 \end{array}\right.$
% \\The following result shows the relation between the transition matrices of $G$ and its complement $\overline{G}.$
\\The preservation of vertex PST under graph complementation was established in \cite[Lemma 15.3]{god1} and \cite[Theorem 2]{alv} for the adjacency matrix of regular graphs and for the Laplacian matrix, respectively. Moreover, the preservation of Laplacian FR between two vertices under graph complementation was shown in \cite[Theorem 9]{chan3}. The observations extend to FR from real pure states as well. We include the proof for convenience.
 \begin{thm}\label{5t15}
 Let $\mathbf{1}$ be an eigenvector of a graph $G$ on $n$ vertices. If $G$ exhibits fractional revival from a state $\u$ to $\v$ at time $\tau$ with respect to $M(G)$ and $n\tau\in 2\pi\mathbb{Z},$ then the complement $\overline G$ exhibits fractional revival from $\u$ to $\v$ relative to $M(\overline{G}).$ 
 \end{thm}
\begin{proof}
   Since $\mathbf{1}$ is an eigenvector of $G$ and $J=\mathbf{1}\mathbf{1}^T,$ it follows that $M(G)$ commutes with $J.$ Therefore, for all $t\in\mathbb{R},$
   \[U_{M(\overline{G})}(t)=\exp{(itM(\overline{G}))=\exp{(i t\zeta)}\exp{(it\delta J)}}\exp{(-itM(G))}.\]
    The spectral decomposition of $J$ gives
   $\exp{(it\delta J)}=[\exp{(itn\delta)}-1]\frac{1}{n}J+I.$ 
   If $n\tau\in 2\pi \mathbb{Z},$ then we obtain $U_{M(\overline{G})}(-\tau)\u=\exp{(-i\tau\zeta)} U_{M(G)}(\tau)\u,$ and the result follows.
\end{proof}
%  \begin{lem}\label{5l5}
%  Let $\mathbf{1}_n$ be an eigenvector of a graph $G$ on $n$ vertices. Suppose $U_{M(G)}(t)$ and $U_{M(\overline{G})}(t)$ be the transition matrices of $G$ and its complement $\overline{G},$ respectively. If $nt\in 2\pi \mathbb{Z},$ then for any vector $\u,$
%  \[U_{M(\overline{G})}(t)\u=\exp{(it\zeta)} U_{M(G)}(-t)\u,\]
%  where $\zeta$ is as defined earlier for $M\in\{A,L,Q\}.$
%  \end{lem}
% \begin{proof}
%    Since $\mathbf{1}_n$ is an eigenvector of $G,$ it follows that $M(G)$ commutes with $J.$ Therefore,
%    \[U_{M(\overline{G})}(t)=\exp{(itM(\overline{G}))=\exp{(i t\zeta)}\exp{(it\delta J)}}\exp{(-itM(G))}.\]
%     The spectral decomposition of $J$ gives
%    \[\exp{(it\delta J)}=[\exp{(itn\delta)}-1]\frac{1}{n}J+I.\]
%    If $nt\in 2\pi \mathbb{Z},$ then we obtain $U_{M(\overline{G})}(t)=\exp{(it\zeta)} U_{M(G)}(-t).$ Consequently, for any vector $\u,$ the result follows.
% \end{proof}
% The preservation of vertex PST under graph complementation was established in \cite[Theorem 2]{alv} and \cite[Lemma 15.3]{god1} for the Laplacian matrix, and for the adjacency matrix in the case of regular graphs, respectively. Moreover, the preservation of Laplacian FR between two vertices under graph complementation was observed in \cite[Theorem 9]{chan3}. Here, we generalize this phenomenon to FR from real pure states.
% \begin{thm}\label{5t15}
%   Let $G$ exhibit fractional revival from a state $\u$ to $\v$ at time $\tau$ with respect to the Laplacian matrix. If $|V(G)|\tau\in 2\pi\mathbb{Z},$ then the complement $\overline G$ exhibits fractional revival from $\u$ to $\v$ relative to the Laplacian matrix. 
% \end{thm}
 It is well known that the path $P_2$ exhibits PST between the end vertices at $\frac{\pi}{2}.$ Using Theorem \ref{5t15}, one may observe that the complete graph on $4n$ vertices with a single missing edge admits Laplacian vertex PST between the non-adjacent vertices. Furthermore, the result in \cite[Corollary 1]{pal8} can be recovered as a direct consequence of Theorem \ref{5t15}.

% \begin{rem}
%  The observation in Theorem \ref{5t15} can be extended to PGFR. More precisely, if the graph $G$ exhibits PGFR from a state $\u$ to $\v$ with respect to the sequence $t_k\in 2\pi\mathbb{Z},$ relative to the Laplacian matrix $L(G),$ then the the complement $\overline G$ also exhibits PGFR from $\u$ to $\v$ relative to $L(\overline{G}).$ 
%  % Remark \ref{5r1} implies that Theorem \ref{5t15} remains valid for a regular graph when the Laplacian matrix is replaced by the adjacency matrix or the signless Laplacian matrix. 
%  Note that if $\mathbf{1}$ is an eigenvector of $G$ and the state $\u$ satisfies $\mathbf{1}^T\u=0,$ then the preservation of state transfer under graph complementation was established in \cite{pal11}. In this case, no additional condition on the time parameter is required. \end{rem}

%  We further observe that if $\u$ is a state in $G$ such that $\mathbf{1}^T\u=0,$ then  
%  \[M(\overline{G})\u=(\zeta I-M(G))\u.\] Consequently, the following result demonstrates the preservation of state transfer under graph complementation for all three choices of $M(G)$, without requiring 
% $G$ to be regular.
% \begin{lem}
%   Let $U_{M(G)}(t)$ and $U_{M(\overline{G})}(t)$ be the transition matrices of $G$ and its complement $\overline{G},$ respectively. If $\u$ is a vector with $\mathbf{1}^T\u=0,$ then
%  \[U_{M(\overline{G})}(t)\u=\exp{(it\zeta)} U_{M(G)}(-t)\u,\]
%  where $\zeta$ is as defined earlier for $M\in\{A,L,Q\}.$
%  \end{lem}
% the following result is obtained as an application of Theorem \ref{5t15} analogous to the observation in \cite[Corollary 1]{alv} and \cite[Corollary 10]{chan3}. 
The \emph{join} of two graphs $G$ and $H$ is defined by $G+H:=\overline{\overline{G}\cup\overline{H}},$ where $\overline{G}\cup\overline{H}$ is the disjoint union of $\overline{G}$ and $\overline{H}.$ An analogous observation to \cite[Corollary 10]{chan3} is obtained for real pure states using Theorem \ref{5t15}. 
 \begin{cor}\label{5c3}
 Let the complement of a graph $G$ exhibit fractional revival from $\u$ to $\v$ at time $\tau$ relative to the Laplacian matrix. For any graph $H,$ if $\tau(|V(G)|+|V(H)|)\in 2\pi\mathbb{Z},$
 then the join $G + H$ has fractional revival from $\u$ to $\v$ relative to the Laplacian matrix.
 \end{cor}
% The \emph{join} of two graphs $G$ and $H$ is defined by $G+H:=\overline{\overline{G}\cup\overline{H}}.$
 Although the path $P_3$ does not admit Laplacian plus PST \cite[Corollary 7.8]{god7}, we use it to find infinitely many graphs with Laplacian plus PST. 
\begin{exm}\label{5ex1}
% Although the path $P_3$ does not admit Laplacian plus PST \cite[Corollary 7.8]{god7}, we use it to find infinitely many graphs with Laplacian plus PST.
Let $P_3$ be a path on three vertices $1,2$ and $3,$ where both $1$ and $3$ are adjacent to $2.$ Since, the path $P_2$ admits Laplacian PST between the end vertices at $\frac{\pi}{2},$ the complement $\overline{P_3}$ admits PST between $\frac{1}{\sqrt{2}}\ob{\e_1+\e_2}$ and $\frac{1}{\sqrt{2}}\ob{\e_3+\e_2}$  at $\frac{\pi}{2}.$ Therefore, by Corollary \ref{5c3},
the graph $P_3+H$ admits Laplacian plus PST at $\frac{\pi}{2}$ between $\frac{1}{\sqrt{2}}\ob{\e_1+\e_2}$ and $\frac{1}{\sqrt{2}}\ob{\e_3+\e_2},$ whenever $|V(H)|\equiv 1 \pmod 4.$ 
\end{exm}
% \begin{proof}
    
% \end{proof}

%A complete multipartite graph, denoted by $K_{n_1,n_2,\ldots,n_k}$ is a simple graph whose vertex set can be partitioned into $k\geq 2$ pairwise disjoint nonempty subsets $V_1, V_2, \ldots, V_k$ such that for all distinct $i,j\in\{1,2,\ldots,k\},$ every vertex in $V_i$ is adjacent to every vertex in $V_j,$ and no two vertices within the same subset $V_i$ are adjacent, where $n_i=|V_i|$ for $1\leq i\leq k.$ 
% The complete multipartite graph $K_{4,n_2,\ldots,n_l}$ can be realized as $\overline{K_4\cup K_{n_2}\cup\cdots \cup K_{n_l}}.$ Since the complete graph $K_4$ has plus PST between $\frac{1}{\sqrt{2}}(\e_1+\e_3)$ and $\frac{1}{\sqrt{2}}(\e_2+\e_4)$ at $\frac{\pi}{4}$ \cite[Example 6.2]{god7}, the following observation hold by Theorem \ref{5t15}.  

% \begin{exm}
% The complete multipartite graph $K_{4,n_2,\ldots,n_l}$ can be realized as $\overline{K_4\cup K_{n_2}\cup\cdots \cup K_{n_l}}.$ Since the complete graph $K_4$ has plus PST between $\frac{1}{\sqrt{2}}(\e_1+\e_3)$ and $\frac{1}{\sqrt{2}}(\e_2+\e_4)$ at $\frac{\pi}{4}$ \cite[Example 6.2]{god7}, the following observation hold by Theorem 
%     The complete multipartite graph $K_{n_1,n_2,\ldots,n_l}$ exhibits Laplacian perfect plus state transfer if $n_1=4$ and $\displaystyle\sum_{j=2}^l n_j\equiv 4\pmod 8.$ 
% \end{exm}
\begin{rem}\label{5remark2}
 The observation in Theorem \ref{5t15} can be extended to PGFR. More precisely, if a graph $G$ on $n$ vertices exhibits PGFR from a state $\u$ to $\v$ with respect to the sequence $t_k$ relative to the Laplacian matrix $L(G),$ and $nt_k\in 2\pi\mathbb{Z},$ then the the complement $\overline G$ also exhibits PGFR from $\u$ to $\v$ relative to $L(\overline{G}).$ 
 \end{rem}
 % Remark \ref{5r1} implies that Theorem \ref{5t15} remains valid for a regular graph when the Laplacian matrix is replaced by the adjacency matrix or the signless Laplacian matrix. 
 % Note that if $\mathbf{1}$ is an eigenvector of $G$ and the state $\u$ satisfies $\mathbf{1}^T\u=0,$ then the preservation of state transfer under graph complementation was established in \cite{pal11}. In this case, no additional condition on the time parameter is required.
For regular graphs, \cite[Lemma 3]{pal11} shows that state transfer is preserved under graph complementation for any real pure state $\u$ with $\mathbf{1}^T\u=0.$ 
% This does not require any additional condition on the time parameter. 
 % If $G$ is not regular, and a state $\u$ in $G$ satisfies $M(G)^2\u=\u$ with both $\u$ and $M(G)\u$ are orthogonal to $\mathbf{1},$ then   
 % \[M(\overline{G})\u=(\zeta I-M(G))\u.\] Consequently, the following result
 We provide a sufficient condition that extends this observation to non-regular graphs,
which does not require any additional condition on the time parameter. 
\begin{thm}\label{5th2}
Let $G$ be a graph on $n$ vertices. For a vector $\u\in\mathbb{R}^n,$
  % Let $U_{M(G)}(t)$ and $U_{M(\overline{G})}(t)$ be the transition matrices of $G$ and its complement $\overline{G},$ respectively. 
  if $\mathbf{1}$ is orthogonal to the set of vectors $\cb{\u,M(G)\u,\ldots, M(G)^{n-1}\u},$ then
  % Suppose $\u$ is a vector with $M(G)^2\u=\u,$ and both $\u$ and $M(G)\u$ are orthogonal to $\mathbf{1}.$ 
   \[U_{M(\overline{G})}(t)\u=\exp{(it\zeta)} U_{M(G)}(-t)\u,\]
 where $U_{M(G)}(t)$ and $U_{M(\overline{G})}(t)$ are the transition matrices of $G$ and $\overline{G}$, respectively, and  the constant $\zeta$ is as defined earlier for $M\in\{A,L,Q\}.$
 \end{thm}
 \begin{proof}
 Since $G$ is a graph of order $n,$ the matrix $M(G)^k$ can be expressed as a linear combination of $I, M(G),\ldots, M(G)^{n-1},$ for all non-negative integers $k.$ As $\mathbf{1}$ is orthogonal to the subspace generated by $\cb{\u,M(G)\u,\ldots, M(G)^{n-1}\u},$ it follows that $JM(G)^k\u=\mathbf{0}.$ Now the result follows from $M(\overline{G})^k\u=(\delta J+\zeta I-M(G))^k\u=(\zeta I-M(G))^k\u.$ 
     % Under the given assumptions on $\u,$ we have $M(\overline{G})^k\u=(\zeta I-M(G))^k\u,$ for every positive integer $k.$ Consequently, the result follows from the definition of the transition matrix.
 \end{proof}
 \begin{exm}\label{5ex2}
     The path $P_5$ exhibits pair PST between $\u=\frac{1}{\sqrt{2}}(\e_1-\e_5)$ and $\v=\frac{1}{\sqrt{2}}(\e_2-\e_4)$ at $\frac{\pi}{2}$ with respect to the adjacency matrix \cite{pal10}. Note that $A(P_5)\u=\v$ and  $A(P_5)^2\u=\u.$ Then by Theorem \ref{5th2}, the graph $\overline{P_5}$ admits PST between $\frac{1}{\sqrt{2}}(\e_1-\e_5)$ and $\frac{1}{\sqrt{2}}(\e_2-\e_4).$
 \end{exm}
% Suppose the premise of Theorem \ref{5th2} holds. Then, for any graph $H$ of order $m,$ and vector $[\u\quad\mathbf{0}]^T\in \mathbb{R}^{m+n},$ the premise also holds for the graph $G\cup H.$ Therefore, Theorem \ref{5th2} gives rise to the following observation. 
As a consequence of Theorem \ref{5th2}, we observe that FR is preserved under graph join.
 \begin{cor}\label{5co2}
     Suppose the premise of Theorem \ref{5th2} holds. If the graph $G$ exhibits fractional revival from $\u$ to $\v$ at $\tau,$ with respect to $M(G),$ then for any graph $H,$ the join $G+H$ admits fractional revival from $[\u^T\quad\mathbf{0}]^T$ to $[\v^T\quad\mathbf{0}]^T$ relative to $M(G+H).$   
 \end{cor}
\begin{proof}
It follows from the proof of Theorem \ref{5th2} that the premise of Theorem \ref{5th2} holds for $\overline{G}.$  For any graph $H$ of order $m,$ observe that the same conditions also hold for the disjoint union $\overline{G}\cup \overline{H}$ with the vector $[\u^T\quad\mathbf{0}]^T\in \mathbb{R}^{m+n}.$   
    Since $G$ admits FR from $\u$ to $\v,$ Theorem \ref{5th2} applies to conclude that $\overline{G}\cup \overline{H}$ admits FR from $[\u^T\quad\mathbf{0}]^T$ to $[\v^T\quad\mathbf{0}]^T.$ Consequently, $G+H$ admits FR from $[\u^T\quad\mathbf{0}]^T$ to $[\v^T\quad\mathbf{0}]^T.$
\end{proof}

\begin{exm} In continuation of  Example \ref{5ex2}, one may apply Corollary \ref{5co2} to show that for any graph $H,$ both $P_5+H$ and $\overline{P_5}+H$ admit PST between $\frac{1}{\sqrt{2}}(\e_1-\e_5)$ and $\frac{1}{\sqrt{2}}(\e_2-\e_4).$  
\end{exm}
 
 
% The sequential join of $H_1, H_2$, and $H_3,$ denoted by $H_1+_s H_2+_s H_3,$ is obtained from the disjoint union of $H_1,H_2$ and $H_3$ by adding all edges between every vertex of $H_i$ and every vertex of $H_{i+1}$ for $i=1,2.$
% The complement of the graph $H_1+_s H_2+_s H_3$ becomes $\ob{\overline{H_1}+\overline{H_3}}\cup \overline{H_2}.$ Corollary \ref{5c3}, yields the following result.
% \begin{cor}
%     Let $H_1, H_2,$ and $H_3$ be graphs. If $\overline{H_2}$ admits fractional revival from $\u$ to $\v$ at $\tau$ relative to the Laplacian matrix and $\tau(|V(H_1)|+|V(H_2)|+|V(H_3)|)\in 2\pi\mathbb{Z},$ then the sequential join $H_1+_s H_2+_s H_3$ exhibits fractional revival from $\u$ to $\v.$
% \end{cor}
% It is well known that $\overline{C_4}$ exhibits plus PST at $\frac{\pi}{2}.$ Therefore, the graph $H_1+_s C_4+_s H_2$ also admits Laplacian plus PST if $|V(H_1)|+|V(H_2)|$ is divisible by $4.$ In the sequential join if $H_1=H_2=H_3,$ then this represents the lexicographic product $P_3[H].$   
% The lexicographic product of two graphs $G$ and $H,$ denoted by $G[H],$ is the graph with vertex set $V(G) \times V(H),$ where two vertices $(a,u)$ and $(b,v)$ are adjacent if either $a$ and $b$ are adjacent in $G,$ or $a=b$ and $u$ and $v$ are adjacent in $H.$ 
% If $H$ is regular, then the existence of vertex PST in $G[H]$ is established in \cite[Lemma 5]{ge}. Here, we provide a similar result for FR from a real pure state with respect to the adjacency matrix.

% \begin{thm}\label{5t16}
% Let $H$ be a regular graph with fractional revival from $\u$ to $\v$ at $\tau.$ Then for a state $\w$ in a graph $G,$ the graph $G[H]$ admits fractional revival from $\w\otimes \u$ to $\w\otimes \v,$ if 
% \[\tau\lambda_j|V(G)|\in{2\pi\mathbb{Z}},\]
% for each eigenvalue $\lambda_j$ of $G.$
% \end{thm}
% \begin{proof}
%  Let $G$ and $H$ be graphs on $m$ and $n$ vertices, respectively, then the adjacency matrix of $G[H]$ is given as follows.
% \[A(G[H])=A(G)\otimes J_m+I_n\otimes A(H)
% \] 
% Since $H$ is a regular graph, $A(H)$ commutes with $J_m,$ therefore,
% \[\exp{(itA(G[H]))}=\exp{(it(A(G)\otimes J_m))}\exp{(I_n\otimes A(H)))}.\]
% Let $G$ have $d$ distinct eigenvalues $\lambda_1,\lambda_2,\ldots, \lambda_d.$ Since the eigenvalues of $J_m$ are $m$ and $0,$ using spectral decomposition we have
% \[\exp{(itA(G[H]))}=\ob{\sum_{j=1}^d\exp{(it\lambda_j m)}E_{\lambda_jm}+E_0}(I_n\otimes \exp{(itA(H))},\]
% where $E_{\lambda_jm}$ and $E_0$ are orthogonal projections of $A(G)\otimes J_m$ corresponding to the eigenvalues $\lambda_jm$ and $0,$ respectively. If $\tau\lambda_jm\in{2\pi\mathbb{Z}},$ then $\exp{(itA(G[H]))}=I_n\otimes \exp{(itA(H))},$ and the result follows.
% \end{proof}

% \begin{exm}
%     The cycle $C_6$ exhibits fractional revival between antipodal vertices $a$ and $b$ at time $\frac{2\pi}{3}.$ The eigenvalues of complete graph $K_n$ are $n-1$ and $-1.$ By Theorem $\ref{5t16},$ if $n$ is divisible by $3,$ then the lexicographic product $K_n[C_6]$ admits fractional revival from $e_g\otimes \e_a$ to $e_g\otimes \e_b,$ for every $g\in V(K_n).$  Moreover, since $C_6$ is  a regular graph, Theorem $\ref{5t15}$ implies that $K_n[\overline{C_6}]$ also exhibits fractional revival from $e_g\otimes \e_a$ to $e_g\otimes \e_b,$ for each $g\in V(K_n),$ whenever $n$ is divisible by $3.$
% \end{exm}

\section{Double covers}\label{5sec3}
Let $X_1$ and $X_2$ be graphs on the same vertex set. The graph $X_1\ltimes X_2$ is defined by the adjacency matrix \[\begin{bmatrix}
    A(X_1) & A(X_2)\\
    A(X_2) & A(X_1)
\end{bmatrix}.\] 
If $X_1$ and $X_2$ have disjoint edge sets, then $X_1\ltimes X_2$ is called a double cover of the graph with adjacency matrix $A(X_1)+A(X_2).$ 
% The graph $X_1\ltimes X_2$ with adjacency matrix 
% \[\begin{bmatrix}
%     A(X_1) & A(X_2)\\
%     A(X_2) & A(X_1)
% \end{bmatrix}\]
% is called a double cover of the graph $G.$ 
Let $D(X_1)$ and $D(X_2)$ be the degree matrices of graphs $X_1$ and $X_2,$ respectively. Consider $D=D(X_1)+D(X_2).$ Then the matrix $M\in\{A,L,Q\}$ associated with the graph $X_1\ltimes X_2$ becomes
\[M(X_1\ltimes X_2)=\begin{bmatrix}
    \eta D+\delta A(X_1) & \delta A(X_2)\\
   \delta A(X_2) & \eta D+\delta A(X_1)
\end{bmatrix},\quad \text{where}\]
 \[\eta=\left\{ \begin{array}{rcl}
 0, & \mbox{if} & M=A, \\ 
 1, & \mbox{if} & M\in\{L,Q\},
 \end{array}\right.\quad 
\text{and} \quad \delta=\left\{ \begin{array}{rcl}
 -1, & \mbox{if} & M =L, \\ 
 1, & \mbox{if} & M \in\{A,Q\}. 
 \end{array}\right.\] 

 We extend \cite[Lemma 5.1]{cou} to obtain a unified expression for the transition matrix of $X_1\ltimes X_2$ with respect to the adjacency, Laplacian, and signless Laplacian matrices. 
 \begin{thm}\label{5t19}
    Let $X_1$ and $X_2$ be graphs on the same vertex set $V.$ 
    % If 
    % \[M(G_{\pm})=\eta D+\delta(A(X_1)\pm A(X_2)),\] then for all $t\in \mathbb{R},$ 
    The transition matrix of  $X_1\ltimes X_2$ is given by 
 \[U_{M({X_1\ltimes X_2})}(t)=\frac{1}{2}\begin{bmatrix}
    U_{M(G_+)}(t)+ U_{M(G_-)}(t) & U_{M(G_+)}(t)- U_{M(G_-)}(t)\\
    U_{M(G_+)}(t)- U_{M(G_-)}(t)& U_{M(G_+)}(t)+ U_{M(G_-)}(t)
\end{bmatrix},\]
where $t\in \mathbb{R},$ and $M(G_{\pm})=\eta D+\delta(A(X_1)\pm A(X_2)).$
\end{thm}
\begin{proof}
  Let $A(K_2)$ be the adjacency matrix of a complete graph on two vertices, then the matrix $M\in\{A,L,Q\}$ associated with the graph $X_1\ltimes X_2$ can be written as
\[M(X_1\ltimes X_2)=I_2\otimes(\eta D+\delta A(X_1))+A(K_2)\otimes  (\delta A(X_2)). \]  
Note that $I_2$ and $A(K_2)$ commute and can be simultaneously diagonalized by the matrix \[H=\frac{1}{\sqrt{2}}\begin{bmatrix}
    1 & 1\\
    1 & -1
\end{bmatrix}.\] 
Since $(H\otimes I_n)^{-1}=H\otimes I_n,$ it follows from \cite[Lemma 3.1]{cou} that
\[(H\otimes I_n)U_{M(X_1\ltimes X_2)}(t)(H\otimes I_n)=\begin{bmatrix}
    U_{M(G_+)}(t) & \mathbf{0}\\
     \mathbf{0} & U_{M(G_-)}(t)
    \end{bmatrix},\]
    where $M(G_{\pm})=\eta D+\delta(A(X_1)\pm A(X_2)),$ and the result follows. 
%     Since $(H\otimes I_n)^{-1}=H\otimes I_n=\frac{1}{\sqrt{2}}\begin{bmatrix}
%     I_n & I_n\\
%     I_n & -I_n
% \end{bmatrix},$ we have the desired result.
\end{proof}
A characterization of vertex PST in the graph $X_1\ltimes X_2$
 was established in \cite[Theorem 5.2]{cou}. Furthermore, Laplacian pair PST was studied in \cite[Theorem 4.2]{jia} when both $X_1$ and $X_2$ are regular. We extend these results to the existence of FR in $X_1\ltimes X_2$ with respect to $M\in\{A,L,Q\}.$ 
 % The result follows from the expression of the transition matrix of $X_1\ltimes X_2$ obtained in Theorem \ref{5t19}.
\begin{thm}\label{5t21}
    Let $X_1$ and $X_2$ be graphs on the same vertex set $V.$ Let $M(G_{\pm})=\eta D+\delta (A(X_1)\pm A(X_2))$ be the corresponding matrices of $G_+$ and $G_-,$ respectively. Then for the states $\u, \v\in \mathbb{R}^{|V|},$ and $\alpha,\beta\in\mathbb{C}$ with $\beta\neq 0,$ the following holds.
   \begin{enumerate}

\item $X_1\ltimes X_2$ admits $(\alpha,\beta)$-fractional revival from $\tb{\u^T \quad \mathbf{0}}^T$
 to $\tb{\mathbf{0} \quad 
\u^T}^T$ at $\tau$ if and only if $\u$ is periodic at $\tau$ with respect to $M(G_+)$ and $M(G_-)$ with phase factor $\alpha+\beta$ and $\alpha-\beta,$ respectively. 

\item $X_1\ltimes X_2$ admits $(\alpha,\beta)$-fractional revival at $\tau$ from state $\tb{\u^T \quad \mathbf{0}}
^T$ to  $\tb{\v^T \quad \mathbf{0}}
^T$ $\ob{\text{or, from} \tb{  \mathbf{0} \quad \u^T}
^T \text{to}   \tb{ 
\mathbf{0} \quad \v^T}^T}$ if and only if  $U_{M(G_\pm)}(\tau)\u=\alpha\u+\beta\v.$ 

\item $X_1\ltimes X_2$ admits $(\alpha,\beta)$-fractional revival at $\tau$ from state $\tb{\u^T \quad \mathbf{0}}
^T$ to $\tb{\mathbf{0} \quad 
\v^T}^T$ $\ob{\text{or, from} \tb{  \mathbf{0} \quad \u^T}
^T \text{to}   \tb{ 
\v^T \quad \mathbf{0}}^T}$ if and only if  $U_{M(G_\pm)}(\tau)\u=\alpha\u\pm \beta\v.$

\item $G_+$ admits $(\alpha,\beta)$-fractional revival at $\tau$ from $\u$ to $\v$ if and only if $X_1\ltimes X_2$ admits $(\alpha,\beta)$-fractional revival at $\tau$ from $\frac{1}{\sqrt{2}}\tb{\u^T \quad \u^T}
^T$ to   $\frac{1}{\sqrt{2}}\tb{\v^T \quad \v^T}
^T.$ 

\item $G_-$ admits $(\alpha,\beta)$-fractional revival at $\tau$ from $\u$ to $\v$ if and only if $X_1\ltimes X_2$ admits $(\alpha,\beta)$-fractional revival at $\tau$ from $\frac{1}{\sqrt{2}}\tb{\u^T \quad -\u^T}
^T$ to   $\frac{1}{\sqrt{2}}\tb{\v^T \quad -\v^T}
^T.$
\end{enumerate} 

\end{thm}

% One may observe Theorem \ref{5t21}
% as a generalization of vertex PST in the double cover $X_1\ltimes X_2$ \cite[Theorem 5.2]{cou}. In particular, if the graph $X_1$ and $X_2$ be $r_1$ and $r_2$-regular, respectively, then \cite[Theorem 4.2]{jia} can be obtained from Theorem \ref{5t21}.
% as for regular graphs the results of perfect state transfer for adjacency and Laplacian matrices are equivalent.
% A connection between fractional revival in a graph and its double cover, independent of the parameters 
% $\alpha$ and $\beta,$ is established.
% \begin{thm}
% Let $X_1$ and $X_2$ be graphs on the same vertex set $V.$ Let $M(G_{\pm})=\eta D+\delta (A(X_1)\pm A(X_2))$ be the corresponding matrices of $G_+$ and $G_-,$ respectively. Then for the states $\u, \v\in \mathbb{R}^{|V|},$ the following hold.
% \begin{enumerate}

% % \item $G_+$ admits fractional revival from $\u$ to $\v$ if and only if $X_1\ltimes X_2$ admits fractional revival from $\frac{1}{\sqrt{2}}\tb{\u^T \quad \u^T}
% % ^T$ to   $\frac{1}{\sqrt{2}}\tb{\v^T \quad \v^T}
% % ^T.$ 

% % \item $G_-$ admits fractional revival from $\u$ to $\v$ if and only if $X_1\ltimes X_2$ admits fractional revival from $\frac{1}{\sqrt{2}}\tb{\u^T \quad -\u^T}
% % ^T$ to   $\frac{1}{\sqrt{2}}\tb{\v^T \quad -\v^T}
% % ^T.$
% \end{enumerate}  
% \end{thm}
\begin{proof}
  Let $X_1\ltimes X_2$ admit $(\alpha,\beta)$-FR from $\tb{\u^T\quad\mathbf{0}}^T$ to $\tb{\mathbf{0}\quad\u^T}^T$ at $\tau.$ Then  
  \begin{equation}\label{5eq5}
      U_{M(X_1\ltimes X_2)}(\tau) \begin{bmatrix}
    \u \\ \mathbf{0}
\end{bmatrix}=\alpha \begin{bmatrix}
    \u \\ \mathbf{0}
\end{bmatrix}+\beta \begin{bmatrix}
   \mathbf{0} \\ \u
\end{bmatrix}=\begin{bmatrix}
      \alpha\u \\ \beta\u
  \end{bmatrix}.
  \end{equation}

Using Theorem \ref{5t19}, we have 
  \begin{equation}\label{5eqn6}
      U_{M({X_1\ltimes X_2})}(\tau)\cdot \begin{bmatrix}
    \u \\ \mathbf{0}
\end{bmatrix}=\frac{1}{2}\begin{bmatrix}
    \ob{U_{M(G_+)}(\tau)+U_{M(G_-)}(\tau)}\u \\ \ob{U_{M(G_+)}(\tau)-U_{M(G_-)}(\tau)}\u
\end{bmatrix}
.\end{equation}
Comparing \eqref{5eq5} and \eqref{5eqn6}, we obtain $U_{M(G_+)}(\tau)\u=(\alpha+\beta)\u$ and $U_{M(G_-)}(\tau)\u=(\alpha-\beta)\u.$
The converse part of the proof is immediate from \eqref{5eqn6}. This proves the first statement. The remaining statements follow by similar arguments.
\end{proof}
In particular, Theorem \ref{5t21}  relates FR from vertex state to FR from plus and pair states.

% if $\u$ and $\v$ are vertex states $\e_a$ and $\e_b,$ respectively, the result follows.
\begin{cor}\label{5cor3}
Suppose the premise of Theorem \ref{5t21} holds. If $\{0,1\}\times V$ is the vertex set of $X_1\ltimes X_2,$ then the following holds. 
\begin{enumerate}
\item $G_+$ admits $(\alpha,\beta)$-fractional revival at $\tau$ from vertex $a$ to $b,$ if and only if $X_1\ltimes X_2$ admits $(\alpha,\beta)$-fractional revival at $\tau$ from $\frac{1}{\sqrt{2}}(\e_{(0,a)}+\e_{(1,a)})$ to   $\frac{1}{\sqrt{2}}(\e_{(0,b)}+\e_{(1,b)}).$ 

\item $G_-$ admits $(\alpha,\beta)$-fractional revival at $\tau$ from vertex $a$ to $b$ if and only if $X_1\ltimes X_2$ admits $(\alpha,\beta)$-fractional revival at $\tau$ from $\frac{1}{\sqrt{2}}(\e_{(0,a)}-\e_{(1,a)})$ to $\frac{1}{\sqrt{2}}(\e_{(0,b)}-\e_{(1,b)}).$ 
\end{enumerate}  
\end{cor}


\begin{figure}[]
\centering

\begin{tabular}{cccc}

\begin{minipage}{0.22\textwidth}\centering
\begin{tikzpicture}[scale=1.2,auto=left]
\tikzstyle{every node}=[circle, thick, black!90, fill=white, scale=0.65]
\node[draw,minimum size=0.55cm, inner sep=0 pt] (1) at (1,1) {$0$};
\node[draw,minimum size=0.55cm, inner sep=0 pt] (2) at (0,1) {$1$};
\node[draw,minimum size=0.55cm, inner sep=0 pt] (3) at (0,0) {$2$};
\node[draw,minimum size=0.55cm, inner sep=0 pt] (4) at (1,0) {$3$};
\draw (1)-- (2)-- (3)--(4)-- (1);
\end{tikzpicture}

$C_4$
\end{minipage}
&
\begin{minipage}{0.22\textwidth}\centering
\begin{tikzpicture}[scale=1.2,auto=left]
\tikzstyle{every node}=[circle, thick, black!90, fill=white, scale=0.65]
\node[draw,minimum size=0.55cm, inner sep=0 pt] (1) at (1,1) {$0$};
\node[draw,minimum size=0.55cm, inner sep=0 pt] (2) at (0,1) {$1$};
\node[draw,minimum size=0.55cm, inner sep=0 pt] (3) at (0,0) {$2$};
\node[draw,minimum size=0.55cm, inner sep=0 pt] (4) at (1,0) {$3$};
\draw (1)--(2)--(3)--(4);
\end{tikzpicture}

$X_1$
\end{minipage}
&
\begin{minipage}{0.22\textwidth}\centering
\begin{tikzpicture}[scale=1.2,auto=left]
\tikzstyle{every node}=[circle, thick, black!90, fill=white, scale=0.65]
\node[draw,minimum size=0.55cm, inner sep=0 pt] (1) at (1,1) {$0$};
\node[draw,minimum size=0.55cm, inner sep=0 pt] (2) at (0,1) {$1$};
\node[draw,minimum size=0.55cm, inner sep=0 pt] (3) at (0,0) {$2$};
\node[draw,minimum size=0.55cm, inner sep=0 pt] (4) at (1,0) {$3$};
\draw (4)-- (1);
\end{tikzpicture}

$X_2$
\end{minipage}
&
\begin{minipage}{0.22\textwidth}\centering
\begin{tikzpicture}[ scale=0.5]
\tikzstyle{every node}=[circle, thick, black!90, fill=white, scale=0.6]
\coordinate (top) at (6,2.1);
\coordinate (bottom) at (6,-2.1);
\node[draw,minimum size=0.6cm,inner sep=0 pt] (0) at (3,0) {$(1,2)$};
\node[draw,minimum size=0.6cm,inner sep=0 pt] (1) at (2.12,2.12) {$(1,3)$};
\node[draw,minimum size=0.6cm, inner sep=0 pt] (2) at (0,3){$(0,0)$};
\node[draw,minimum size=0.6cm, inner sep=0 pt] (3) at (-2.12,2.12) {$(0,1)$};
\node[draw,minimum size=0.6cm, inner sep=0 pt] (4) at (-3,0) {$(0,2)$};
\node[draw,minimum size=0.6cm, inner sep=0 pt] (5) at (-2.12,-2.12) {$(0,3)$};
\node[draw,minimum size=0.6cm, inner sep=0 pt] (6) at (0,-3) {$(1,0)$};
\node[draw,minimum size=0.6cm, inner sep=0 pt] (7) at (2.12,-2.12) {$(1,1)$};
\draw (0)--(1)--(2)--(3)--(4)--(5)--(6)--(7)--(0);
\end{tikzpicture}

$X_1 \ltimes X_2$
\end{minipage}

\end{tabular}

\caption{$C_8$ as double cover of $C_4.$}
\label{5fig3}
\end{figure}

% Let $G$ be a graph with an edge $e$ connecting the vertices $a$ and $b.$ Let $X_1=G\backslash \{e\}$ and let $X_2$ be the spanning subgraph of $G$ consisting only of the edge $e.$ The double cover $X_1\ltimes X_2$ is obtained by taking two disjoint copies of $G\backslash \{e\}$ and adding two edges between vertex $a$ in one copy to vertex $b$ in other. Then 
% Let $G$ be the graph with adjacency matrix $A(X_1)+A(X_2).$ 
One may observe from Corollary~\ref{5cor3} that if a graph $G$ admits vertex PST with respect to $M(G)$, then every double cover $X_1\ltimes X_2$ of $G$ exhibits plus PST with respect to $M(X_1\ltimes X_2).$ Consider the cycle $C_4$, shown in Figure~\ref{5fig3}, together with two spanning subgraphs $X_1$ and $X_2$. It is well known that $C_4$ has vertex PST between the vertices $0$ and $2$ at time $\frac{\pi}{2}$. By Corollary~\ref{5cor3}(1), the double cover $X_1\ltimes X_2$ of $C_4$, which is the cycle $C_8$, exhibits plus PST at time $\frac{\pi}{2}$ between $\frac{1}{\sqrt{2}}\big(\e_{(0,0)}+\e_{(1,0)}\big)$ and $\frac{1}{\sqrt{2}}\big(\e_{(0,2)}+\e_{(1,2)}\big),$  which is consistent with \cite[Theorem 6.5]{kim}.

\section{Plus PGST in cycles}\label{5sec5}

% \begin{lem}\label{5l7}
% Let a graph $G$ admit pretty good plus state transfer between $\frac{1}{\sqrt{2}}\ob{\e_a+\e_b}$ and $\frac{1}{\sqrt{2}}\ob{\e_c+\e_d}.$ If $P$ is a permutation matrix corresponding to an automorphism of $G$ such that $P\e_a=\e_b$ and $P\e_b\neq \e_a,$ then the graph $G$ admits pretty good pair state transfer.

% \end{lem}
% \begin{proof}
%  Suppose the graph $G$ admits plus PGST between $\frac{1}{\sqrt{2}}\ob{\e_a+\e_b}$ and $\frac{1}{\sqrt{2}}\ob{\e_c+\e_d},$ then by \eqref{5eq2}, we have 
% % Then there is a sequence of time $t_k\in\mathbb{R}$ and a unit complex number $\gamma$ such that  
%     \begin{equation}\label{5eqn7}
%         \lim_{k \to \infty}U_G\ob{t_k}\ob{\e_a+\e_b}=\gamma\ob{\e_c+\e_d}.
%     \end{equation}
%    Multiplying $P$ on both sides of \eqref{5eqn7} and subtracting the resulting equation from \eqref{5eqn7}, yields
%    \begin{equation}\label{5eqn8}
%        \lim_{k \to \infty}U_G\ob{t_k}\ob{\e_a-P\e_b}=\gamma\ob{\e_c+\e_d-P\e_c-P\e_d}.
%    \end{equation}
%    Since $U_G(t)$ is unitary, \eqref{5eqn8} implies that $\left\|\frac{1}{\sqrt{2}}\ob{\e_b-P\e_b}\right\|=\left\|\frac{1}{\sqrt{2}}\ob{\e_c+\e_d-P\e_c-P\e_d}\right\|.$ This equality is possible only if $\frac{1}{\sqrt{2}}\ob{\e_c+\e_d-P\e_c-P\e_d}$ represents a pair state. 
% \end{proof}
Let $(\Gamma,+)$ be a finite abelian group, and let $S \subseteq\Gamma\backslash\{0\}$ be a connection set satisfying $\left\lbrace -s:s\in S\right\rbrace=S.$ The \emph{Cayley graph}, denoted by $\text{Cay}(\Gamma, S),$ has the vertex set $\Gamma$ where two vertices $a$ and $b$ are adjacent if and only if $a-b\in S$. The Cayley graph over $\mathbb{Z}_n,$ the group of integers modulo $n,$ is called a \emph{circulant graph}. In particular, the cycle $C_n$ is a circulant graph over $\mathbb{Z}_n$ with the connection set $\{1,n-1\}.$
Since a circulant graph is regular,
we investigate plus PGST with respect to the adjacency matrix.
% We investigate the existence of plus PGST in circulant graphs with respect to the adjacency matrix. 
% The eigenvalues and eigenvectors of a cycle are well known. Suppose $\omega_n=\exp {\ob{\frac{2\pi i}{n}}}$
% is
% the primitive $n$-th root of unity. Then the eigenvalues of $C_n$ are
% $\lambda_r=2\cos\ob{{\frac{2r\pi}{n}}},$ where $ 0\leq r\leq n-1,   
% $
% and the corresponding eigenvectors are \begin{equation}\label{5eq7}\x_r =\tb{1\quad \omega_n^r\quad\cdots\quad\omega_n^{r(n-1)}}^T.
% \end{equation}
% The complement $\overline{C_n}$ has eigenvalues $\mu_0=n-3,$ and $\mu_r=-1-\lambda_r,$ for $1\leq r \leq n-1,$ corresponding to the eigenvectors as given in \eqref{5eq7}. 
% Before we investigate plus PGST in the circulant graphs, we include the following observation.
%  \begin{lem}\label{5l7}
%  Let a graph $G$ admit pretty good plus state transfer between $\frac{1}{\sqrt{2}}\ob{\e_a+\e_b}$ and $\frac{1}{\sqrt{2}}\ob{\e_c+\e_d}.$ Then the graph $G$ admits pretty good pair state transfer if there exists an automorphism of $G$ with  permutation matrix 
%  $P$ satisfying one of the following conditions.
%  \begin{enumerate}
%     \item $P\e_a=\e_a$ and $P\e_b\neq \e_b,$ 
%     \item  $P\e_a=\e_b$ and $P\e_b\neq \e_a.$ 
% \end{enumerate}
%  \end{lem}
%  \begin{proof}
%   Suppose the graph $G$ admits plus PGST between $\frac{1}{\sqrt{2}}\ob{\e_a+\e_b}$ and $\frac{1}{\sqrt{2}}\ob{\e_c+\e_d}.$ Then by \eqref{5eq2}, we have 
%   \begin{equation}\label{5eqn7}
%          \lim_{k \to \infty}U_G\ob{t_k}\ob{\e_a+\e_b}=\gamma\ob{\e_c+\e_d}.
%      \end{equation}
%     Suppose condition-1 holds. Then
%      multiplying $P$ on both sides of \eqref{5eqn7} and subtracting the resulting equation from \eqref{5eqn7}, yields
%  \begin{equation}\label{5eqn8}
%         \lim_{k \to \infty}U_G\ob{t_k}\ob{\e_b-P\e_b}=\gamma\ob{\e_c+\e_d-P\e_c-P\e_d}.
%    \end{equation}
%     Since $U_G(t)$ is unitary, \eqref{5eqn8} implies that $\left\|\frac{1}{\sqrt{2}}\ob{\e_b-P\e_b}\right\|=\left\|\frac{1}{\sqrt{2}}\ob{\e_c+\e_d-P\e_c-P\e_d}\right\|.$ This equality is possible only if $\frac{1}{\sqrt{2}}\ob{\e_c+\e_d-P\e_c-P\e_d}$ represents a pair state. Hence, $G$ admits pair PGST.  
%     Similarly, if condition-2 holds, an analogous argument shows that  
%     $G$ admits pair PGST.
%  \end{proof}
 
% \begin{rem}
 If $n$ is odd, then there exists an automorphism of $\text{Cay}(\mathbb{Z}_n, S)$ fixing only one vertex. Hence, by Theorem \ref{5l7}, if $\text{Cay}(\mathbb{Z}_n, S)$ admits plus PGST, then the graph must exhibit pair PGST. However, there is no pair PGST in a Cayley graph over an abelian group of odd order \cite [Theorem 10]{pal10}, leading to the following result.   

% Any two vertices $a$ and $b$ in a circulant graph $\text{Cay}(\mathbb{Z}_n, S)$ of odd order are not antipodal. Consequently, Lemma \ref{5l7} implies that if $\text{Cay}(\mathbb{Z}_n, S)$ admits plus PGST from $\frac{1}{\sqrt{2}}\ob{\e_a+\e_b}$, then the graph must exhibit pair PGST. However, this is ruled out by \cite [Theorem 10]{pal10}, leading to the following result.   
% If $P$ is a permutation matrix corresponding to an automorphism of $G$ such that $P\e_a=\e_b$ and $P\e_b\neq \e_a,$ then Lemma \ref{5l7} holds for $G.$ In particular, Lemma \ref{5l7} holds for the circulant graphs.
% \end{rem}
% Note that Lemma \ref{5l7} also holds for the complement $\overline{C_n}.$
% Since an odd cycle and its complement have no antipodal vertices, Lemma \ref{5l7} implies that if they admit plus PGST, then they must exhibit pair PGST. However, this is ruled out by \cite [Theorem 10]{pal10}, leading to the following result.   
\begin{lem}\label{5l8}
    A circulant graph of odd order does not exhibit pretty good plus state transfer.
\end{lem}



% A complete characterization of vertex PGST and pair PGST in $C_n$ is given in \cite[Theorem 13]{pal4} and \cite[Theorem 10]{pal10} as follows.


% \begin{thm}\cite{pal4}\label{5t24} A cycle $C_n$ as well as its complement $\overline{C_n}$ admit pretty good state transfer
% if and only if $n=2^k$ for $k\geq 2,$ and it occurs between every pair of antipodal vertices.
% \end{thm}

% \begin{thm}\cite{pal10}\label{5th25}
%  A cycle $C_n$ on $n$ vertices admits pretty good pair state transfer if and only
% if either $n = 2^k$ or $n = 2^kp,$ where $k$ is a positive integer and $p$ is an odd prime.
% \end{thm}
% We obtain the following as an immediate consequence of Theorem \ref{5t24}.
% \begin{thm}\label{5t26}
% Let $C_n$ be a cycle with two vertices $a$ and
% $b,$ where $n=2^k,$ with $k\geq 2.$ Pretty good plus state transfer occurs in $C_n$ and its complement $\overline{C_n}$ between $\frac{1}{\sqrt{2}}\ob{\e_a+\e_b}$ and $\frac{1}{\sqrt{2}}\ob{\e_{a+\frac{n}{2}}+\e_{b+\frac{n}{2}}}$ if and only
% if $a-b\neq \frac{n}{2}.$
% \end{thm}

 % The following result shows the existence of plus PGST in cycles implies the existence of pair PGST.

% \begin{lem}\label{5l7}
% Let the cycle $C_n$ exhibit pretty good plus state transfer between $\frac{1}{\sqrt{2}}\ob{\e_a+\e_b}$ and $\frac{1}{\sqrt{2}}\ob{\e_c+\e_d}.$ If $a$ and $b$ are not antipodal vertices, then $C_n$ admits pretty good pair state transfer.
% \end{lem}
% \begin{proof}
% Let $a$ and $b$ be non-antipodal vertices in $C_n.$ If the cycle $C_n$ exhibits plus PGST between $\frac{1}{\sqrt{2}}\ob{\e_a+\e_b}$ and $\frac{1}{\sqrt{2}}\ob{\e_c+\e_d},$ then by \eqref{5eq2}, we have 
% % Then there is a sequence of time $t_k\in\mathbb{R}$ and a unit complex number $\gamma$ such that  
%     \begin{equation}\label{5e11}
%         \lim_{k \to \infty}U_G\ob{t_k}\ob{\e_a+\e_b}=\gamma\ob{\e_c+\e_d}.
%     \end{equation}
%     Let $P$ be a permutataion matrix of an automorphism of $C_n$ such that $P\e_a=\e_b$ and $P\e_b\neq\e_a.$
%    Multiplying $P$ on both sides of \eqref{5e11} and subtracting the resulting equation from \eqref{5e11}, we obtain
%    \begin{equation}\label{5e12}
%        \lim_{k \to \infty}U_G\ob{t_k}\ob{\e_a-P\e_b}=\gamma\ob{\e_c+\e_d-P\e_c-P\e_d}.
%    \end{equation}
%    Since $U_G(t)$ is unitary, \eqref{5e12} implies that $\left\|\frac{1}{\sqrt{2}}\ob{\e_b-P\e_b}\right\|=\left\|\frac{1}{\sqrt{2}}\ob{\e_c+\e_d-P\e_c-P\e_d}\right\|.$ This equality is possible only if $\frac{1}{\sqrt{2}}\ob{\e_c+\e_d-P\e_c-P\e_d}$ represents a pair state. 
% \end{proof}
% \begin{rem}
% Let $a$ and $b$ be two vertices of a graph $G.$ If $P$ is a permutation matrix corresponding to an automorphism of $G$ such that $P\e_a=\e_b$ and $P\e_b\neq \e_a,$ then Lemma \ref{5l7} holds for $G.$ In particular, Lemma \ref{5l7} holds for the circulant graphs.
% \end{rem}
% Note that Lemma \ref{5l7} also holds for the complement $\overline{C_n}.$
% Since an odd cycle and its complement have no antipodal vertices, Lemma \ref{5l7} implies that if they admit plus PGST, then they must exhibit pair PGST. However, this is ruled out by \cite [Theorem 10]{pal10}, leading to the following result.   
% \begin{lem}\label{5l8}
%     The odd cycle and its complement do not admit pretty good plus state transfer.
% \end{lem}

A pair of vertices $a$ and $b$ in a circulant graph $\text{Cay}(\mathbb{Z}_n, S)$ are called \emph{antipodal vertices} if $a-b=\frac{n}{2}.$ The existence of plus PGST from a pair of non-antipodal vertices in the circulant graph of even order is shown in the following result. 
 % Let $n$ be an even number, and $a$ and $b$ are non-antipodal vertices in $\text{Cay}(\mathbb{Z}_n, S).$  The existence of plus PGST from $\frac{1}{\sqrt{2}}\ob{\e_a+\e_b}$ in  $\text{Cay}(\mathbb{Z}_n, S)$ is shown in the following result. 
\begin{lem}\label{5l9}
Let $a$ and $b$ be non-antipodal vertices of the circulant graph $\text{Cay}(\mathbb{Z}_n, S),$ where $n$ is even. If $\text{Cay}(\mathbb{Z}_n, S)$ exhibits pretty good plus state transfer between $\frac{1}{\sqrt{2}}\ob{\e_a+\e_b}$ and $\frac{1}{\sqrt{2}}\ob{\e_c+\e_d},$ then $\e_c+\e_d=\e_{a+\frac{n}{2}}+\e_{b+\frac{n}{2}}.$
\end{lem}
\begin{proof}
     Without loss of generality, let $a=0$ and $b<\frac{n}{2}.$ Suppose $\text{Cay}(\mathbb{Z}_n, S)$ admits plus PGST between $\frac{1}{\sqrt{2}}\ob{\e_0+\e_b}$ and $\frac{1}{\sqrt{2}}\ob{\e_c+\e_d}.$   Consider the automorphism of $\text{Cay}(\mathbb{Z}_n, S)$ with permutation matrix $P$ satisfying $P\e_j=\e_{n-j}$ that fixes only $0$ and $\frac{n}{2}.$ Since $b\neq \frac{n}{2},$ it follows from \eqref{5eqn8} as argued in the proof of Theorem \ref{5l7}, that $\frac{1}{\sqrt{2}}\ob{\e_c+\e_d-P\e_c-P\e_d}$ must be a pair state. If $\e_c=P\e_d,$ then $\e_d=P\e_c,$ which yields  $\frac{1}{\sqrt{2}}\ob{\e_c+\e_d-P\e_c-P\e_d}=0,$ a contradiction.  Consequently, the automorphism must fix either $c$ or $d,$ that is, one of them lies in the set $\{0, \frac{n}{2}\}.$ 
Therefore, without loss of generality, assume that $c\in\{0,\frac{n}{2}\}.$ 
\\\textit{Case I} $(c=0)$:
When $c=0,$ the graph $\text{Cay}(\mathbb{Z}_n, S)$  admits plus PGST between $\frac{1}{\sqrt{2}}\ob{\e_0+\e_b}$ and $\frac{1}{\sqrt{2}}\ob{\e_0+\e_d}.$
Consider an automorphism of $\text{Cay}(\mathbb{Z}_n, S)$ with permutation matrix $P'$ such that $P'\e_j=\e_{2b-j}$ which fixes only $b$ and $ b+\frac{n}{2}.$ Using a similar argument as before, we observe that $\frac{1}{\sqrt{2}}\ob{\e_0+\e_d-P'\e_0-P'\e_d}$ must be a pair state. Since $P'$ does not fix $\e_0,$ it must fix $\e_d.$ 
% we obtain $\e_d=\e_{b+\frac{n}{2}}.$
Consequently, $d=b+\frac{n}{2},$ and plus PGST occurs in $\text{Cay}(\mathbb{Z}_n, S)$  between $\frac{1}{\sqrt{2}}\ob{\e_0+\e_b}$ and $\frac{1}{\sqrt{2}}\ob{\e_0+\e_{b+\frac{n}{2}}}.$ Now, consider another automorphism with permutation matrix $P''$ such that $P''\e_j=\e_{b+j}.$ Since $b<\frac{n}{2},$ we have $ P''\e_b\neq\e_0.$ Using an argument similar to that in the proof of Theorem \ref{5l7}, we obtain 
   \begin{equation}\label{5eq18}
       \lim_{k \to \infty}U\ob{t_k}\ob{\e_0-\e_{2b}}=\gamma\ob{\e_0+\e_{b+\frac{n}{2}}-\e_b-\e_{2b+\frac{n}{2}}}.
   \end{equation} 
   Since $\e_0\notin\cb{ \e_b,\e_{b+\frac{n}{2}}},$ we have  $\e_{b+\frac{n}{2}}\neq \e_{2b+\frac{n}{2}}.$ Then the equality in \eqref{5eq18} holds only if $\e_0=\e_{2b+\frac{n}{2}},$ which implies $b=\frac{n}{4}.$ Consequently, pair PGST occurs in $\text{Cay}(\mathbb{Z}_n, S)$ from $\frac{1}{\sqrt{2}}\ob{\e_0-\e_\frac{n}{2}},$ which is a contradiction to \cite[Lemma 1]{pal10} as there is an automorphism in $\text{Cay}(\mathbb{Z}_n, S)$ that fixes only $0$ and $\frac{n}{2}.$ 
  % This is a contradiction since an argument analogous to that used in \cite[Lemma 7]{pal10}, shows that  
   % Hence, $c\neq 0,$ and therefore, $c$ must be $\frac{n}{2}.$
 \\\textit{Case II} $(c=\frac{n}{2}):$ If $c=\frac{n}{2},$ then $\text{Cay}(\mathbb{Z}_n, S)$  admits plus PGST between $\frac{1}{\sqrt{2}}\ob{\e_0+\e_b}$ and $\frac{1}{\sqrt{2}}\ob{\e_\frac{n}{2}+\e_d}.$ Since the automorphism with permutation matrix $P'$ fixes only $b$ and $b+\frac{n}{2},$ it follows that $\frac{1}{\sqrt{2}}\ob{\e_\frac{n}{2}+\e_d-\e_{2b-\frac{n}{2}}-\e_{2b-d}}$ is a pair state as observed from \eqref{5eqn8}. If  $d= 2b-\frac{n}{2},$ then  the state is $\mathbf{0},$ a contradiction. Since $\e_\frac{n}{2}\neq \e_{2b-\frac{n}{2}},$ then the only possibility is either $d=b$ or $d= b+\frac{n}{2}.$ If $d=b,$ then plus PGST occurs between $\frac{1}{\sqrt{2}}\ob{\e_0+\e_b}$ and $\frac{1}{\sqrt{2}}\ob{\e_\frac{n}{2}+\e_b}.$ Using the permutation matrix $P''$ and applying a similar argument as before, we find pair PGST occurs from $\frac{1}{\sqrt{2}}\ob{\e_\frac{n}{4}-\e_\frac{3n}{4}},$ that is from a pair of antipodal vertices. This is again a contradiction to \cite[Lemma 1]{pal10}. Therefore, we must have $d=b+\frac{n}{2}.$
  \end{proof} 
 % Since $a$ and $b$ are arbitrary non-antipodal vertices and $\text{Cay}(\mathbb{Z}_n, S)$ is vertex transitive, the result follows.
     % Then \eqref{5eq2} gives
     % there exists a sequence $t_k\in\mathbb{R}$ and a complex number $\gamma$ of unit modulus such that 
% \begin{equation}\label{5eq13}
%        \lim_{k \to \infty}U\ob{t_k}\ob{\e_0+\e_b}=\gamma\ob{\e_c+\e_d}.
%    \end{equation}


We now characterize plus PGST for the case when 
$a$ and 
$b$ are antipodal vertices in a circulant graph of even order.
\begin{lem}\label{5l10}
    Let $n$ be even and $a$ and $b$ be antipodal vertices in the circulant graph $\text{Cay}(\mathbb{Z}_n, S).$ If there is pretty good plus state transfer from $\frac{1}{\sqrt{2}}\ob{\e_a+\e_b},$ then $n$ is divisible by $4$ and pretty good plus state transfer occurs between $\frac{1}{\sqrt{2}}\ob{\e_a+\e_b}$ and $\frac{1}{\sqrt{2}}\ob{\e_{a+\frac{n}{4}}+\e_{b+\frac{n}{4}}}.$  
\end{lem}

\begin{proof}
    Without loss of generality, let $a=0$ and $b=\frac{n}{2}.$ Suppose $\text{Cay}(\mathbb{Z}_n, S)$ admits plus PGST between $\frac{1}{\sqrt{2}}\ob{\e_0+\e_{\frac{n}{2}}}$ and $\frac{1}{\sqrt{2}}\ob{\e_c+\e_d}.$Then by \eqref{5eq2}, we have 
\begin{equation}\label{5e13}
       \lim_{k \to \infty}U\ob{t_k}\ob{\e_0+\e_{\frac{n}{2}}}=\gamma\ob{\e_c+\e_d}.
   \end{equation}
   % Let $P$ be a permutation matrix corresponding to an automorphism of $\text{Cay}(\mathbb{Z}_n, S)$ satisfying $P\e_j=\e_{j+\frac{n}{2}}.$ Multiplying  
    Let $P$ be the permutation matrix corresponding to the automorphism of $\text{Cay}(\mathbb{Z}_n, S)$ satisfying $P\e_j=\e_{j+\frac{n}{2}}.$ Since $P(\e_0+\e_{\frac{n}{2}})=\e_0+\e_{\frac{n}{2}},$ and the limit in \eqref{5e13} is unique, we have
$P(\e_c+\e_d)=\e_c+\e_d.$
Therefore, $c$ and $d$ are antipodal vertices. Let $P'$ be the permutation matrix corresponding to the automorphism of $\text{Cay}(\mathbb{Z}_n, S)$ such that $P'\e_j=\e_{n-j} $ that fixes only $0$ and $\frac{n}{2}.$ Multiplying $P'$ on both sides of \eqref{5e13}, gives 
\begin{equation}\label{5e15}
P'(\e_c+\e_d)=\e_c+\e_d.
\end{equation}
The only antipodal vertices satisfying $\eqref{5e15},$ are $c=\frac{n}{4}$ and $d=\frac{3n}{4}.$ 
% Since 
% $a$ and $b$ are arbitrary antipodal vertices and $\text{Cay}(\mathbb{Z}_n, S)$
% is vertex transitive, the result follows.
\end{proof}
% \begin{rem}
%     Lemma \ref{5l9} and Lemma \ref{5l10} also holds for $\overline{C_n},$ where $n$ is even.
% \end{rem}
 The eigenvalues and eigenvectors of a cycle are well known. Suppose $\omega_n=\exp {\ob{\frac{2\pi i}{n}}}$
is
the primitive $n$-th root of unity. Then the eigenvalues of $C_n$ \cite{bro} are
$\lambda_l=2\cos\ob{{\frac{2l\pi}{n}}},$ where $ 0\leq l\leq n-1,   
$
and the corresponding eigenvectors are \begin{equation}\label{5eq7}\x_l =\tb{1\quad \omega_n^l\quad\cdots\quad\omega_n^{l(n-1)}}^T.
\end{equation}
The complement $\overline{C_n}$ has eigenvalues $\mu_0=n-\lambda_0-1,$ and $\mu_l=-1-\lambda_l,$ for $1\leq l \leq n-1,$ corresponding to the eigenvectors as given in \eqref{5eq7}.
A complete characterization of vertex PGST and pair PGST in $C_n$ is given in \cite[Theorem 13]{pal4} and \cite[Theorem 10]{pal10} as follows.


\begin{thm}\cite{pal4}\label{5t24} A cycle $C_n$ as well as its complement $\overline{C_n}$ admit pretty good state transfer
if and only if $n=2^k$ for $k\geq 2,$ and it occurs between every pair of antipodal vertices.
\end{thm}

\begin{thm}\cite{pal10}\label{5th25}
 A cycle $C_n$ on $n$ vertices admits pretty good pair state transfer if and only
if either $n = 2^k$ or $n = 2^kp,$ where $k$ is a positive integer and $p$ is an odd prime.
\end{thm}
We obtain the following as an immediate consequence of Theorem \ref{5t24}.
\begin{lem}\label{5t26}
Let $C_n$ be a cycle with two vertices $a$ and
$b,$ where $n=2^k,$ with $k\geq 2.$ Pretty good plus state transfer occurs in $C_n$ and its complement $\overline{C_n}$ between $\frac{1}{\sqrt{2}}\ob{\e_a+\e_b}$ and $\frac{1}{\sqrt{2}}\ob{\e_{a+\frac{n}{2}}+\e_{b+\frac{n}{2}}}$ if and only
if $a-b\neq \frac{n}{2}.$
\end{lem}
An even cycle $C_n$ can be realized as a double cover of $C_{\frac{n}{2}}$ as demonstrated in Figure \ref{5fig3}. Let $a$ and $b$ be antipodal vertices in $C_n.$ Corollary \ref{5cor3}(1) shows that a cycle $C_n$ admits PGST between $a$ and $b$ if and only if $C_{2n}$ admits PGST between $\frac{1}{\sqrt{2}}\ob{\e_{(0,a)}+\e_{(1,a)}}$ and $\frac{1}{\sqrt{2}}\ob{\e_{(0,b)}+\e_{(1,b)}}.$ Here, the vertices $(0,a)$ and $(1,a)$ are antipodal vertices in $C_{2n}.$ Including the cycle $C_4$ that admits plus PST between $\frac{1}{\sqrt{2}}\ob{\e_0+\e_2}$ and $\frac{1}{\sqrt{2}}\ob{\e_1+\e_3}$ at $\frac{\pi}{4}$ as shown in \cite[Theorem 6.5]{kim}, and using Lemma \ref{5l10} and Theorem \ref{5t24}, we observe the following.

\begin{lem}\label{5lemma5}
    Let $a$ and $b$ be antipodal vertices in $C_n.$ The cycle $C_n$ admits PGST between $\frac{1}{\sqrt{2}}\ob{\e_a+\e_b}$ and $\frac{1}{\sqrt{2}}\ob{\e_{a+\frac{n}{4}}+\e_{b+\frac{n}{4}}}$ if and only if $n=2^k$ with $k\geq 2.$ 
\end{lem}


% Using Corollary \ref{5cor3}(1) and Theorem \ref{5t24}, we find that if an even cycle $C_n$ exhibits plus PGST from $\frac{1}{\sqrt{2}}\ob{\e_a+\e_b},$ where $a$ and $b$ are antipodal vertices then $n=2^k$ with $k\geq 3.$
% Theorem \ref{5l7} together with Theorem \ref{5th25} indicates that  plus PGST may occur from a pair of non-antipodal in $C_n$ with $n=2^kp$, where $k$ is a positive integer and $p$ is an odd prime. The next result rules out this possibility.
Now we investigate plus PGST in $C_n,$ whenever $n$ has an odd prime factor.  
\begin{lem}\label{5l11}
   A cycle $C_n$ does not admit pretty good plus state transfer 
   % between $\frac{1}{\sqrt{2}}\ob{\e_0+\e_s}$ and $\frac{1}{\sqrt{2}}\ob{\e_{\frac{n}{2}}+\e_{\frac{n}{2}+s}},$ 
   % for $0<s<\frac{n}{2},$ 
   whenever $n=mp,$ where $m$ is a positive integer and $p$ is an odd prime.  
\end{lem}
\begin{proof}
If $m$ is odd, then the result follows from Lemma \ref{5l8}. Suppose that $m$ is even. Then Lemma \ref{5t26}
 and Lemma \ref{5lemma5} implies that if $C_n$ admits plus PGST from $\frac{1}{\sqrt{2}}\ob{\e_a+\e_b},$ then the vertices $a$ and $b$ must be non-antipodal, and PGST occurs between $\frac{1}{\sqrt{2}}\ob{\e_a+\e_b}$ and $\frac{1}{\sqrt{2}}\ob{\e_{a+\frac{n}{2}}+\e_{b+\frac{n}{2}}}.$ 
 Without loss of generaity let $a=0.$ The spectral decomposition of the adjacency matrix of $C_n$ gives $A(C_n)=\sum_{l=0}^{n-1}\lambda_l E_{\lambda_l}$  where $\lambda_l=2\cos{(\frac{2l\pi}{n})},$ $E_{\lambda_l}=\frac{1}{n}\x_l\x_l^*,$ and $\x_l(j)=\omega_n^{lj}$ for $0\leq j\leq n-1$ as given in \eqref{5eq7}. Consider the automorphism of $C_n$ with permutation matrix $P$ satisfying $P\e_j=\e_{j+\frac{n}{2}}.$
  Therefore, we have
  \[(P\x_l)(j)=\x_l\ob{j+\frac{n}{2}}=\omega_n^{l\ob{j+\frac{n}{2}}}=(-1)^l\x_l(j).\]
  This shows that $P\x_l=(-1)^l \x_l,$ and consequently, $PE_{\lambda_l}=(-1)^lE_{\lambda_l}.$
    Since $P^2=I,$ it follows that $P^T=P$ and $PE_{\lambda_l}=E_{\lambda_l}P.$ Hence,
    \begin{equation}\label{5eq22}
        E_{\lambda_l}\ob{\e_0+\e_b}=PE_{\lambda_l}\ob{e_{\frac{n}{2}}+\e_{b+\frac{n}{2}}}=(-1)^lE_{\lambda_l}\ob{e_{\frac{n}{2}}+\e_{b+\frac{n}{2}}}.
         \end{equation}
In the proof of \cite[Lemma 11]{pal4}, we observe that for an odd prime $p,$ the eigenvalues of $C_n$ satisfy the following identity.
         
%     For an odd prime $p,$ we have the following identity.
%  \[1+\omega_p+\omega_p^2+\ldots+\omega_p^{p-1}=0.\]
%  Therefore, considering the real parts of the above identity gives
%  \begin{equation}\label{5e22}
% 1+2\sum_{r=1}^\frac{p-1}{2}\cos{\ob{\frac{2r\pi}{p}}}=0.
%  \end{equation}
%  Let $m=2^k,$ thus $n=mp,$ and $\cos{\ob{\frac{2r\pi}{p}}}=\cos{\ob{\frac{2rm\pi}{n}}}.$ Multiplying both sides of \eqref{5e22} by $2\cos{\ob{\frac{2\pi}{n}}},$ yields
%  \begin{equation}\label{5e23}
%     \lambda_1+\sum_{r=1}^\frac{p-1}{2}\lambda_{mr+1}+ \sum_{r=1}^\frac{p-1}{2}\lambda_{mr-1}=0.
%  \end{equation}
%    Similarly, multiplying both sides of \eqref{5e22} by $2\cos{\ob{\frac{4\pi}{n}}},$ gives
%    \begin{equation}\label{5e24}
%     \lambda_2+\sum_{r=1}^\frac{p-1}{2}\lambda_{mr+2}+ \sum_{r=1}^\frac{p-1}{2}\lambda_{mr-2}=0.
%  \end{equation}
%  Subtracting \eqref{5e23} from \eqref{5e24}, we obtain
  \begin{equation}\label{5e25}
    (\lambda_2-\lambda_1)+\sum_{r=1}^\frac{p-1}{2}(\lambda_{mr+2}-\lambda_{mr+1})+ \sum_{r=1}^\frac{p-1}{2}(\lambda_{mr-2}-\lambda_{mr-1})=0.
 \end{equation}
 Suppose $C_n$ admits plus PGST between $\frac{1}{\sqrt{2}}\ob{\e_0+\e_b}$ and $\frac{1}{\sqrt{2}}\ob{\e_{\frac{n}{2}}+\e_{b+\frac{n}{2}}},$ then there exists a sequence $t_k\in\mathbb{R}$ and a complex number $\gamma$ of unit modulus such that 
\begin{equation}\label{5eq27}
      \gamma \sum_{l=0}^{n-1}E_{\lambda_l}\ob{\e_{\frac{n}{2}}+\e_{b+\frac{n}{2}}}=\gamma\ob{\e_{\frac{n}{2}}+\e_{b+\frac{n}{2}}}=\lim_{k \to \infty}\sum_{l=0}^{n-1}\exp{(it_k\lambda_l)}E_{\lambda_l}\ob{\e_0+\e_b}.
   \end{equation}
 Multiplying both sides of \eqref{5eq27} by $E_{\lambda_l}$ and using \eqref{5eq22}, we obtain  $\displaystyle\lim_{k \to \infty}\exp{(it_k\lambda_l)}=(-1)^l\gamma.$ 
 Therefore, 
$ \displaystyle\lim_{k \to \infty}\exp{(it_k(\lambda_{l+1}-\lambda_l))}=-1. $ Denoting the term on the left-hand side of \eqref{5e25} as $W,$ we obtain  
  $\displaystyle\lim_{k \to \infty}\exp{\ob{it_kW}}=-1.$
  This contradicts \eqref{5e25} as $W=0,$ implying no plus PGST between $\frac{1}{\sqrt{2}}\ob{\e_0+\e_b}$ and $\frac{1}{\sqrt{2}}\ob{\e_{\frac{n}{2}}+\e_{b+\frac{n}{2}}}$ in $C_n$.
\end{proof}
Combining Lemma \ref{5t26}, Lemma \ref{5lemma5} and Lemma \ref{5l11}
together with the case of $C_4$ and $C_8$ that exhibit plus PST, we obtain the following result on plus PGST in cycles.
\begin{thm}\label{5thm9}
    A cycle $C_n$ admits pretty good plus state transfer if and only if $n=2^k$ with $k\geq 2.$ Moreover, every plus state in $C_{2^k}$ exhibits pretty good plus state transfer. 
\end{thm}

Next, we investigate plus PGST in $\overline{C_n}$ whenever $n$ has an odd prime factor.

\begin{lem}\label{5lemm6}
There is no pretty good plus state transfer in the complement $\overline{C_n},$ whenever $n=mp,$ where $m$ is a positive integer and $p$ is an odd prime.
% Let $n$ be even with an odd prime factor. If $a$ and $b$ are a pair of antipodal vertices in $\overline{C_n},$ the complement of the cycle $C_n,$ then $\overline{C_n}$ does not admit pretty good plus state transfer from $\frac{1}{\sqrt{2}}\ob{\e_a+\e_b}.$
% Let $a$ and $b$ be a pair of antipodal vertices in $\overline{C_n},$ the complement of the cycle $C_n.$ The complement  $\overline{C_n}$ does not admit pretty good plus state transfer from $\frac{1}{\sqrt{2}}\ob{\e_a+\e_b},$ whenever $n$ has an odd prime factor.
    % If $n$ is divisible by $4,$ and has an odd prime factor $p,$ then the complement of the cycle graph $\overline{C_n}$ does not admit pretty good plus state transfer between $\frac{1}{\sqrt{2}}\ob{\e_0+\e_{\frac{n}{2}}}$ and $\frac{1}{\sqrt{2}}\ob{\e_{\frac{n}{4}}+\e_{\frac{3n}{4}}}.$
\end{lem}
\begin{proof}
If $m$ is odd, then the result follows from Lemma \ref{5l8}. Thus, assume that $m$ is even. 
The spectral decomposition of the adjacency matrix of
 $\overline{C_n}$ gives $A(\overline{C_n})=\displaystyle\sum_{l=0}^{n-1}\mu_l E_{\mu_l},$ where   
     $E_{\mu_l}=\frac{1}{n}\x_l\x_l^*,$ and $\x_l(j)=\omega_n^{lj},$ for $0\leq j\leq n-1,$ as given in \eqref{5eq7}.
\\\textit{Case I} $(a-b\neq\frac{n}{2})$:   If  $\overline{C_n}$ admits plus PGST from $\frac{1}{\sqrt{2}}\ob{\e_a+\e_b},$ where $a$ and 
$b$ are non-antipodal vertices, then by Lemma \ref{5l9}, it occurs between $\frac{1}{\sqrt{2}}\ob{\e_a+\e_b}$ and $\frac{1}{\sqrt{2}}\ob{\e_{a+\frac{n}{2}}+\e_{b+\frac{n}{2}}}.$ Without loss of generality, let $a=0$ and $0<b<\frac{n}{2}.$
% The complement of a cycle graph $\overline{C_n}$ does not admit pretty good plus state transfer between $\frac{1}{\sqrt{2}}\ob{\e_0+\e_s}$ and $\frac{1}{\sqrt{2}}\ob{\e_{\frac{n}{2}}+\e_{\frac{n}{2}+s}},$ for $0<s<\frac{n}{2},$ whenever $n=2^k p,$ where $k$ is a positive integer and $p$ is an odd prime.  
 Let $P$ be the permutation matrix corresponding to the automorphism of $\overline{C_n}$ defined by $P\e_j=\e_{j+\frac{n}{2}}.$ Using a similar approach as in the proof of Lemma \ref{5l11}, we obtain 
\begin{equation}\label{5eq33}
        E_{\mu_l}(e_0+\e_b)=(-1)^l E_{\mu_l}\ob{e_{\frac{n}{2}}+\e_{b+\frac{n}{2}}}.
         \end{equation}
         For an odd prime $p,$ using \eqref{5e25} and the fact that $\lambda_0=n-\mu_0-1,$ and $\lambda_l=-1-\mu_l,$ for $1 \leq l \leq n - 1,$ we have
  \begin{equation}\label{5eqn17}
    (\mu_2-\mu_1)+\sum_{r=1}^\frac{p-1}{2}(\mu_{mr+2}-\mu_{mr+1})+ \sum_{r=1}^\frac{p-1}{2}(\mu_{mr-2}-\mu_{mr-1})=\left\{ \begin{array}{rcl}
 2p, & \mbox{if} & m =2, \\ 
 0, & \mbox{if} & m\geq 3.
 \end{array}\right.
 \end{equation}
 If $\overline{C_n}$ admits plus PGST between $\frac{1}{\sqrt{2}}\ob{\e_0+\e_b}$ and $\frac{1}{\sqrt{2}}\ob{\e_{\frac{n}{2}}+\e_{b+\frac{n}{2}}},$ then a similar argument as in the proof of Lemma \ref{5l11} gives $\displaystyle\lim_{k \to \infty}\exp{(it_k\mu_l)}=(-1)^l\gamma.$ 
 It follows that
$ \displaystyle\lim_{k \to \infty}\exp{(it_k(\mu_{l+1}-\mu_l))}=-1. $ Let $W'$ be the expression on the left-hand side of \eqref{5eqn17}. Then  
  $\displaystyle\lim_{k \to \infty}\exp{\ob{it_kW'}}=-1,$ which contradicts \eqref{5eqn17} for $m\geq 3$ as $W'=0.$

Next we consider the case $m=2.$ Since $\mu_p=1,$ it follows that $\displaystyle\lim_{k \to \infty}\exp{(it_k)}=-\gamma.$ Since $W'=2p,$ we have $\displaystyle\lim_{k \to \infty}\exp(it_k 2p)=-1,$ and consequently $\gamma^{2p}=-1.$ For the eigenvalues $\mu_1$ and $\mu_{p-1},$ we obtain $\displaystyle\lim_{k \to \infty}\exp(it_k(\mu_1+\mu_{p-1}) )=-\gamma^2.$ Since $\mu_1+\mu_{p-1}=-2,$ we have $\gamma^4=-1,$ which is a contradiction to $\gamma^{2p}=-1$ as $p$ is an odd prime.
% Then proceeding in a similar way as discussed for $m=2,$ in the proof of \cite[Corollary 12]{pal4}, we arrive at a contradiction.
  % This is a contradiction as argued in the proof of Lemma \ref{5l11}. Hence, there is no plus PGST between $\frac{1}{\sqrt{2}}\ob{\e_0+\e_b}$ and $\frac{1}{\sqrt{2}}\ob{\e_{\frac{n}{2}}+\e_{b+\frac{n}{2}}}$ in $\overline{C_n}.$ 
 \\\textit{Case II} $(a-b=\frac{n}{2})$: If $\overline{C_n}$ admits plus PGST from $\frac{1}{\sqrt{2}}\ob{\e_a+\e_b},$ where $a$ and $b$ are antipodal vertices, then by Lemma \ref{5l10}, $n$ is divisible by $4$ and plus PGST occurs between $\frac{1}{\sqrt{2}}\ob{\e_a+\e_b}$ and $\frac{1}{\sqrt{2}}\ob{\e_{a+\frac{n}{4}}+\e_{b+\frac{n}{4}}}.$ Without loss of generality, let $a=0$ and $b=\frac{n}{2}.$
%    The spectral decomposition of the adjacency matrix of
% $\overline{C_n}$ gives $A(\overline{C_n})=\displaystyle\sum_{r=0}^{n-1}\mu_r E_{\mu_r},$ where   
%     $E_{\mu_r}=\frac{1}{n}\x_r\x_r^*,$ and $\x_r(j)=\omega_n^{rj},$ for $0\leq j\leq n-1,$ as given in \eqref{5eq7}.
Let $P'$ be the permutation matrix corresponding to the automorphism of $\overline{C_n}$ defined by $P'\e_j=\e_{j+\frac{n}{4}}.$
   Therefore, we have
   \[(P'\x_l)(j)=\x_l\ob{j+\frac{n}{4}}=\omega_n^{l\ob{j+\frac{n}{4}}}=(i)^l\x_l(j).\]
   This shows that $P'\x_l=(i)^l \x_l,$ and consequently, $P'E_{\mu_l}=(i)^lE_{\mu_l}.$
  Since $P'^T=P'^{-1}$ and $\x_l^*P'^T=(i)^{-l}\x_l^*,$ we have $\x_l^*P'=(i)^l\x_l^*.$ It follows that $P'E_{\mu_l}=E_{\mu_l}P'$ and 
% Since the adjacency matrix commute with $P$ and $E_{\mu_r}$ is a polynomial in $A,$ it follows that $PE_{\mu_r}=E_{\mu_r}P$ and
     \begin{equation}\label{5eq29}
         E_{\mu_l}\ob{e_0+\e_{\frac{n}{2}}}=P'E_{\mu_l}\ob{e_{\frac{n}{4}}+\e_{\frac{3n}{4}}}=(i)^lE_{\mu_l}\ob{e_{\frac{n}{4}}+\e_{\frac{3n}{4}}}.
          \end{equation}
          From the proof of \cite[Lemma 11]{pal4}, we observe that for an odd prime $p,$ 
\begin{equation}\label{5e22}
 1+2\sum_{r=1}^\frac{p-1}{2}\cos{\ob{\frac{2r\pi}{p}}}=0.
  \end{equation}
 Multiplying both sides of \eqref{5e22} by $2\cos{\ob{\frac{2\pi}{n}}},$ and $2\cos{\ob{\frac{6\pi}{n}}},$ yields
  \[
     \lambda_1+\sum_{r=1}^\frac{p-1}{2}\lambda_{mr+1}+ \sum_{r=1}^\frac{p-1}{2}\lambda_{mr-1}=0, \quad \text{and} \quad \lambda_3+\sum_{r=1}^\frac{p-1}{2}\lambda_{mr+3}+ \sum_{r=1}^\frac{p-1}{2}\lambda_{mr-3}=0,
  \]
respectively. Consequently, we have
  %   \begin{equation}\label{5e24}
  %    \lambda_2+\sum_{r=1}^\frac{p-1}{2}\lambda_{mr+2}+ \sum_{r=1}^\frac{p-1}{2}\lambda_{mr-2}=0.
  % \end{equation}
%  Subtracting \eqref{5e23} from \eqref{5e24}, we obtain
 \begin{equation}\label{5eqn21}
    (\lambda_3-\lambda_1)+\sum_{r=1}^\frac{p-1}{2}(\lambda_{mr+3}-\lambda_{mr+1})+ \sum_{r=1}^\frac{p-1}{2}(\lambda_{mr-3}-\lambda_{mr-1})=0.
 \end{equation}
          Substituting $\lambda_l=-1-\mu_l$ in \eqref{5eqn21}, we obtain
  \begin{equation}\label{5eq26}
    (\mu_3-\mu_1)+\sum_{r=1}^\frac{p-1}{2}(\mu_{mr+3}-\mu_{mr+1})+ \sum_{r=1}^\frac{p-1}{2}(\mu_{mr-3}-\mu_{mr-1})=0.
 \end{equation}
 Suppose $\overline{C_n}$ admits plus PGST between $\frac{1}{\sqrt{2}}\ob{\e_0+\e_\frac{n}{2}}$ and $\frac{1}{\sqrt{2}}\ob{\e_{\frac{n}{4}}+\e_{\frac{3n}{4}}},$ then there exists a sequence $t_k\in\mathbb{R}$ and a complex number $\gamma$ of unit modulus such that 
\begin{equation}\label{5eq31}
      \gamma \sum_{l=0}^{n-1}E_{\mu_l}\ob{\e_{\frac{n}{4}}+\e_{\frac{3n}{4}}}=\gamma\ob{\e_{\frac{n}{4}}+\e_{\frac{3n}{4}}}=\lim_{k \to \infty}\sum_{l=0}^{n-1}\exp{(it_k\mu_l)}E_{\mu_l}\ob{\e_0+\e_{\frac{n}{2}}}.
   \end{equation}
  Multiplying both sides of \eqref{5eq31} by $E_{\mu_l}$ and using \eqref{5eq29}, we obtain 
$\displaystyle\lim_{k \to \infty}\exp{(it_k\mu_l)}=(-i)^l\gamma.$
 Therefore, 
$ \displaystyle\lim_{k \to \infty}\exp{(it_k(\mu_{l+2}-\mu_l))}=-1. $ Denoting the term on the left-hand side of \eqref{5eq26} as $W'',$ we arrive at a contradiction
  $\displaystyle\lim_{k \to \infty}\exp{\ob{it_kW''}}=-1$
as $W''=0.$ Hence, no plus PGST occurs between $\frac{1}{\sqrt{2}}\ob{\e_0+\e_{\frac{n}{2}}}$ and $\frac{1}{\sqrt{2}}\ob{\e_{\frac{n}{4}}+\e_{\frac{3n}{4}}}.$ 
\end{proof}
% Now we provide a complete characterization of plus PGST in $\overline{C_n}$ from a pair of antipodal vertices.
Now we determine the conditions under which plus PGST occurs from a pair of antipodal vertices in the complement $\overline{C_n}.$
\begin{lem}\label{5lem6}
Let $a$ and $b$ be a pair of antipodal vertices in $\overline{C_n},$ the complement of the cycle $C_n.$
The complement $\overline{C_n}$ admits pretty good plus state transfer from $\frac{1}{\sqrt{2}}\ob{\e_a+\e_b}$ if and only if $n=2^k$ with $k\geq 3.$ 
\end{lem}
\begin{proof}
Without loss of generality let $a=0$ and $b=\frac{n}{2}.$
    If $n=2^k,$ with $k\geq 3,$ then $C_n$ as well as its complement admit vertex PGST with respect to the same sequence $t_l$ in $2\pi\mathbb{Z}$ \cite[Theorem 7]{pal4}. Since $C_n$ can be realized as a double cover of $C_{\frac{n}{2}},$ by Corollary \ref{5cor3}(1), the cycle $C_{2^k}$ admits plus PGST between $\frac{1}{\sqrt{2}}\ob{\e_0+\e_{\frac{n}{2}}}$ and $\frac{1}{\sqrt{2}}\ob{\e_{\frac{n}{4}}+\e_{\frac{3n}{4}}},$ whenever $k\geq 4$ with respect to the same sequence. As argued in Remark \ref{5remark2}, since $2^kt_l\in2\pi\mathbb{Z},$ $\overline{C_{2^k}}$
admits plus PGST between $\frac{1}{\sqrt{2}}\ob{\e_0+\e_{\frac{n}{2}}}$ and $\frac{1}{\sqrt{2}}\ob{\e_{\frac{n}{4}}+\e_{\frac{3n}{4}}},$ whenever $k\geq 4.$ For $k=3,$ the cycle $C_8$ has plus PST between $\frac{1}{\sqrt{2}}\ob{\e_0+\e_4}$ and $\frac{1}{\sqrt{2}}\ob{\e_2+\e_6}$ at $\frac{\pi}{2}.$ Hence, by Theorem \ref{5t15}, plus PST occurs  between those states in $\overline{C_8}$.

Conversely, if $\overline{C_n}$ admits plus PGST from $\frac{1}{\sqrt{2}}\ob{\e_0+\e_\frac{n}{2}},$ then by Lemma \ref{5l10}, it occurs between $\frac{1}{\sqrt{2}}\ob{\e_0+\e_\frac{n}{2}}$ and $\frac{1}{\sqrt{2}}\ob{\e_\frac{n}{4}+\e_\frac{3n}{4}}.$  Using the observation above and Lemma \ref{5lemm6}, we find that the only possible case is when $n=2^k$ with $k\geq 2.$
The graph $\overline{C_4}$ does not admit plus PGST between $\frac{1}{\sqrt{2}}\ob{\e_0+\e_2}$ and $\frac{1}{\sqrt{2}}\ob{\e_1+\e_3}$ as $\frac{1}{\sqrt{2}}\ob{\e_0+\e_2}$ is a fixed state in $\overline{C_4}.$ Hence, the result follows.   
\end{proof}
 % Next, we check the existence of plus PGST between $\frac{1}{\sqrt{2}}\ob{\e_0+\e_{\frac{n}{2}}}$ and $\frac{1}{\sqrt{2}}\ob{\e_{\frac{n}{4}}+\e_{\frac{3n}{4}}},$ whenever $n$ is divisible by $4$ and has an odd prime factor $p.$  
% \begin{lem}
%     If $n$ is divisible by $4,$ and has an odd prime factor $p,$ then the complement of the cycle graph $\overline{C_n}$ does not admit pretty good plus state transfer between $\frac{1}{\sqrt{2}}\ob{\e_0+\e_{\frac{n}{2}}}$ and $\frac{1}{\sqrt{2}}\ob{\e_{\frac{n}{4}}+\e_{\frac{3n}{4}}}.$
% \end{lem}
% \begin{proof}
%     Let $A(\overline{C_n})=\displaystyle\sum_{r=0}^{n-1}\mu_r E_{\mu_r}$ be the spectral decomposition of adjacency matrix of $\overline{C_n},$ where 
%     % $\mu_0=n-3,$ and $\mu_r=-1-\lambda_r,$ where $\lambda_r$ are eigenvalues of $C_n$ for $1\leq r\leq n-1.$  
%     $E_{\mu_r}=\frac{1}{n}\x_r\x_r^*,$ and $\x_r(j)=\omega_n^{rj}$ as given in \eqref{5eq7}. Let $P$ be the permutation matrix corresponding to an automorphism of $\overline{C_n}$ defined by $P\e_j=\e_{j+\frac{n}{4}}.$
%    Therefore, we have
%    \[P\x_r(j)=\x_r\ob{j+\frac{n}{4}}=\omega_n^{r\ob{j+\frac{n}{4}}}=(i)^r\x_r(j).\]
%    This shows that $P\x_r=(i)^r \x_r,$ and consequently, $PE_{\mu_r}=(i)^rE_{\mu_r}.$
% Since the adjacency matrix commute with $P$ and $E_{\mu_r}$ is a polynomial in $A,$ it follows that $PE_{\mu_r}=E_{\mu_r}P$ and
%      \begin{equation}\label{5eq29}
%          E_{\mu_r}\ob{e_0+\e_{\frac{n}{2}}}=PE_{\mu_r}\ob{e_{\frac{n}{4}}+\e_{\frac{3n}{4}}}=(i)^rE_{\mu_r}\ob{e_{\frac{n}{4}}+\e_{\frac{3n}{4}}}.
%           \end{equation}
%    Let $n=mp,$ for some positive integer $m.$ For an odd prime $p,$ using a similar argument as given in the proof of Lemma \ref{5l11} and using the fact that $\lambda_r=-1-\mu_r,$ we obtain
%   \begin{equation}\label{5eq26}
%     (\mu_3-\mu_1)+\sum_{r=1}^\frac{p-1}{2}(\mu_{mr+3}-\mu_{mr+1})+ \sum_{r=1}^\frac{p-1}{2}(\mu_{mr-3}-\mu_{mr-1})=0.
%  \end{equation}
%  Suppose $\overline{C_n}$ admits plus PGST between $\frac{1}{\sqrt{2}}\ob{\e_0+\e_\frac{n}{2}}$ and $\frac{1}{\sqrt{2}}\ob{\e_{\frac{n}{4}}+\e_{\frac{3n}{4}}},$ then there exists a sequence $t_k\in\mathbb{R}$ and a complex number $\gamma$ of unit modulus such that 
% \begin{equation}\label{5eq31}
%        \lim_{k \to \infty}\sum_{r=0}^{n-1}\exp{(it_k\mu_r)}E_{\mu_r}\ob{\e_0+\e_{\frac{n}{2}}}=\gamma\ob{\e_{\frac{n}{4}}+\e_{\frac{3n}{4}}}.
%    \end{equation}
%  From \eqref{5eq29} and \eqref{5eq31}, we obtain 
%  \begin{equation}\label{5e27}
%        \lim_{k \to \infty}\exp{(it_k\mu_r)}=(-i)^r\gamma. 
%   \end{equation} 
%  Therefore, 
% $ \displaystyle\lim_{k \to \infty}\exp{(it_k(\mu_{r+2}-\mu_r))}=-1. $ Denoting the term on the left-hand side of \eqref{5eq26} as $L,$ we obtain 
%   \[\lim_{k \to \infty}\exp{\ob{it_kL}}=-1.\]
%   This contradicts \eqref{5eq26}, implying no plus PGST between $\frac{1}{\sqrt{2}}\ob{\e_0+\e_{\frac{n}{2}}}$ and $\frac{1}{\sqrt{2}}\ob{\e_{\frac{n}{4}}+\e_{\frac{3n}{4}}}.$ 
% \end{proof}
 % Since $\overline{C_n}$ has pair PGST if and only if $C_n$ has pair PGST, by Theorem \ref{5l7}, it may be possible to have plus PGST in $\overline{C_n},$ whenever $n=2^kp,$ for a positive integer $k$ and an odd prime $p.$ However, this is ruled out by the following observation.
% The following result identifies $n$ for which $\overline{C_n}$ does not admit plus PGST.
 % \begin{lem}\label{5lemma9}
% The complement $\overline{C_n}$ does not admit pretty good plus state transfer, whenever $n=mp,$ where $m$ is a positive integer and $p$ is an odd prime. 
 % \end{lem} 
% \begin{proof} 
% If $m$ is odd, then the result follows from Lemma \ref{5l8}. Thus, assume that $m$ is even. By Lemma \ref{5lemm6}, plus PGST does not occur from a pair of antipodal vertices, and hence it remains to consider the case where 
% $a$ and 
% $b$ are non-antipodal vertices in $\overline{C_n}.$ If  $\overline{C_n}$ admits plus PGST from $\frac{1}{\sqrt{2}}\ob{\e_a+\e_b},$ then by Lemma \ref{5l9}, it occurs between $\frac{1}{\sqrt{2}}\ob{\e_a+\e_b}$ and $\frac{1}{\sqrt{2}}\ob{\e_{a+\frac{n}{2}}+\e_{b+\frac{n}{2}}}.$ Without loss of generality, let $a=0$ and $0<b<\frac{n}{2}.$
% % The complement of a cycle graph $\overline{C_n}$ does not admit pretty good plus state transfer between $\frac{1}{\sqrt{2}}\ob{\e_0+\e_s}$ and $\frac{1}{\sqrt{2}}\ob{\e_{\frac{n}{2}}+\e_{\frac{n}{2}+s}},$ for $0<s<\frac{n}{2},$ whenever $n=2^k p,$ where $k$ is a positive integer and $p$ is an odd prime.  
%  Let $P$ be the permutation matrix corresponding to an automorphism of $\overline{C_n}$ defined by $P\e_j=\e_{j+\frac{n}{2}}.$ Using a similar approach as in the proof of Lemma \ref{5l11}, we obtain 
% \begin{equation}\label{5eq33}
%         E_{\mu_r}(e_0+\e_b)=(-1)^r E_{\mu_r}\ob{e_{\frac{n}{2}}+\e_{b+\frac{n}{2}}}.
%          \end{equation}
%          For an odd prime $p,$ using a similar argument as given in the proof of Lemma \ref{5l11} and using the fact that $\lambda_r=-1-\mu_r,$ we have
%   \begin{equation}\label{5eq30}
%     (\mu_2-\mu_1)+\sum_{r=1}^\frac{p-1}{2}(\mu_{mr+2}-\mu_{mr+1})+ \sum_{r=1}^\frac{p-1}{2}(\mu_{mr-2}-\mu_{mr-1})=0.
%  \end{equation}
%  A similar argument as presented in the proof of Lemma \ref{5l11}, contradicts the existence of plus PGST between $\frac{1}{\sqrt{2}}\ob{\e_0+\e_b}$ and $\frac{1}{\sqrt{2}}\ob{\e_{\frac{n}{2}}+\e_{b+\frac{n}{2}}}$ in $\overline{C_n}.$ 
 % \end{proof}

 Now we provide a complete characterization of plus PGST in the cycle $C_n$ and its complement $\overline{C_n}$ by combining Theorem \ref{5thm9}, Lemma \ref{5lemm6} and Lemma \ref{5lem6}. 
 \begin{thm}\label{5th10}
      A cycle $C_n$ as well as its complement $\overline{C_n}$ admit pretty good plus state transfer if and only if $n=2^k$ with $k\geq 2.$ Moreover, for $k\geq 2,$ every plus state in $C_{2^k}$ and for $k\geq 3,$ every plus state in $\overline{C_{2^k}}$ exhibit pretty good plus state transfer.
 \end{thm}
% \begin{lem}\label{5l7}
% Let a graph $G$ admit pretty good plus state transfer between $\frac{1}{\sqrt{2}}\ob{\e_a+\e_b}$ and $\frac{1}{\sqrt{2}}\ob{\e_c+\e_d}.$ If $P$ is a permutation matrix corresponding to an automorphism of $G$ such that $P\e_a=\e_b$ and $P\e_b\neq \e_a,$ then the graph $G$ admits pretty good pair state transfer.

% \end{lem}
% \begin{proof}
%  Suppose the graph $G$ admits plus PGST between $\frac{1}{\sqrt{2}}\ob{\e_a+\e_b}$ and $\frac{1}{\sqrt{2}}\ob{\e_c+\e_d},$ then by \eqref{5eq2}, we have 
% % Then there is a sequence of time $t_k\in\mathbb{R}$ and a unit complex number $\gamma$ such that  
%     \begin{equation}\label{5eqn7}
%         \lim_{k \to \infty}U_G\ob{t_k}\ob{\e_a+\e_b}=\gamma\ob{\e_c+\e_d}.
%     \end{equation}
%    Multiplying $P$ on both sides of \eqref{5eqn7} and subtracting the resulting equation from \eqref{5eqn7}, yields
%    \begin{equation}\label{5eqn8}
%        \lim_{k \to \infty}U_G\ob{t_k}\ob{\e_a-P\e_b}=\gamma\ob{\e_c+\e_d-P\e_c-P\e_d}.
%    \end{equation}
%    Since $U_G(t)$ is unitary, \eqref{5eqn8} implies that $\left\|\frac{1}{\sqrt{2}}\ob{\e_b-P\e_b}\right\|=\left\|\frac{1}{\sqrt{2}}\ob{\e_c+\e_d-P\e_c-P\e_d}\right\|.$ This equality is possible only if $\frac{1}{\sqrt{2}}\ob{\e_c+\e_d-P\e_c-P\e_d}$ represents a pair state. 
% \end{proof}

% We investigate the existence of plus PGST in cycles with respect to the adjacency matrix. 
% The eigenvalues and eigenvectors of a cycle are well known. Suppose $\omega_n=\exp {\ob{\frac{2\pi i}{n}}}$
% is
% the primitive $n$-th root of unity. Then the eigenvalues of $C_n$ are
% $\lambda_r=2\cos\ob{{\frac{2r\pi}{n}}},$ where $ 0\leq r\leq n-1,   
% $
% and the corresponding eigenvectors are \begin{equation}\label{5eq7}\x_r =\tb{1\quad \omega_n^r\quad\cdots\quad\omega_n^{r(n-1)}}^T.
% \end{equation}
% The complement $\overline{C_n}$ has eigenvalues $\mu_0=n-3,$ and $\mu_r=-1-\lambda_r,$ for $1\leq r \leq n-1,$ corresponding to the eigenvectors as given in \eqref{5eq7}.
% A complete characterization of vertex PGST and pair PGST in $C_n$ is given in \cite[Theorem 13]{pal4} and \cite[Theorem 10]{pal10} as follows.


% \begin{thm}\cite{pal4}\label{5t24} A cycle $C_n$ as well as its complement $\overline{C_n}$ admit pretty good state transfer
% if and only if $n=2^k$ for $k\geq 2,$ and it occurs between every pair of antipodal vertices.
% \end{thm}

% \begin{thm}\cite{pal10}\label{5th25}
%  A cycle $C_n$ on $n$ vertices admits pretty good pair state transfer if and only
% if either $n = 2^k$ or $n = 2^kp,$ where $k$ is a positive integer and $p$ is an odd prime.
% \end{thm}
% We obtain the following as an immediate consequence of Theorem \ref{5t24}.
% \begin{thm}\label{5t26}
% Let $C_n$ be a cycle with two vertices $a$ and
% $b,$ where $n=2^k,$ with $k\geq 2.$ Pretty good plus state transfer occurs in $C_n$ and its complement $\overline{C_n}$ between $\frac{1}{\sqrt{2}}\ob{\e_a+\e_b}$ and $\frac{1}{\sqrt{2}}\ob{\e_{a+\frac{n}{2}}+\e_{b+\frac{n}{2}}}$ if and only
% if $a-b\neq \frac{n}{2}.$
% \end{thm}

%  The following result shows the existence of plus PGST in cycles implies the existence of pair PGST.

% \begin{lem}\label{5l7}
% Let the cycle $C_n$ exhibit pretty good plus state transfer between $\frac{1}{\sqrt{2}}\ob{\e_a+\e_b}$ and $\frac{1}{\sqrt{2}}\ob{\e_c+\e_d}.$ If $a$ and $b$ are not antipodal vertices, then $C_n$ admits pretty good pair state transfer.
% \end{lem}
% \begin{proof}
% Let $a$ and $b$ be non-antipodal vertices in $C_n.$ If the cycle $C_n$ exhibits plus PGST between $\frac{1}{\sqrt{2}}\ob{\e_a+\e_b}$ and $\frac{1}{\sqrt{2}}\ob{\e_c+\e_d},$ then by \eqref{5eq2}, we have 
% % Then there is a sequence of time $t_k\in\mathbb{R}$ and a unit complex number $\gamma$ such that  
%     \begin{equation}\label{5e11}
%         \lim_{k \to \infty}U_G\ob{t_k}\ob{\e_a+\e_b}=\gamma\ob{\e_c+\e_d}.
%     \end{equation}
%     Let $P$ be a permutataion matrix of an automorphism of $C_n$ such that $P\e_a=\e_b$ and $P\e_b\neq\e_a.$
%    Multiplying $P$ on both sides of \eqref{5e11} and subtracting the resulting equation from \eqref{5e11}, we obtain
%    \begin{equation}\label{5e12}
%        \lim_{k \to \infty}U_G\ob{t_k}\ob{\e_a-P\e_b}=\gamma\ob{\e_c+\e_d-P\e_c-P\e_d}.
%    \end{equation}
%    Since $U_G(t)$ is unitary, \eqref{5e12} implies that $\left\|\frac{1}{\sqrt{2}}\ob{\e_b-P\e_b}\right\|=\left\|\frac{1}{\sqrt{2}}\ob{\e_c+\e_d-P\e_c-P\e_d}\right\|.$ This equality is possible only if $\frac{1}{\sqrt{2}}\ob{\e_c+\e_d-P\e_c-P\e_d}$ represents a pair state. 
% \end{proof}
% \begin{rem}
% Let $a$ and $b$ be two vertices of a graph $G.$ If $P$ is a permutation matrix corresponding to an automorphism of $G$ such that $P\e_a=\e_b$ and $P\e_b\neq \e_a,$ then Lemma \ref{5l7} holds for $G.$ In particular, Lemma \ref{5l7} holds for the circulant graphs.
% \end{rem}
% % Note that Lemma \ref{5l7} also holds for the complement $\overline{C_n}.$
% Since an odd cycle and its complement have no antipodal vertices, Lemma \ref{5l7} implies that if they admit plus PGST, then they must exhibit pair PGST. However, this is ruled out by \cite [Theorem 10]{pal10}, leading to the following result.   
% \begin{lem}\label{5l8}
%     The odd cycle and its complement do not admit pretty good plus state transfer.
% \end{lem}
%  The existence of plus PGST from $\frac{1}{\sqrt{2}}\ob{\e_a+\e_b}$ in an even cycle $C_n$ when $a$ and $b$ are not antipodal vertices is shown in the following result. 
% \begin{lem}\label{5l9}
%     Let an even cycle $C_n$ exhibit pretty good plus state transfer between $\frac{1}{\sqrt{2}}\ob{\e_a+\e_b}$ and $\frac{1}{\sqrt{2}}\ob{\e_c+\e_d}.$ If $a$ and $b$ are not antipodal vertices, then $\e_c+\e_d=\e_{a+\frac{n}{2}}+\e_{b+\frac{n}{2}}.$
% \end{lem}
% \begin{proof}
%      Without loss of generality, let $a=0$ and $b<\frac{n}{2}.$ Suppose $C_n$ admits plus PGST between $\frac{1}{\sqrt{2}}\ob{\e_0+\e_b}$ and $\frac{1}{\sqrt{2}}\ob{\e_c+\e_d}.$ Then \eqref{5eq2} gives
%      % there exists a sequence $t_k\in\mathbb{R}$ and a complex number $\gamma$ of unit modulus such that 
% \begin{equation}\label{5eq13}
%        \lim_{k \to \infty}U_G\ob{t_k}\ob{\e_0+\e_b}=\gamma\ob{\e_c+\e_d}.
%    \end{equation}
%  Let $P$ be a permutation matrix corresponding to an automorphism of $C_n$ satisfying $P\e_j=\e_{-j}$ which fixes only $0$ and $\frac{n}{2}.$ Multiplying $P$ on both sides of \eqref{5eq13} and subtracting the resulting equation from \eqref{5eq13}, we obtain
%    \begin{equation}\label{5eq14}
%        \lim_{k \to \infty}U_G\ob{t_k}\ob{\e_b-P\e_b}=\gamma\ob{\e_c+\e_d-P\e_c-P\e_d}.
%    \end{equation}
% This equality holds only if either one of $c$ or $d$ lies in the set $\{0, \frac{n}{2}\}.$ Therefore, without loss of generality, assume that $c\in\{0,\frac{n}{2}\}.$
% In the case, when $c=0$,  \eqref{5eq13} becomes
% \begin{equation}\label{5eq15}
%        \lim_{k \to \infty}U_G\ob{t_k}\ob{\e_0+\e_b}=\gamma\ob{\e_0+\e_d}.
%    \end{equation}
%    Let $P'$ be a permutation matrix  defined by $P'\e_j=\e_{2b-j}.$ Then  $P'$ fixes only $b$ and $ b+\frac{n}{2}.$
%    Multiplying $P'$ on both sides of \eqref{5eq15} and subtracting the resulting equation from \eqref{5eq15}, yields
%    \begin{equation}\label{5eq16}
%        \lim_{k \to \infty}U_G\ob{t_k}\ob{\e_0-P'\e_0}=\gamma\ob{\e_0+\e_d-P'\e_0-P'\e_d}.
%    \end{equation}
% This equality holds only when $d=b+\frac{n}{2}.$ Hence, \eqref{5eq15} becomes
% \begin{equation}\label{5eq17}
%        \lim_{k \to \infty}U_G\ob{t_k}\ob{\e_0+\e_b}=\gamma\ob{\e_0+\e_{b+\frac{n}{2}}}.
%    \end{equation}
%    Next, let $P''$ be a permutation matrix  satisfying $P''\e_0=\e_b.$ Multiplying $P''$ on both sides of \eqref{5eq17} and subtracting the resulting equation from \eqref{5eq17}, gives
%    \begin{equation}\label{5eq18}
%        \lim_{k \to \infty}U_G\ob{t_k}\ob{\e_0-\e_{2b}}=\gamma\ob{\e_0+\e_{b+\frac{n}{2}}-\e_b-\e_{2b+\frac{n}{2}}}.
%    \end{equation} 
%    Since $\e_0\neq \e_b,$ $ \e_{b+\frac{n}{2}}\neq \e_b$ and $\e_{b+\frac{n}{2}}\neq \e_{2b+\frac{n}{2}},$ \eqref{5eq18} holds only when $\e_0=\e_{2b+\frac{n}{2}},$ which implies $b=\frac{n}{4}.$ Consequently, pair PGST occurs from a pair of antipodal vertices, contradicting  \cite[Lemma 7]{pal10}. Therefore, $c$ must be $\frac{n}{2}.$
%  If $c=\frac{n}{2},$ then multiplying $P'$ on both sides of \eqref{5eq13} and subtracting the resulting equation from \eqref{5eq13} implies that
% $d=b+\frac{n}{2}.$

%  Since $a$ and $b$ are arbitrary non-antipodal vertices and $C_n$ is vertex transitive, the result follows.
% \end{proof}

% We next characterize plus PGST when 
% $a$ and $b$ are antipodal vertices of the even cycle.
% \begin{lem}\label{5l10}
%     Let $a$ and $b$ be antipodal vertices in an even cycle $C_n.$ If there is pretty good plus state transfer from $\frac{1}{\sqrt{2}}\ob{\e_a+\e_b},$ then $n$ is divisible by $4$ and pretty good plus state transfer occurs between $\frac{1}{\sqrt{2}}\ob{\e_a+\e_b}$ and $\frac{1}{\sqrt{2}}\ob{\e_{a+\frac{n}{4}}+\e_{b+\frac{n}{4}}}.$  
% \end{lem}

% \begin{proof}
%     Without loss of generality, let $a=0$ and $b=\frac{n}{2}.$ Suppose $C_n$ admits plus PGST between $\frac{1}{\sqrt{2}}\ob{\e_0+\e_{\frac{n}{2}}}$ and $\frac{1}{\sqrt{2}}\ob{\e_c+\e_d}.$ Then there exists a sequence $t_k\in\mathbb{R}$ and a complex number $\gamma$ of unit modulus such that 
% \begin{equation}\label{5e13}
%        \lim_{k \to \infty}U_G\ob{t_k}\ob{\e_0+\e_{\frac{n}{2}}}=\gamma\ob{\e_c+\e_d}.
%    \end{equation}
%    Let $P$ be a permutation matrix corresponding to an automorphism of $C_n$ satisfying $P\e_j=\e_{j+\frac{n}{2}}.$ Multiplying  
   
%  Since a permutation matrix $P$ corresponding to an automorphism of $C_n$ satisfying $P\e_j=\e_{j+\frac{n}{2}},$ gives $P(\e_0+\e_{\frac{n}{2}})=\e_0+\e_{\frac{n}{2}},$ and the limit in \eqref{5e13} is unique, we have
%  \begin{equation}\label{5e14}
% P(\e_c+\e_d)=\e_c+\e_d=\e_{c+\frac{n}{2}}+\e_{d+\frac{n}{2}}.
% \end{equation}
% Therefore, $c$ and $d$ are antipodal vertices. Let $P'$ be a permutation matrix corresponding to an automorphism of $C_n$ such that $P'\e_j=\e_{-j} $ that fixes only $0$ and $\frac{n}{2}.$ Multiplying $P'$ on both sides of \eqref{5e13}, gives 
% \begin{equation}\label{5e15}
% P'(\e_c+\e_d)=\e_c+\e_d.
% \end{equation}
% From \eqref{5e14} and $\eqref{5e15},$ we obtain $c=\frac{n}{4}$ and $d=\frac{3n}{4}.$ Since 
% $a$ and $b$ are arbitrary antipodal vertices and $C_n$
% is vertex transitive, the result follows.
% \end{proof}
% \begin{rem}
%     Lemma \ref{5l9} and Lemma \ref{5l10} also holds for $\overline{C_n},$ where $n$ is even.
% \end{rem}
% An even cycle $C_n$ can be seen as a double cover of $C_{\frac{n}{2}}$ (see Figure \ref{5fig3}). Using Corollary \ref{5cor3} (1) and Theorem \ref{5t24}, we find that if an even cycle $C_n$ exhibits plus PGST from $\frac{1}{\sqrt{2}}\ob{\e_a+\e_b},$ where $a$ and $b$ are antipodal vertices then $n=2^k$ with $k\geq 2.$ Lemma \ref{5l7} together with Theorem \ref{5th25} indicates that $C_n$ may admit plus PGST whenever $n=2^kp$, where $k$ is a positive integer and $p$ is an odd prime. The next result together with Lemma \ref{5l9} rules out this possibility.

% \begin{lem}\label{5l11}
%    A cycle $C_n$ does not admit pretty good plus state transfer between $\frac{1}{\sqrt{2}}\ob{\e_0+\e_s}$ and $\frac{1}{\sqrt{2}}\ob{\e_{\frac{n}{2}}+\e_{\frac{n}{2}+s}},$ for $0<s<\frac{n}{2},$ whenever $n=2^k p,$ where $k$ is a positive integer and $p$ is an odd prime.  
% \end{lem}
% \begin{proof}
% Let $A(C_n)=\sum_{r=0}^{n-1}\lambda_r E_{\lambda_r}$ be the spectral decomposition of adjacency matrix of $C_n,$ where $\lambda_r=2\cos{(\frac{2r\pi}{n})},$ $E_{\lambda_r}=\frac{1}{n}\x_r\x_r^*,$ and $\x_r(j)=\omega_n^{rj}$ as given in \eqref{5eq7}. Let $P$ be the permutation matrix corresponding to an automorphism of $C_n$ defined by $P\e_j=\e_{j+\frac{n}{2}}.$
%   Therefore, we have
%   \[P\x_r(j)=\x_r\ob{j+\frac{n}{2}}=\omega_n^{r\ob{j+\frac{n}{2}}}=(-1)^r\x_r(j).\]
%   This shows that $P\x_r=(-1)^r \x_r,$ and consequently, $PE_{\lambda_r}=(-1)^rE_{\lambda_r}.$
%     Since the adjacency matrix commute with $P$ and $E_{\lambda_r}$ is a polynomial in $A,$ it follows that $PE_{\lambda_r}=E_{\lambda_r}P$ and
%     \begin{equation}\label{5eq22}
%         E_{\lambda_r}\ob{e_{\frac{n}{2}}+\e_{\frac{n}{2}+s}}=PE_{\lambda_r}(e_0+\e_s)=(-1)^rE_{\lambda_r}(e_0+\e_s).
%          \end{equation}
%     For an odd prime $p,$ we have the following identity.
%  \[1+\omega_p+\omega_p^2+\ldots+\omega_p^{p-1}=0.\]
%  Therefore, considering the real parts of the above identity gives
%  \begin{equation}\label{5e22}
% 1+2\sum_{r=1}^\frac{p-1}{2}\cos{\ob{\frac{2r\pi}{p}}}=0.
%  \end{equation}
%  Let $m=2^k,$ thus $n=mp,$ and $\cos{\ob{\frac{2r\pi}{p}}}=\cos{\ob{\frac{2rm\pi}{n}}}.$ Multiplying both sides of \eqref{5e22} by $2\cos{\ob{\frac{2\pi}{n}}},$ yields
%  \begin{equation}\label{5e23}
%     \lambda_1+\sum_{r=1}^\frac{p-1}{2}\lambda_{mr+1}+ \sum_{r=1}^\frac{p-1}{2}\lambda_{mr-1}=0.
%  \end{equation}
%    Similarly, multiplying both sides of \eqref{5e22} by $2\cos{\ob{\frac{4\pi}{n}}},$ gives
%    \begin{equation}\label{5e24}
%     \lambda_2+\sum_{r=1}^\frac{p-1}{2}\lambda_{mr+2}+ \sum_{r=1}^\frac{p-1}{2}\lambda_{mr-2}=0.
%  \end{equation}
%  Subtracting \eqref{5e23} from \eqref{5e24}, we obtain
%   \begin{equation}\label{5e25}
%     (\lambda_2-\lambda_1)+\sum_{r=1}^\frac{p-1}{2}(\lambda_{mr+2}-\lambda_{mr+1})+ \sum_{r=1}^\frac{p-1}{2}(\lambda_{mr-2}-\lambda_{mr-1})=0.
%  \end{equation}
%  Suppose $C_n$ admits plus PGST between $\frac{1}{\sqrt{2}}\ob{\e_0+\e_s}$ and $\frac{1}{\sqrt{2}}\ob{\e_{\frac{n}{2}}+\e_{\frac{n}{2}+s}},$ then there exists a sequence $t_k\in\mathbb{R}$ and a complex number $\gamma$ of unit modulus such that 
% \begin{equation}\label{5e27}
%        \lim_{k \to \infty}\sum_{r=0}^{n-1}\exp{(it_k\lambda_r)}E_{\lambda_r}\ob{\e_0+\e_s}=\gamma\ob{\e_{\frac{n}{2}}+\e_{\frac{n}{2}+s}}.
%    \end{equation}
%  From \eqref{5eq22} and \eqref{5e27}, we obtain 
%  \begin{equation}\label{5e27}
%        \lim_{k \to \infty}\exp{(it_k\lambda_r)}=(-1)^r\gamma. 
%   \end{equation} 
%  Therefore, 
% $ \displaystyle\lim_{k \to \infty}\exp{(it_k(\lambda_{r+1}-\lambda_r))}=-1. $ Consequently, 
%   \[\lim_{k \to \infty}\exp{\ob{it_k\ob{(\lambda_2-\lambda_1)+\sum_{r=1}^\frac{p-1}{2}(\lambda_{mr+2}-\lambda_{mr+1})+ \sum_{r=1}^\frac{p-1}{2}(\lambda_{mr-2}-\lambda_{mr-1})}}}=-1.\]
%   This contradicts \eqref{5e25}, implying no plus PGST between $\frac{1}{\sqrt{2}}\ob{\e_0+\e_s}$ and $\frac{1}{\sqrt{2}}\ob{\e_{\frac{n}{2}}+\e_{\frac{n}{2}+s}}$ in $C_n$.
% \end{proof}
% Including the case of $C_4$ and $C_8$ that exhibit plus PST, the following result gives a complete characterization of plus PGST in cycles.
% \begin{thm}
%     A cycle $C_n$ admits pretty good plus state transfer if and only if $n=2^k$ with $k\geq 2.$ Moreover, for $k\geq 3,$ every plus state in $C_{2^k}$ exhibits pretty good plus state transfer. 
% \end{thm}

% Next, we investigate the existence of plus PGST in $\overline{C_n},$ whenever $n$ is even. If $n=2^k,$ with $k\geq 3,$ then $C_n$ as well as its complement admit pretty good vertex state
% transfer with respect to the same sequence in $2\pi\mathbb{Z}$ \cite[Theorem 7]{pal4}. Since $C_n$ can be seen as a double cover of $C_{\frac{n}{2}},$ by Corollary \ref{5cor3}(1), the cycle $C_{2^k}$ admits plus PGST between $\frac{1}{\sqrt{2}}\ob{\e_0+\e_{\frac{n}{2}}}$ and $\frac{1}{\sqrt{2}}\ob{\e_{\frac{n}{4}}+\e_{\frac{3n}{4}}},$ whenever $k\geq 4$ with respect to the same sequence. Since, in this case, $|V(C_{2^k})|t_k\in2\pi\mathbb{Z},$ using Theorem \ref{5t15}, $\overline{C_{2^k}}$
% admits plus PGST between $\frac{1}{\sqrt{2}}\ob{\e_0+\e_{\frac{n}{2}}}$ and $\frac{1}{\sqrt{2}}\ob{\e_{\frac{n}{4}}+\e_{\frac{3n}{4}}},$ whenever $k\geq 4.$ For $k=3,$ the cycle $C_8$ has plus PST between $\frac{1}{\sqrt{2}}\ob{\e_0+\e_4}$ and $\frac{1}{\sqrt{2}}\ob{\e_2+\e_6}$ at $\frac{\pi}{2}.$ Hence,  by Theorem \ref{5t15}, the complement $\overline{C_8}$ also has plus PST between those states. Next, we check the existence of plus PGST between $\frac{1}{\sqrt{2}}\ob{\e_0+\e_{\frac{n}{2}}}$ and $\frac{1}{\sqrt{2}}\ob{\e_{\frac{n}{4}}+\e_{\frac{3n}{4}}},$ whenever $n$ is divisible by $4$ and has an odd prime factor $p.$  

% \begin{lem}
%     If $n$ is divisible by $4,$ and has an odd prime factor $p,$ then the complement of the cycle graph $\overline{C_n}$ does not admit pretty good plus state transfer between $\frac{1}{\sqrt{2}}\ob{\e_0+\e_{\frac{n}{2}}}$ and $\frac{1}{\sqrt{2}}\ob{\e_{\frac{n}{4}}+\e_{\frac{3n}{4}}}.$
% \end{lem}
% \begin{proof}
%     Let $A(\overline{C_n})=\displaystyle\sum_{r=0}^{n-1}\mu_r E_{\mu_r}$ be the spectral decomposition of adjacency matrix of $\overline{C_n},$ where 
%     % $\mu_0=n-3,$ and $\mu_r=-1-\lambda_r,$ where $\lambda_r$ are eigenvalues of $C_n$ for $1\leq r\leq n-1.$  
%     $E_{\mu_r}=\frac{1}{n}\x_r\x_r^*,$ and $\x_r(j)=\omega_n^{rj}$ as given in \eqref{5eq7}. Let $P$ be the permutation matrix corresponding to an automorphism of $\overline{C_n}$ defined by $P\e_j=\e_{j+\frac{n}{4}}.$
%    Therefore, we have
%    \[P\x_r(j)=\x_r\ob{j+\frac{n}{4}}=\omega_n^{r\ob{j+\frac{n}{4}}}=(i)^r\x_r(j).\]
%    This shows that $P\x_r=(i)^r \x_r,$ and consequently, $PE_{\mu_r}=(i)^rE_{\mu_r}.$
% Since the adjacency matrix commute with $P$ and $E_{\mu_r}$ is a polynomial in $A,$ it follows that $PE_{\mu_r}=E_{\mu_r}P$ and
%      \begin{equation}\label{5eq29}
%          E_{\mu_r}\ob{e_0+\e_{\frac{n}{2}}}=PE_{\mu_r}\ob{e_{\frac{n}{4}}+\e_{\frac{3n}{4}}}=(i)^rE_{\mu_r}\ob{e_{\frac{n}{4}}+\e_{\frac{3n}{4}}}.
%           \end{equation}
%    Let $n=mp,$ for some positive integer $m.$ For an odd prime $p,$ using a similar argument as given in the proof of Lemma \ref{5l11} and using the fact that $\lambda_r=-1-\mu_r,$ we obtain
%   \begin{equation}\label{5eq26}
%     (\mu_3-\mu_1)+\sum_{r=1}^\frac{p-1}{2}(\mu_{mr+3}-\mu_{mr+1})+ \sum_{r=1}^\frac{p-1}{2}(\mu_{mr-3}-\mu_{mr-1})=0.
%  \end{equation}
%  Suppose $\overline{C_n}$ admits plus PGST between $\frac{1}{\sqrt{2}}\ob{\e_0+\e_\frac{n}{2}}$ and $\frac{1}{\sqrt{2}}\ob{\e_{\frac{n}{4}}+\e_{\frac{3n}{4}}},$ then there exists a sequence $t_k\in\mathbb{R}$ and a complex number $\gamma$ of unit modulus such that 
% \begin{equation}\label{5eq31}
%        \lim_{k \to \infty}\sum_{r=0}^{n-1}\exp{(it_k\mu_r)}E_{\mu_r}\ob{\e_0+\e_{\frac{n}{2}}}=\gamma\ob{\e_{\frac{n}{4}}+\e_{\frac{3n}{4}}}.
%    \end{equation}
%  From \eqref{5eq29} and \eqref{5eq31}, we obtain 
%  \begin{equation}\label{5e27}
%        \lim_{k \to \infty}\exp{(it_k\mu_r)}=(-i)^r\gamma. 
%   \end{equation} 
%  Therefore, 
% $ \displaystyle\lim_{k \to \infty}\exp{(it_k(\mu_{r+2}-\mu_r))}=-1. $ Denoting the term on the left-hand side of \eqref{5eq26} as $L,$ we obtain 
%   \[\lim_{k \to \infty}\exp{\ob{it_kL}}=-1.\]
%   This contradicts \eqref{5eq26}, implying no plus PGST between $\frac{1}{\sqrt{2}}\ob{\e_0+\e_{\frac{n}{2}}}$ and $\frac{1}{\sqrt{2}}\ob{\e_{\frac{n}{4}}+\e_{\frac{3n}{4}}}.$ 
% \end{proof}

%  Note that Lemma \ref{5l7} also holds for the complement $\overline{C_n}.$ Since $\overline{C_n}$ has pair PGST if and only if $C_n$ has pair PGST, it may be possible to have plus PGST in $\overline{C_n},$ whenever $n=2^kp,$ for a positive integer $k$ and an odd prime $p.$ However, this is rules out by Lemma \ref{5l9} together with the following observation.

%  \begin{lem}
%      The complement of a cycle graph $\overline{C_n}$ does not admit pretty good plus state transfer between $\frac{1}{\sqrt{2}}\ob{\e_0+\e_s}$ and $\frac{1}{\sqrt{2}}\ob{\e_{\frac{n}{2}}+\e_{\frac{n}{2}+s}},$ for $0<s<\frac{n}{2},$ whenever $n=2^k p,$ where $k$ is a positive integer and $p$ is an odd prime.  
%  \end{lem}
%  \begin{proof}
%      Let $P$ be the permutation matrix corresponding to an automorphism of $\overline{C_n}$ defined by $P\e_j=\e_{j+\frac{n}{2}}.$ Using a similar approach as in the proof of Lemma \ref{5l11}, we obtain 
% \begin{equation}\label{5eq33}
%         E_{\mu_r}(e_0+\e_s)=(-1)^r E_{\mu_r}\ob{e_{\frac{n}{2}}+\e_{\frac{n}{2}+s}}.
%          \end{equation}
%          For an odd prime $p,$ using a similar argument as given in the proof of Lemma \ref{5l11} and using the fact that $\lambda_r=-1-\mu_r,$ we obtain
%   \begin{equation}\label{5eq30}
%     (\mu_2-\mu_1)+\sum_{r=1}^\frac{p-1}{2}(\mu_{mr+2}-\mu_{mr+1})+ \sum_{r=1}^\frac{p-1}{2}(\mu_{mr-2}-\mu_{mr-1})=0.
%  \end{equation}
%  A similar argument as presented in the proof of Lemma \ref{5l11}, contradicts the existence of plus PGST between $\frac{1}{\sqrt{2}}\ob{\e_0+\e_s}$ and $\frac{1}{\sqrt{2}}\ob{\e_{\frac{n}{2}}+\e_{\frac{n}{2}+s}}$ in $\overline{C_n}.$ 
%  \end{proof}

%  Now, we give a complete characterization of plus PGST in the cycle graph $C_n$ and its complement $\overline{C_n}.$
%  \begin{thm}\label{5th10}
%       A cycle $C_n$ as well as its complement $\overline{C_n}$ admit pretty good plus state transfer if and only if $n=2^k$ with $k\geq 2.$ Moreover, for $k\geq 3,$ every plus state in $C_{2^k}$ and $\overline{C_{2^k}}$ exhibit pretty good plus state transfer.
%  \end{thm}
 \section{PGST in weighted paths with potential}\label{5sec6}
 Let $P_n$ be a path on $n$ vertices with vertex set $\{1,2,\ldots,n\},$ where two vertices $j$ and $k$ are adjacent if and only if $|j-k|=1.$ 
  We study the existence of vertex PGST in certain weighted paths with potential using equitable partitions. A partition $\Pi$ of the vertex set of a graph $G$ with disjoint cells $V_1, V_2, \ldots, V_d$ is said to be \textit{equitable} if each vertex in $V_j$ is adjacent to exactly $c_{jk}$ vertices in $V_k,$ where 
% \[\sum_{b \in V_k} w(a,b)=c_{jk}~~~\text{for each}~ a\in V_j,\]
$c_{jk}$ is a constant depending only on the cells $V_j$ and $V_k$. 
% The \textit{characteristic matrix} $\mathcal{C}$ associated to $\Pi$ is an $n \times d$ matrix where the columns represent the normalized characteristic vectors of $V_1, V_2,\dots, V_d.$ 
The graph with the $d$ cells of $\Pi$ as its vertices having adjacency matrix $A(G/\Pi),$ where $A(G/\Pi)_{j,k}=\sqrt{c_{jk}c_{kj}},$ is called the \textit{symmetrized quotient graph} of $G$ over $\Pi,$ and is denoted by $G/\Pi.$ A relation between vertex PGST in $G/\Pi$ and plus PGST in $G$ is observed in \cite[Remark 5]{god8}, which implies that for $V_1=\{a,b\}$ and $V_2=\{c,d\},$ PGST occurs between $V_1$ and $V_2$ in $G/\Pi$ if and only if PGST occurs between $\frac{1}{\sqrt{2}}\ob{\e_a+\e_b}$ and $\frac{1}{\sqrt{2}}\ob{\e_c+\e_d}$ in $G.$

It is well known that PST occurs only on the paths $P_2$ and $P_3$ with respect to the adjacency matrix \cite{god1}. This raises a question of whether PST could be induced between the endpoints of
a path of arbitrary length by placing a suitable potential on the endpoints. Kempton et al. \cite{kem1} showed that there is no potential that induces PST between the end
points of a path of length at least $4.$ However, for a path $P_n$ of any length, there is some choice of potential such that by
placing the potential on each endpoint of $P_n,$ there is PGST between the endpoints \cite{kem2}. 
 % Consider a path on more than three vertices with potential $w$ on the end vertices only.
  Further, it is observed in \cite{kirk5} that for paths on more than two vertices, the set of weights $w\geq 1$ assigned to the end vertices, for which the resulting path does not admit PGST, forms a dense subset of $[1,\infty).$ 
 % such that if $w$ is chosen from that subset,
 % then PGST between the end vertices is impossible.
Here, we present a few observations on the existence of PGST on weighted paths with potentials.

\begin{figure}[]
\centering
% ---------------- LEFT SIDE (First Figure) ----------------
\begin{minipage}{0.48\textwidth}
\centering
\begin{tikzpicture}[scale=0.6]
\tikzstyle{every node}=[circle, thick, black!90, fill=white, scale=0.65]
\node[draw,minimum size=0.55cm, inner sep=0 pt] (0) at (0,0) {$0$};
\node[draw,minimum size=0.55cm, inner sep=0 pt] (1) at (2,-2) {$1$};
\node[draw,minimum size=0.55cm, inner sep=0 pt] (2) at (4,-2) {$2$};
\fill (5.1,-2) circle (2.5pt);
\fill (5.5,-2) circle (2.5pt);
\fill (5.9,-2) circle (2.5pt);
\node[draw,minimum size=0.55cm, inner sep=0 pt] (3) at (7,-2) {};
\node[draw,minimum size=0.55cm, inner sep=1.5 pt] (4) at (9,-2) {$n-1$};
\node[draw,minimum size=0.55cm, inner sep=0 pt] (5) at (9,2) {$n$};
\node[draw,minimum size=0.55cm, inner sep=0 pt] (6) at (7,2) {};
\node[draw,minimum size=0.55cm, inner sep=0 pt] (7) at (4,2) {};
\fill (5.1,2) circle (2.5pt);
\fill (5.5,2) circle (2.5pt);
\fill (5.9,2) circle (2.5pt);
\node[draw,minimum size=0.55cm, inner sep=1.5 pt] (8) at (2,2) {$\scriptsize 2n-2$};
\draw (0)--(1)--(2);
\draw (3)--(4)--(5)--(6);
\draw (7)--(8)--(0);
\end{tikzpicture}
\subcaption{}
\end{minipage}
\hfill
% ---------------- RIGHT SIDE (Second & Third) ----------------
\begin{minipage}{0.45\textwidth}
\centering
% -------- Second Figure (Top) --------
\begin{tikzpicture}[scale=1.1]
\tikzstyle{every node}=[circle, thick, black!90, fill=white, scale=0.65]
\node[draw,minimum size=0.55cm, inner sep=0 pt] (1) at (0,0) {$1$};
\node[draw,minimum size=0.55cm, inner sep=0 pt] (2) at (1,0) {$2$};
\node[draw,minimum size=0.55cm, inner sep=0 pt] (3) at (2,0) {$3$};
\fill (2.6,0) circle (1.5pt);
\fill (3,0) circle (1.5pt);
\fill (3.4,0) circle (1.5pt);
\node[draw,minimum size=0.55cm, inner sep=0 pt] (4) at (4,0) {};
\node[draw,minimum size=0.55cm, inner sep=0 pt] (5) at (5,0) {$n$};
\draw[thick]
(5) .. controls +(60:1.5cm) and +(120:1.5cm) ..
node[pos=0.3, above right] {$1$}
(5);
\draw (1) -- (2) node[midway, above, yshift=0.5pt] {$\sqrt{2}$};
\draw (2) -- (3);
\draw (4) -- (5);
\end{tikzpicture}
\subcaption{}
\vspace{1cm} % space between top and bottom figure
% -------- Third Figure (Bottom) --------
\begin{tikzpicture}[scale=1]
\tikzstyle{every node}=[circle, thick, black!90, fill=white, scale=0.65]
\node[draw,minimum size=0.55cm, inner sep=0 pt] (1) at (-1,0.6) {$a$};
\node[draw,minimum size=0.55cm, inner sep=0 pt] (2) at (-1,-0.6) {$b$};
\node[draw,minimum size=0.55cm, inner sep=0 pt] (3) at (0,0) {$1$};
\node[draw,minimum size=0.55cm, inner sep=0 pt] (4) at (1,0) {$2$};
\fill (1.6,0) circle (1.5pt);
\fill (2,0) circle (1.5pt);
\fill (2.4,0) circle (1.5pt);
\node[draw,minimum size=0.55cm, inner sep=0 pt] (5) at (3,0) {};
\node[draw,minimum size=0.55cm, inner sep=1.5 pt] (6) at (4,0) {$n-1$};
\draw[thick]
(6) .. controls +(60:1.5cm) and +(120:1.5cm) ..
node[pos=0.3, above right] {$1$}
(6);
\draw (1)--(3)--(2);
\draw (3)--(4);
\draw (5)--(6);
\end{tikzpicture}
\subcaption{}
\end{minipage}
\caption{(a) Odd cycle $C_{2n-1},$ (b) Symmetrized quotient graph of $C_{2n-1},$ (c) A graph with symmetrized quotient matrix equal to the adjacency matrix of Figure (b).}
\label{5fig4}
\end{figure}
%  The odd cycle $C_{2n-1}$ given in Figure \ref{5fig4}(a) has an equitable partition $\Pi$ with cells $V_1=\{0\}, V_2=\{1, 2n-2\},\ldots, V_{n}=\{n-1,n\}.$
%  Figure \ref{5fig4}(b) represents the symmetrized quotient graph corresponding to $\Pi.$ Lemma \ref{5l8} together with \cite[Remark 5]{god8} demonstrates that there is no PGST between any pair of vertices in the set $\{2,3,\ldots,n\}.$ in the symmetrized quotient graph of $C_{2n-1}$ given in Figure \ref{5fig4}(b). Suppose the graph has PGST between $1$ and $j\neq 1.$ Then by \cite[Corollary 5]{god8}, the cycle in Figure \ref{5fig4}(a)
%  has PGST between $\e_0$ and $\frac{1}{\sqrt{2}}\ob{\e_a+\e_b},$ where $a$ and $b$ are in $V_j$ for some $j\neq 1.$ Then there is a sequence of time $t_k\in\mathbb{R}$ and a complex number $\gamma$ of unit modulus such that  
%     \begin{equation}\label{5e29}
%         \lim_{k \to \infty}U_G\ob{t_k}\e_0=\frac{\gamma}{\sqrt{2}}\ob{\e_a+\e_b}.
%     \end{equation}
%  Let $P$ be a permutation matrix corresponding to an automorphism of $C_n$ which fixes only the vertex $a.$ Now, multiplying $P$ on both sides of \eqref{5e29} and subtracting the resulting equation from it, gives
%  \[
%         \lim_{k \to \infty}U_G\ob{t_k}(\e_0-P\e_0)=\frac{\gamma}{\sqrt{2}}\ob{\e_b-P\e_b},
%     \]
%  which is a contradiction as $\left\|\e_0-P\e_0\right\|\neq \left\|\frac{1}{\sqrt{2}}\ob{\e_b-P\e_b}\right\|.$ 
% Therefore, we have the following result.
\begin{thm}
    There is no pretty good vertex state transfer in a weighted path on $n$ vertices, with potential $1$ only at the end vertex $n$ and $w(1,2)=\sqrt{2}$ and all other edges having weight $1.$ 
\end{thm}
\begin{proof}
    The odd cycle $C_{2n-1}$ given in Figure \ref{5fig4}(a) has an equitable partition $\Pi$ with cells $V_1=\{0\}, V_2=\{1, 2n-2\},\ldots, V_{n}=\{n-1,n\}.$
  Lemma \ref{5l8} together with \cite[Remark 5]{god8} demonstrates that there is no PGST between any pair of vertices in  $\{2,3,\ldots,n\}$ in the symmetrized quotient graph of $C_{2n-1}$ given in Figure \ref{5fig4}(b). Suppose the graph has PGST between $1$ and $j,$ for some $j\in\{2,3,\ldots,n\}.$ Then by \cite[Remark 5]{god8}, the cycle in Figure \ref{5fig4}(a)
 has PGST between $\e_0$ and $\frac{1}{\sqrt{2}}\ob{\e_a+\e_b},$ where $a$ and $b$ are in $V_j.$ This gives a contradiction to Theorem \ref{5th1} as there is an automorphism of $C_n$ which fixes only the vertex $a.$ Hence the result follows.  
\end{proof} 
 
 
%  Then there is a sequence of time $t_k\in\mathbb{R}$ and a complex number $\gamma$ of unit modulus such that  
%     \begin{equation}\label{5e29}
%         \lim_{k \to \infty}U_G\ob{t_k}\e_0=\frac{\gamma}{\sqrt{2}}\ob{\e_a+\e_b}.
%     \end{equation}
%  Let $P$ be a permutation matrix corresponding to an automorphism of $C_n$ which fixes only the vertex $a.$ Now, multiplying $P$ on both sides of \eqref{5e29} and subtracting the resulting equation from it, gives
%  \[
%         \lim_{k \to \infty}U_G\ob{t_k}(\e_0-P\e_0)=\frac{\gamma}{\sqrt{2}}\ob{\e_b-P\e_b},
%     \]
%  which is a contradiction as $\left\|\e_0-P\e_0\right\|\neq \left\|\frac{1}{\sqrt{2}}\ob{\e_b-P\e_b}\right\|.$ 
% Therefore, we have the desired result.

Note that \cite[Corollary 9.2]{god1} also holds for PGST. Therefore, if $G$ admits PGST between $a$ and $b,$ then each automorphism fixing $a$ must fix $b.$ Suppose there is PGST between $a$ and $j,$ for some $j\in\{1,2,\ldots, n-1\},$ in Figure \ref{5fig4}(c). Since there is an automorphism of the graph in Figure \ref{5fig4}(c) that fixes $j$ but not $a,$ there is no PGST between $a$ and $j.$
 The graph in Figure \ref{5fig4}(c) has an equitable partition $\Pi'$ with cells $V_1=\{a,b\}, V_j=\{j-1\},$ for $2\leq j\leq n.$ The symmetrized quotient matrix relative to $\Pi'$ is equal to the adjacency matrix of the graph given in Figure \ref{5fig4}(b). Therefore, using \cite[Remark 5]{god8}, the result follows.
\begin{cor}
Let $G$ be the graph obtained from $P_n$ by joining two pendant vertices to vertex 
$1$ and assigning a potential 
$1$ only at vertex $n.$ Then there is no pretty good vertex state transfer in
$G$ from the vertices of $P_n.$ 
\end{cor}

\begin{figure}[]
\centering

% ---------------- LEFT SIDE (First Figure) ----------------
\begin{minipage}{0.48\textwidth}
\centering
\begin{tikzpicture}[scale=0.6]
\tikzstyle{every node}=[circle, thick, black!90, fill=white, scale=0.65]

\node[draw,minimum size=0.55cm, inner sep=0 pt] (0) at (0,0) {$0$};
\node[draw,minimum size=0.55cm, inner sep=0 pt] (1) at (2,-2) {$1$};
\node[draw,minimum size=0.55cm, inner sep=0 pt] (2) at (4,-2) {$2$};
\fill (5.1,-2) circle (2.5pt);
\fill (5.5,-2) circle (2.5pt);
\fill (5.9,-2) circle (2.5pt);
\node[draw,minimum size=0.55cm, inner sep=0 pt] (3) at (7,-2) {};
\node[draw,minimum size=0.55cm, inner sep=1.5 pt] (4) at (9,-2) {$n-2$};
\node[draw,minimum size=0.55cm, inner sep=1.5 pt] (5) at (11,0) {$\scriptsize n-1$};
\node[draw,minimum size=0.55cm, inner sep=0 pt] (6) at (9,2) {$\scriptsize n$};
\node[draw,minimum size=0.55cm, inner sep=0 pt] (7) at (7,2) {};
\node[draw,minimum size=0.55cm, inner sep=0 pt] (8) at (4,2) {};
\fill (5.1,2) circle (2.5pt);
\fill (5.5,2) circle (2.5pt);
\fill (5.9,2) circle (2.5pt);
\node[draw,minimum size=0.55cm, inner sep=1.5 pt] (9) at (2,2) {$\scriptsize 2n-3$};

\draw (0)--(1)--(2);
\draw (3)--(4)--(5)--(6)--(7);
\draw (8)--(9)--(0);
\end{tikzpicture}
\subcaption{}
\end{minipage}
\hfill
% ---------------- RIGHT SIDE (Second & Third) ----------------
\begin{minipage}{0.45\textwidth}
\centering

% -------- Second Figure (Top) --------
\begin{tikzpicture}[scale=1.2]
\tikzstyle{every node}=[circle, thick, black!90, fill=white, scale=0.65]

\node[draw,minimum size=0.55cm, inner sep=0 pt] (1) at (0,0) {$1$};
\node[draw,minimum size=0.55cm, inner sep=0 pt] (2) at (1,0) {$2$};
\node[draw,minimum size=0.55cm, inner sep=0 pt] (3) at (2,0) {$3$};
\fill (2.6,0) circle (1.2pt);
\fill (3,0) circle (1.2pt);
\fill (3.4,0) circle (1.2pt);
\node[draw,minimum size=0.55cm, inner sep=0 pt] (4) at (4,0) {};
\node[draw,minimum size=0.55cm, inner sep=0 pt] (5) at (5,0) {$n$};
\draw (1) -- (2) node[midway, above, yshift=0.5pt] {$\sqrt{2}$};
\draw (4) -- (5) node[midway, above, yshift=0.5pt] {$\sqrt{2}$};
\draw (2) -- (3);
\draw (4) -- (5);

\end{tikzpicture}
\subcaption{}
\vspace{1cm} % space between top and bottom figure

% -------- Third Figure (Bottom) --------
\begin{tikzpicture}[scale=1]
\tikzstyle{every node}=[circle, thick, black!90, fill=white, scale=0.65]

\node[draw,minimum size=0.55cm, inner sep=0 pt] (1) at (-1,0.6) {$a$};
\node[draw,minimum size=0.55cm, inner sep=0 pt] (2) at (-1,-0.6) {$b$};
\node[draw,minimum size=0.55cm, inner sep=0 pt] (3) at (0,0) {$1$};
\node[draw,minimum size=0.55cm, inner sep=0 pt] (4) at (1,0) {$2$};
\fill (1.6,0) circle (1.5pt);
\fill (2,0) circle (1.5pt);
\fill (2.4,0) circle (1.5pt);
\node[draw,minimum size=0.55cm, inner sep=0 pt] (5) at (3,0) {};
\node[draw,minimum size=0.55cm, inner sep=1.5 pt] (6) at (4,0) {$\scriptsize n-2$};
\node[draw,minimum size=0.55cm, inner sep=0 pt] (7) at (5,0.6) {$c$};
\node[draw,minimum size=0.55cm, inner sep=0 pt] (8) at (5,-0.6) {$d$};
% \draw[thick]
% (6) .. controls +(60:1cm) and +(120:1cm) ..
% node[pos=0.3, above right] {$1$}
% (6);
\draw (1)--(3)--(2);
\draw (3)--(4);
\draw (5)--(6)--(7);
\draw (6)--(8);

\end{tikzpicture}
\subcaption{}
\end{minipage}

\caption{(a) Even cycle $C_{2n-2},$ (b) Symmetrized quotient graph of $C_{2n-2},$ (c) A graph
with symmetrized quotient matrix equal to the adjacency matrix of Figure (b).}
\label{5fig5}

\end{figure}

% Now, consider an even cycle $C_{2n-2}$ as given in Figure \ref{5fig5}(a). The cells $V_1=\{0\},V_2=\{1, 2n-3\},\ldots,V_{n+1}=\{n-2,n\},V_{n+2}=\{n-1\}$ form an equitable partion of $V(C_{2n+2})$ giving rise to the symmetrized quotient graph as given in Figure \ref{5fig5}(b). 
% By the same argument as in the case of the odd cycle $C_{2n-1},$ PGST does not occur between $\e_0$ and $\frac{1}{\sqrt{2}}\ob{\e_a+\e_b},$ whenever $a,b\in V_j$ with $j\notin\{1,n+2\}.$ The result now follows from \cite[Remark 5]{god8}
%  and Theorem \ref{5th10}.
Now we consider the symmetrized quotient graph of an even cycle to have the following observations.
 \begin{thm}
    Let $P_n$ be a weighted path  with $w(1,2)=w(n-1,n)=\sqrt{2}$ and the remaining edges have weight $1.$ Then $P_n$ admits pretty good vertex state transfer if and only if $n-1=2^k$ with $k\geq 1.$
\end{thm}
\begin{proof}
    The even cycle $C_{2n-2}$ given in Figure \ref{5fig5}(a) has an equitable partition with the cells $V_1=\{0\},V_2=\{1, 2n-3\},\ldots,V_{n-1}=\{n-2,n\},V_{n}=\{n-1\}.$ Using Theorem \ref{5th10} and \cite[Remark 5]{god8}, we have the weighted path exhibits PGST if and only if $2n-2=2^k$ with $k\geq 2.$ Hence the result follows.
  % Let $a,b\in V_j$ with $j\notin\{1,n\}.$ By, Theorem \ref{5th1}, there is no PGST between  between $\e_0$ and $\frac{1}{\sqrt{2}}\ob{\e_a+\e_b}.$     
 %    By the same argument as in the case of the odd cycle $C_{2n-1},$ PGST does not occur between $\e_0$ and $\frac{1}{\sqrt{2}}\ob{\e_a+\e_b},$ whenever The result now follows from \cite[Remark 5]{god8}
 % and Theorem \ref{5th10}.
\end{proof}
  The graph in Figure \ref{5fig5}(c) has an equitable partition $\Pi''$ with cells $V_1=\{a,b\}, V_j=\{j-1\}, V_n=\{c,d\}$ for $2\leq j \leq n-1.$ The symmetrized quotient matrix relative to $\Pi''$ is equal to the adjacency matrix of the graph given in Figure \ref{5fig5}(b). Therefore, by \cite[Remark 5]{god8}, the result follows.
 \begin{cor}\label{5t31}
 Let $G$ be the graph obtained from $P_n$ by attaching two pendant vertices to each end vertex of $P_n.$ Then pretty good vertex state transfer occurs in $G$ from the vertices of 
 $P_n$ if and only if $n+1=2^k$ with $k\geq 1.$ 
 \end{cor}
For $n=2,$ the graph $G$ described in Corollary \ref{5t31} becomes a double star $S_{2,2},$ and the result is consistent with the result \cite[Theorem 5.3(b)]{fan}. Furthermore, when $n=3,$ the graph $G$ is the 1-sum of star graph $F_{2,2},$ and the conclusion agrees with \cite[Lemma 3.3]{hou}.

Note that a path $P_n$ with potential $1$ only at the end vertices can
 be realized as a symmetrized quotient graph of the even cycle. Hence,  using Theorem \ref{5th10} and \cite[Remark 5]{god8}, we have the following observation.

 \begin{cor}
     Let $P_n$ be a path on $n$ vertices with potential $1$ only at the end vertices. Then $P_n$ admits pretty good vertex state transfer if and only if $n=2^k,$ with $k\geq 1.$
 \end{cor}
% \section{Fractional revival in complementary prisms}\label{5sec4}
% %  \begin{thm}
% %     Let $M$ be a real symmetric matrix associated with a graph $G.$ Let a state $\psi$ have only two eigenvalues $\lambda_1,\lambda_2$ in its support. Then $G$ has fractional revival from $\psi$ to $\phi=(E_{\lambda_1}-E_{\lambda_2})\psi.$ 
% % \end{thm}
% % The complement of a graph $G,$ denoted by $\overline{G},$ is the graph with the same vertex set as $G,$ where two distinct vertices are adjacent in $\overline {G}$ if and only if they are not adjacent in $G.$
% The complementary prism $G\overline{G}$ of a graph $G$ is obtained from the disjoint union of $G$ and its complement $\overline{G}$ by adding an edge for each pair of vertices $(a,a'),$ where $a$ is in $G$ and its copy $a'$ is in $\overline{G}.$ 
% % Let $M(G)$ denote the adjacency, Laplacian, or signless Laplacian matrix of $G.$ 
% Then for $M\in\{A,L,Q\},$ the matrix representing the complementary prism $G\overline{G}$ of $G$ is defined as follows.
% \[M(G\overline{G})=\begin{bmatrix}
%     M(G)+\delta I & \tau I\\
%     \tau I & M(\overline{G})+\delta I
% \end{bmatrix},\]
% where $\delta=\left\{ \begin{array}{rcl}
%  0, & \mbox{if M = A} \\ 
%  1, & \text{otherwise} 
%  \end{array}\right.$,\quad and \quad $\tau=\left\{ \begin{array}{rcl}
%  -1, & \mbox{if M = L} \\ 
%  1, & \text{otherwise} 
%  \end{array}\right..$ 
%  % and $I$ is the identity matrix of appropriate order. 
% \\The relationships between the eigenvalues and eigenvectors of the adjacency, Laplacian, and signless Laplacian matrices of the complementary prism $G\overline{G}$ and those of the graph $G$ are investigated in \cite{card1}. An eigenvalue $\lambda$ of $M(G)$ is said to be a main eigenvalue if there is an associated eigenvector $v$ which is not orthogonal to $\mathbf{1}.$ Otherwise, we say that $\lambda$ is non-main. For further definitions and related concepts, the reader is referred to \cite{bro,cve1}.
% % Let $M(G)$ have $d$ distinct eigenvalues. We denote the spectrum of $M(G)$ by 
% % \[\text{spec}(M(G))=\{\lambda_1^{[m_1]},\lambda_2^{[m_2]},\ldots,\lambda_d^{[m_d]}\},\] 
% % where $[m_j]$ indicates the multiplicity of the eigenvalue $\lambda_j.$ 
% Here, we investigate the existence of FR in the complementary prism graph $G\overline{G}$ of a graph $G$ with respect to the  adjacency, Laplacian and signless Laplacian matrices.
% % \subsection{FR with respect to adjacency matrix}
% % Let $G$ be a connected and regular graph. 
% The following result determines the full set of eigenvalues and corresponding eigenvectors of the complementary prism $G\overline{G}$ of a connected and regular graph $G$ with respect to the adjacency matrix. In what follows, $\alpha_{1,2}(\lambda_j)$ denotes the pair $\alpha_1(\lambda_j)$ and $\alpha_2(\lambda_j),$ while $\beta_{1,2}(n,k)$ denotes $\beta_1(n,k)$ and $\beta_2(n,k).$
% \begin{cor}\cite{card1}(Adjacency matrix)\label{5cor1}
%     If $G$ is a connected $k$-regular graph on $n$ vertices with spectrum $\cb{k,\lambda_2^{\ob{m_2}}, \ldots, \lambda_d^{\ob{m_d}}},$ and $\x_j\perp \mathbf{1}$ is an eigenvector corresponding to $\lambda_j$ then the eigenvalues and eigenvectors of complementary prism $G\overline{G}$ are: 
%     \begin{enumerate}
%         \item $\alpha_{1,2}(\lambda_j)=\dfrac{-1\pm\sqrt{(2\lambda_j+1)^2+4}}{2}$ with multiplicity $m_j,$ for $2\leq j\leq d$ and eigenvectors $\begin{bmatrix}
%             \x_j\\
%             \ob{\alpha_{1,2}(\lambda_j)-\lambda_j}\x_j
%         \end{bmatrix}$ which are orthogonal to the all one vector.
%  \item$\beta_{1,2}(n,k)=\dfrac{n-1\pm\sqrt{(n-1-2k)^2+4}}{2}$ with eigenvectors $\begin{bmatrix}
%             \mathbf{1}\\
%             \ob{\beta_{1,2}(n,k)-k}\mathbf{1}
%         \end{bmatrix}.$ $\beta_{1,2}(n,k)$ are the two main eigenvalues when $G\overline{G}$ is non-regular.
%     \end{enumerate}
% \end{cor}
% Now, we have the following observation.
% \begin{lem}\label{5t1}
%     Let $\u$ be a fixed state in a regular graph $G$ such that $\mathbf{1}^T\u=0.$ Then for all $t\in\mathbb{R},$ 
%     \[U_{A(G\overline{G})}(t)\begin{bmatrix}
%         \u\\\mathbf{0}
% \end{bmatrix}=\alpha(t)\begin{bmatrix}
%         \u\\\mathbf{0}
%     \end{bmatrix}+\beta(t)\begin{bmatrix}
%        \mathbf{0}\\\u
%     \end{bmatrix},~~~\text{for some $\alpha(t)$ and $\beta(t).$}
% \]
% \end{lem}

% \begin{proof}
%     Since $\u$ is a fixed state in $G,$ the eigenvalue support of $\u$ contains only one eigenvalue $\lambda,$ say. By Corollary \ref{5cor1}, the eigenvalue support of $\ob{\u^T,\mathbf{0}}^T$ in $G\overline{G}$ contains the eigenvalue $\alpha_{1,2}(\lambda).$ Let $w_1=\begin{bmatrix}
%             \u\\
%             \gamma_1\u
%         \end{bmatrix}$ and $w_2=\begin{bmatrix}
%             \u\\
%             \gamma_2\u
%         \end{bmatrix}$ be eigenvectors corresponding to $\alpha_1(\lambda)$ and $\alpha_2(\lambda),$ respectively,
%         with $\gamma_1=\alpha_1(\lambda)-\lambda$ and $\gamma_2=\alpha_2(\lambda)-\lambda.$ Note that, since $\gamma_1\neq\gamma_2,$ the state $\begin{bmatrix}
%             \u\\
%            \mathbf{0}
%         \end{bmatrix}\in span\cb{w_1,w_2}.$  For some scalars $c_1$ and $c_2$ let $\begin{bmatrix}
%             \u\\
%             \mathbf{0}
%         \end{bmatrix}=c_1w_1+c_2w_2.$ Now,
% \begin{equation*}
% \begin{split}
% \exp{(itA(G\overline{G}))}\begin{bmatrix}
%             \u\\
%             \mathbf{0}
%         \end{bmatrix}&=\exp{\ob{it\alpha_1}}E_{\alpha_1}c_1w_1
% +\exp{\ob{it\alpha_2}}E_{\alpha_2}c_2w_2\\
% &=\exp{\ob{it\alpha_1}}c_1w_1+\exp{\ob{it\alpha_2}}c_2w_2\\
% &=\ob{c_1\exp{\ob{it\alpha_1}}+c_2\exp{\ob{it\alpha_2}}}\begin{bmatrix}
%             \u\\
%              \mathbf{0} 
% \end{bmatrix}+\ob{c_1\gamma_1\exp{\ob{it\alpha_1}}+c_2\gamma_2\exp{\ob{it\alpha_2}}}\begin{bmatrix}
%             \mathbf{0}\\
%               \u
% \end{bmatrix}
% \end{split}
% \end{equation*}
% Here, the second equality follows since $\begin{bmatrix}
%             \u\\
%             \gamma_j\u
%         \end{bmatrix}$ is an eigenvector corresponding to $\alpha_j,$ for $j=1,2.$ Considering $\alpha(t)=c_1\exp{\ob{it\alpha_1}+c_2\exp{\ob{it\alpha_2}}}$ and $\beta(t)=c_1\gamma_1\exp{\ob{it\alpha_1}}+c_2\gamma_2\exp{\ob{it\alpha_2}},$ the result follows.
% \end{proof}

% From Lemma \ref{5t1}, the following result is now immediate.
% \begin{thm}\label{5t2}
%      Let $\u$ be a fixed state in a regular graph $G$ such that $\mathbf{1}^T\u=0.$ Then the complementary prism graph $G\overline{G}$  exhibits fractional revival from the state $\ob{\u^T,\mathbf{0}}^T$ to $\ob{\mathbf{0},\u^T}^T$ for all time $t$ with respect to the adjacency matrix. 
%      \end{thm}

% % \begin{proof}
% % Since $||U_{A({G\overline{G}})}(t)\begin{bmatrix}\psi\\0\end{bmatrix}||^2=1$ and the states $\begin{bmatrix}
% %     \psi\\0
% % \end{bmatrix}$ and  $\begin{bmatrix}
% %     0\\\psi
% % \end{bmatrix}$ are orthonormal unit vectors, the proof of Theorem \ref{5t1} gives $|\alpha(t)|^2+|\beta(t)|^2=1$ for all $t.$  
% % \end{proof}

% The following result gives the Laplacian spectrum of the complementary prism graph.
% % \subsection{FR with respect to Laplacian matrix}
% \begin{thm}\cite{card1}(Laplacian matrix)\label{5t3}
%     Let $G$ be a graph on $n$ vertices with Laplacian eigenvalues $\lambda_1\geq \lambda_2\geq\cdots\geq \lambda_{n-1}\geq\lambda_n=0.$ For each $j =1, 2, \ldots, n -1,$ if $\ob{\lambda_j,\x_j}$ is a Laplacian eigenpair of $G,$ then
% \[\alpha_{1,2}(\lambda_j) = \dfrac{n + 2 \pm
% \sqrt{(n - 2\lambda_j)^2+4}}{2}\]
% are Laplacian eigenvalues of $G\overline{G}$ with associated eigenvectors $\begin{bmatrix}
%             \x_j\\
%             \ob{\lambda_j-\alpha_{1,2}\ob{\lambda_j}}\x_j
%         \end{bmatrix}.$ The other two eigenvalues are $2$ and $0$ with  eigenvectors  $\begin{bmatrix}
%             \mathbf{1}\\
%             -\mathbf{1}
%         \end{bmatrix}$ and $\begin{bmatrix}
%             \mathbf{1}\\
%             \mathbf{1}
%         \end{bmatrix},$ respectively.   
% \end{thm}
% The next result characterizes FR and PST in complementary prism graphs with respect to the Laplacian matrix.
%     \begin{thm}\label{5t4}
%      Let $\u$ be a fixed state in $G$ such that $\mathbf{1}^T\u=0.$ Then the complementary prism graph $G\overline{G}$  exhibits fractional revival from the state $\ob{\u^T,\mathbf{0}}^T$ to $\ob{\mathbf{0},\u^T}^T$ for all time $t$ with respect to the Laplacian matrix. Moreover, the graph $G\overline{G}$ exhibits perfect state transfer between 
%      $\ob{\mathbf{1}^T,\mathbf{0}}^T$ and $\ob{\mathbf{0},\mathbf{1}^T}^T$
%      % \begin{bmatrix}
%      %     \mathbf{j}\\0
%      % \end{bmatrix}$ and $\begin{bmatrix}
%      %     0\\j
%      % \end{bmatrix}$ 
%     at odd multiple of $\frac{\pi}{2}.$
% \end{thm}
% \begin{proof}
%     The proof of the first part can be obtained using Theorem \ref{5t3}, and a similar approach as in Lemma \ref{5t1}. The eigenvalue support of the state $(\mathbf{1}^T,\mathbf{0})^T$ contains the eigenvalues $2$ and $0.$ 
%     From Theorem \ref{5t3}, one may observe that the states $(\mathbf{1}^T,\mathbf{0})^T$ and $(\mathbf{0},\mathbf{1}^T)^T$ are strongly cospectral. Hence, by \cite[Theorem 5.2]{god7}, the result follows.
% %     Therefore,
% %  \begin{equation}\label{5e3}
% % \begin{split}
% % \exp{(itL(G\overline{G}))}\begin{bmatrix}
% %             j\\
% %             0
% %         \end{bmatrix}&=\exp{(2it)}E_2\begin{bmatrix}
% %             j\\0
% %         \end{bmatrix}+E_{0}\begin{bmatrix}
% %             j\\0
% %         \end{bmatrix}\\
% % &=\frac{1}{2}\exp{(2it)}E_2\ob{\begin{bmatrix}
% %             j\\j
% % \end{bmatrix}+\begin{bmatrix}
% %             j\\-j
% % \end{bmatrix}}+\frac{1}{2}E_0\ob{\begin{bmatrix}
% %             j\\j
% % \end{bmatrix}+\begin{bmatrix}
% %             j\\-j
% % \end{bmatrix}}\\
% % &=\frac{1}{2}\exp{(2it)}\begin{bmatrix}
% %             j\\-j
% % \end{bmatrix}+\frac{1}{2}\begin{bmatrix}
% %             j\\j
% % \end{bmatrix}
% % \end{split}
% % \end{equation}   
% % If $t$ is an odd multiple of $\frac{\pi}{2},$ then the equality in $\eqref{5e3}$ gives $U_{L(G\overline{G})}(t)(j^T,0)^T=(0,j^T)^T.$
% \end{proof}
% % \textbf{Another way:}
% % The eigenvalue support of the state $(j^T,0)^T$ contains the eigenvalues $2$ and $0.$ From Theorem \ref{5t3}, one may observe that the states $(j^T,0)^T$ and $(0,j^T)^T$ are strongly cospectral. Hence, by \cite[Theorem 23(1)]{god7}, the result follows.

% \begin{exm}
%     Let $P_n$ be a path on $n$ vertices. One may observe $P_4$ as a complementary prism of empty graph on two vertices. Then Theorem \ref{5t4} yields perfect state transfer between $\frac{1}{\sqrt{2}}(\e_1+\e_4)$ and $\frac{1}{\sqrt{2}}(\e_2+\e_3)$ at $\frac{\pi}{2},$ which is consistent with the result \cite[Corollary 7.8]{god7}. 
% \end{exm}



% % \subsection{FR with respect to signless Laplacian matrix}
% % The following result gives the spectra of the complementary prism graph relative to the signless Laplacian matrix.
% We now turn to the signless Laplacian spectrum of the complementary prism graph of a connected regular graph.
% \begin{thm}\cite{card1}(Signless Laplacian matrix)\label{5t5}
%     If $G$ is a connected $k$-regular graph on $n$ vertices with spectrum $\cb{2k, \lambda_2^{(m_2)},\ldots, \lambda_d^{(m_d)}},$ and $\x_j\perp \mathbf{1}$ is an eigenvector corresponding to $\lambda_j,$ then the eigenvalues and eigenvectors of $G\overline{G}$ are: 
% \begin{enumerate}
%         \item $\alpha_{1,2}(\lambda_j)=\dfrac{n\pm\sqrt{\ob{n-2(\lambda_j+1)}^2+4}}{2}$ with multiplicity $m_j,$ for $2\leq j\leq d$ and eigenvectors $\begin{bmatrix}
%             \x_j\\
%             \ob{\alpha_{1,2}(\lambda_j)-(\lambda_j+1)}\x_j
%         \end{bmatrix}$ that are orthogonal to the all one vector.
        
%         % \item $\lambda_{2,i}(\gamma_i)=\dfrac{n-\sqrt{\ob{n-2(\gamma_i+1)}^2+4}}{2}$ with multiplicity $m_i,$ for $2\leq i\leq d$ and eigenvector $\begin{bmatrix}
%         %     u\\
%         %     \ob{\alpha_{2,i}(\gamma_i)-(\gamma_i+1)}u
%         % \end{bmatrix};$
        
%         \item$\beta_{1,2}(n,k)=n\pm\sqrt{(n-1-2k)^2+1}$ with eigenvectors $\begin{bmatrix}
%             \mathbf{1}\\
%             \ob{\beta_{1,2}(n,k)-2k-1}\mathbf{1}
%         \end{bmatrix}.$ $\beta_{1,2}(n,k)$ are the two main eigenvalues when $G\overline{G}$ is non-regular.

% \end{enumerate}
% \end{thm}
% By Theorem \ref{5t5} and an argument analogous to that used for the adjacency matrix in Lemma \ref{5t1}, we obtain the following result.
% % By using Theorem \ref{5t5}, and a similar approach as in Theorem \ref{5t1} and \ref{5t2}, we obtain the following result.
% \begin{thm}\label{5t6}
%      Let $\u\in\mathbb{R}^n$ be a fixed state in a regular graph $G$ such that $\mathbf{1}^T\u=0.$ Then the complementary prism graph $G\overline{G}$  exhibits fractional revival from the state $\ob{\u^T,0}^T$ to $\ob{0,\u^T}^T$ for all time $t$ with respect to the signless Laplacian matrix.
% \end{thm}
%  Theorem \ref{5t2}, Theorem \ref{5t4} and Theorem \ref{5t6} together yield the following result, which holds for the adjacency, Laplacian, and signless Laplacian matrices whenever 
% $G$ is regular, even though $G\overline{G}$ need not be regular.
% \begin{thm}\label{5t7}
%     Let $\u\in\mathbb{R}^n$ be a fixed state in a regular graph $G$ such that $\mathbf{1}^T\u=0.$ Then the complementary prism graph $G\overline{G}$  exhibits fractional revival from the state $\ob{\u^T,\mathbf{0}}^T$ to $\ob{\mathbf{0},\u^T}^T$ for all time $t$ with respect to the adjacency, Laplacian, and signless Laplacian matrices.
% \end{thm}

%  \begin{figure}
% \begin{center}
% \begin{tikzpicture}[scale=.3,auto=left]
%                        \tikzstyle{every node}=[circle, thick, fill=white, scale=0.6]
                       
% 		        \node[draw] (1) at (-4,4) {$a$};		        
% 		        \node[draw,minimum size=0.65cm, inner sep=0 pt] (2) at (-4, -4) {};
% 		        \node[draw,minimum size=0.65cm, inner sep=0 pt] (3) at (4, -4) {$b$};
% 		        \node[draw,minimum size=0.65cm, inner sep=0 pt] (4) at (4, 4) {};
% 		        \node[draw,minimum size=0.65cm, inner sep=0 pt] (5) at (-1.5, 1.5) {$a'$};		 \node[draw,minimum size=0.65cm, inner sep=0 pt] (6) at (-1.5, -1.5) {};       
% 	\node[draw,minimum size=0.65cm, inner sep=0 pt] (7) at (1.5, -1.5) {$b'$};
%     \node[draw,minimum size=0.65cm, inner sep=0 pt] (8) at (1.5, 1.5) {};
				
								
% 				\draw [thick, black!70] (1)--(2)--(3)--(4)--(1);
				
%                 \draw [thick, black!70] (1)--(5)--(7)--(3);
% \draw [thick, black!70] (2)--(6)--(8)--(4);
		  
% 				\end{tikzpicture}
				
% \end{center}	
% \caption{\label{5fig1} The complementary prism $C_4\overline{C_4}.$}
% \end{figure}
% \begin{exm}
% Let $a$ and $b$ be a pair of antipodal vertices in $C_4.$ Then the state $\frac{1}{\sqrt{2}}(\e_a-\e_b)$ is a fixed state in $C_4.$ Therefore, by Theorem \ref{5t7} the complementary prism $C_4\overline{C_4}$ exhibits FR for all time $t$ from $\frac{1}{\sqrt{2}}(\e_a-\e_b)$ to $\frac{1}{\sqrt{2}}(\e_{a'}-\e_{b'})$ [see Figure \ref{5fig1}] with respect to the adjacency, Laplacian, and signless Laplacian matrices. In particular, at $t=\frac{\pi}{2},$ $C_4\overline{C_4}$ admits PST between $\frac{1}{\sqrt{2}}(\e_a-\e_b)$ and $\frac{1}{\sqrt{2}}(\e_{a'}-\e_{b'})$ with respect to the Laplacian matrix.
% \end{exm}
% A pair of vertices $a$ and $b$ in $G$ are called twin vertices if $N_G(a)\backslash\{b\}=N_G(b)\backslash\{a\},$ where $N_G(a)$ denote the set of vertices that are adjacent to vertex $a$ in $G.$ It follows that, the state $\frac{1}{\sqrt{2}}(\e_a-\e_b)$ is a fixed state in $G$ \cite{god7}. By Theorem \ref{5t7}, if $a$ and $b$ are twin vertices in a graph $G,$ then the complementary prism $G\overline{G}$ exhibits fractional revival from $\frac{1}{\sqrt{2}}(\e_a-\e_b)$ for all time $t$ with respect to the Laplacian matrix. Moreover, if $G$ is regular, then fractional revival occurs relative to the adjacency, Laplacian and signless Laplacian matrices.   

% \begin{exm}\label{5ex3}
%     Let $K_n$ be a complete graph on $n$ vertices. Since any two distinct vertices $a$ and $b$ in $K_n$ are twin vertices, the state $\frac{1}{\sqrt{2}}(\e_a-\e_b)$ is a fixed state in $K_n.$ The corona product of $K_n$ and $K_1$ is the complementary prism $K_n\overline{K_n}.$ Therefore, by Theorem \ref{5t7}, the corona product of $K_n$ and $K_1$ exhibits fractional revival from the state $\frac{1}{\sqrt{2}}(\e_a-\e_b)$ for all time $t$ with respect to the adjacency, Laplacian, and signless Laplacian matrices.
% \end{exm}
% \begin{exm}
%   Let $K_{m,n}$ be a complete bipartite graph, and $a$ and $b$ be any two vertices in the same partite set. Then $\frac{1}{\sqrt{2}}(\e_a-\e_b)$ becomes a fixed state in $K_{m,n}.$ Hence, by Theorem \ref{5t4}, the complementary prism $K_{m,n}\overline{K_{m,n}}$  exhibits fractional revival from the state $\frac{1}{\sqrt{2}}(\e_a-\e_b)$ to $\frac{1}{\sqrt{2}}(\e_{a'}-\e_{b'})$ for all time $t$ with respect to the Laplacian matrix. In particular, if $m=n,$ then fractional revival occurs relative to the adjacency, Laplacian, and signless Laplacian matrices.
% \end{exm}

% % It is worth noting that the complement of $G\overline{G}$ need not itself be the complementary prism of any graph. By \cite[Lemma 3]{pal11}, one may observe that if $\u$ is a fixed state in $G$ such that $\mathbf{1}^T\u=0,$ then the complement of the complementary prism graph $G\overline{G}$  exhibits fractional revival from the state $\ob{\u^T,\mathbf{0}}^T$ to $\ob{\mathbf{0},\u^T}^T$ for all time $t$ with respect to the Laplacian matrix. Here, we observe that the complement of the complementary prism graph $G\overline{G}$ also preserves the fractional revival property whenever $G$ is regular, relative to the adjacency and signless Laplacian matrices. Let $G$ be a graph on $n$ vertices, and
% % $H=\overline{G\overline{G}}$ be the complement of the complementary prism graph. Since the number of vertices in $G\overline{G}$ is $2n,$ we have  $A(H)=J-I-A(G\overline{G}),$ and $D(H)=(2n-1)I-D(G\overline{G}).$ Therefore, for $M\in\{A,Q\},$
% %  \[M(H)=\zeta I+J-M(G\overline{G}),\]
% %     where $\zeta=-1$ and $2n-2$ for $M=A$ and $Q,$ respectively. Including the observation for the Laplacian matrix, we have the following result.
%   %   $\zeta_M=\left\{ \begin{array}{rcl}
%   % -1, & \mbox{if} & M ~\text{is the adjacency matrix}, \\ 
%   % 2n, & \mbox{if} & M ~\text{is the Laplacian matrix}, \\
%   %  2n-2, & \mbox{if} & M ~\text{is the signless Laplacian matrix},
%   % \end{array}\right.$

% %   and

% %  \vspace{0.5cm}
% %  \hspace{0.3cm} $\sigma_M=\left\{ \begin{array}{rcl}
% %   -1, & \mbox{if} & M ~\text{is the Laplacian matrix}, \\ 
% %   1, & \mbox{if} & M ~\text{is the adjacency or signless Laplacian matrix}. 
% %   \end{array}\right.$
% %  \\
% % It is worth noting that the complement of $G\overline{G}$ need not itself be the complementary prism of any graph. However, the complement of  $G\overline{G}$ still exhibits fractional revival between the states $\ob{\psi^T,0}^T$ and $\ob{0,\psi^T}^T$ for all time $t$ with respect to adjacency, Laplacian, and signless Laplacian matrix.
% % \begin{thm}\label{5t8}
% %  Let $\u$ be a fixed state in a regular graph $G$ such that $\mathbf{1}^T\u=0.$ Then the complement of a complementary prism graph $G\overline{G}$  exhibits fractional revival from the state $\ob{\u^T,\mathbf{0}}^T$ to $\ob{\mathbf{0},\u^T}^T$ for all time $t$ with respect to the adjacency, Laplacian and signless Laplacian matrix.    
% % \end{thm}
% %   \begin{proof}
% %   From \cite[Lemma 3]{pal11}, the result follows for the Laplacian matrix. 
% % Let $\u$ be a fixed state in $G$ such that $\mathbf{1}^T\u=0.$ Then by Theorem \ref{5t7}, for each time $t$ there is some $\alpha(t)$ and $\beta(t)$ with respect to adjacency and signless Laplacian matrix such that 
% %     \begin{equation}\label{5e4}
% %      U_{M(G\overline{G})}(t)(\u^T, \mathbf{0})^T=\alpha(t)(\u^T, \mathbf{0})^T+\beta(t)(\mathbf{0},\u^T)^T.
% %      \end{equation}
% % %     where 
% % % $|\alpha(t)|^2+|\beta(t)|^2=1.$ 
% % Let $H=\overline{G\overline{G}}$ be the complement of the complementary prism graph. 
% % Since $\mathbf{1}^T\u=0,$ we have $M(H)(\u^T,\mathbf{0})^T=(\zeta I-M(G\overline{G}))(\u^T,\mathbf{0})^T.$ Therefore, for all $t\in\mathbb{R},$
% %   \begin{equation}\label{5e5}
% %       U_{M(H)}(t)(\psi^T,0)^T=\exp{(i\zeta_M t)}U_{M(G\overline{G})}(-t)(\psi^T,0)^T.
% %  \end{equation}
% %   Now, \eqref{5e4} and \eqref{5e5} gives 
% %   \[U_{M(H)}(t)(\u^T,\mathbf{0})^T=\exp{(i\zeta t)}\tb{\overline{\alpha(t)}(\u^T, \mathbf{0})^T+\overline{\beta(t)}(\mathbf{0},\u^T)^T}.\]
% %   Since $|\exp{(i\zeta_M t)}\overline{\alpha(t)}|^2+|\exp{(i\zeta_M t)}\overline{\beta(t)}|^2=1,$ for all $t\in\mathbb{R},$ the result follows.
% %  \end{proof}
% % \subsection{Complete bipartite graph}
% % \begin{lem}
% %    Let $a$ and $b$ be two vertices in the same partite set of a complete bipartite graph $K_{m,n}.$ The two states $\frac{1}{\sqrt{2}}(\e_a-\e_b)$ and $\frac{1}{\sqrt{2}}(\e_{a'}-\e_{b'})$ in the complementary prism $K_{m,n}\overline{K_{m,n}}$ are strongly cospectral with respect to the Laplacian matrix if and only if $m=n.$ 
% % \end{lem}
% % \begin{proof}
    
% % \end{proof}
% % \begin{thm}
% %     Let $K_{m,n}$ be a complete bipartite graph. Suppose $m=n$ and let $a$ and $b$ be two vertices in the same partite set of $K_{m,n}$, then the complementary prism $K_{m,n}\overline{K_{m,n}}$ exhibits perfect state transfer between $\frac{1}{\sqrt{2}}(\e_a-\e_b)$ and $\frac{1}{\sqrt{2}}(\e_{a'}-\e_{b'})$ with respect to the Laplacian matrix.
% % \end{thm}
% % \begin{proof}
    
% % \end{proof}
% % This gives an infinite family of graphs having perfect state transfer between an edge pair and non-edge pair.
% % \subsection{FR in graphs with generalized cluster}
% % \begin{defn}\cite{pal11}
% %    Let $G$ be an undirected graph having two subsets of vertices $C$ and $S$ where
% %  $|C| \geq 2.$ The vertices in $C$ are pairwise non-adjacent and share the same neighborhood $S.$
% %  For each $v \in S,$ the edge weight $w(u,v)$ is constant for all $u \in C.$ Then $(C,S)$ is called a generalized cluster in $G$ of order $|C|$ and degree $|S|.$ 
% % \end{defn}

% % \begin{defn}\cite{pal11}
% %     Let $G$ be a connected graph of order n containing a generalized cluster
% %  $(C,S),$ and let $H$ be a graph with vertex set $C.$ Then the graph $G(H)$ is defined to have vertex set $V(G)$ and edge set $E(G)\cup E(H),$ where the edge weights from both $G$ and $H$
% %  are preserved.
% % \end{defn}
% % \begin{thm}\label{5t9}
% %     Let $\psi \in \mathbb{R}^c$ satisfying $\bf{1}_c^T\psi = 0.$ Suppose $H$ is a graph for which $\bf{1}_c$ is an
% %  eigenvector of $M(H).$ The graph $H$ exhibits fractional revival between $\psi$ and $\phi$ if and only if $G(H)$
% %  exhibits fractional revival between $(\psi^T,0)^T$ and $(\phi^T,0)^T.$
% % \end{thm}
% % \begin{figure}
% % \begin{center}
% % \begin{tikzpicture}[scale=.75,auto=left]
% %                        \tikzstyle{every node}=[circle, thick, fill=white, scale=0.5]
                       
% % 		        \node[draw] (1) at (-1,1) {$1$};		        
% % 		        \node[draw,minimum size=0.7cm, inner sep=0 pt] (2) at (0, 0) {$2$};
% % 		        \node[draw,minimum size=0.7cm, inner sep=0 pt] (3) at (-1, -1) {$3$};
% % 		        \node[draw,minimum size=0.7cm, inner sep=0 pt] (4) at (-2, 2) {$4$};
% % 		        \node[draw,minimum size=0.7cm, inner sep=0 pt] (5) at (-2, -2) {$5$};	
% %                 \node[draw,minimum size=0.7cm, inner sep=0 pt] (6) at (1.5, 0.3) {$6$};
% %                 \node[draw,minimum size=0.7cm, inner sep=0 pt] (7) at (4.5, 0) {$7$};
% %                 \node[draw,minimum size=0.7cm, inner sep=0 pt] (8) at (5.5, 0) {};
% %                 \node[draw,minimum size=0.7cm, inner sep=0 pt] (9) at (6.5, 0) {};
% %                 \node[draw,minimum size=0.7cm, inner sep=0 pt] (10) at (9.5, 0) {};
% % 				\node[draw,minimum size=0.7cm, inner sep=0 pt] (11) at (10.5, 0) {};
% %                 \node at (7.5,0) {$\bullet$};
% % \node at (8.0,0) {$\bullet$};
% % \node at (8.5,0) {$\bullet$};
% % \draw [thick, black!70] (5)--(4)--(1)--(2)--(3)--(5);
% % \draw [thick, black!70] (4)--(7)--(6)--(2);
% % \draw [thick, black!70] (1)--(7)--(2);
% % \draw [thick, black!70] (3)--(7)--(5);	\draw [thick, black!70] (7)--(8)--(9);
% % \draw [thick, black!70] (10)--(11);

		  
% % 				\end{tikzpicture}
				
% % \end{center}	
% % \caption{\label{5fig2} Another signed version of the graph on the left in Figure \ref{fig1}}
% % \end{figure}


% % The above result can be used to construct a variety of graphs that exhibit fractional revival. In
% %  particular, if $\psi$ is a pair state in $H,$ then it satisfies $\bf{1}_c^T\psi= 0.$ Moreover, if $H$ has fractional revival then so does the graph $G(H).$ In Figure \ref{5fig2}, the subgraph H induced by
% %  the vertices 1,2,3,4,5 and 6 in the graph $G$ is a complementary prism $P_3\overline{P_3}$ that exhibits fractional revival between $\e_1-\e_3$ and $\e_4-\e_5$ at all time $t$ with respect to adjacency, Laplacian, and signless Laplacian matrix.  Since $\{V(H),{7}\}$ forms a cluster in $G-E(H),$ the graph $G$ admits fractional revival 
% %   between $((\e_1-\e_3)^T,0)^T$ and $((\e_4-\e_5)^T,0)^T$ irrespective of the choice of matrix $M(G),$ whether
% %  it be the adjacency matrix, Laplacian matrix, or the signless Laplacian matrix. This
% %  provides an infinite family of graphs that exhibit fractional revival with respect to each of the aforementioned matrices, between the same pair
% %  states and at all time despite the fact that these graphs are not regular. 
% % \subsection{Cartesian product}
% %  The Cartesian product of two graphs $G$ and $H,$ denoted by $G\square H,$ is defined as the graph with vertex set $V (G)\times V (H),$ where two vertices $(a_1,a_2)$ and $(b_1,b_2)$ are adjacent in $G\square H,$ if and only if either $a_1 = b_1$ and $(a_2,b_2)\in E(H),$ or $a_2 =b_2$ and $(a_1,b_1) \in E(G).$ If the graphs $G$ and $H$ have $m$ and $n$ vertices, respectively, then for $M\in\{A,L,Q\},$ the following holds.
% %  \[M(G\square H)=M(G)\otimes I_n+I_m\otimes M(H).\]
% %  Since $M$ $(G)\otimes I_n$ and $I_m\otimes M(H)$ commute, the transition matrix of $G\square H$ is given by
% %  \begin{equation}\label{5e6}
% %      U_{G\square H}(t) = U_G
% %  (t) \otimes U_H(t).
% %  \end{equation}
% %  The existence of FR from vertex state relative to adjacency and Laplacian matrices is studied in \cite{cha2, chan3}. Here, we characterize FR from real pure states. A similar characterization for PST from real pure states was obtained in \cite[Theorem 8.1]{god7}. 
% % \begin{thm}\label{5t10}
% %     Let $G$ and $H$ be graphs with the adjacency, Laplacian, or signless Laplacian matrices, respectively. Let $\u_1$ and $\u_2$ be states in $G$ and $\v_1$ and $v_2$ be states in $H.$ Then fractional revival occurs from $\u_1\otimes \v_1$ to $\u_2\otimes \v_2$ in $G\square H$ at $\tau$ if one of the followings hold:
% %     \begin{enumerate}
% %         \item fractional revival between $\u_1$ and $\u_2$ in $G$ and fractional revival between $\v_1$ and $\v_2$ in $H$ both at the same time $\tau.$
% %         \item fractional revival occurs between $\u_1$ and $\u_2$ at time $\tau$ and $\v_1=\v_2$ is periodic in $H$ at the same time $\tau.$
% %     \end{enumerate}
    
% % \end{thm}

% % % \begin{exm}
% % %   Let $G$ be a graph with fixed state $\u$ with $\bf{1}^T\u=0.$ Then by Theorem \ref{5t7} the complementary prism $G\overline G$ exhibits fractional revival between $(\u^T,\mathbf{0})^T$ and $(\mathbf{0},\u^T)^T$ for all time $t$ with respect to adjacency, Laplacian, and signless Laplacian matrix. Therefore, by Theorem \ref{5t10}, we have the cartesian product $G\overline{G}^{\square m}$ exhibits fractional revival between $((\u^T,\mathbf{0})^T)^{\otimes m}$ and $((\mathbf{0},\u^T)^T)^{\otimes m}$ for all time $t.$
% % % \end{exm}
% % % Next we consider $G$ as a complete graph and have the following observation.
% % \begin{exm}
% %   In Example \ref{5ex3}, we find that for any two vertices $a$ and $b$ in $K_n,$ the graph $K_n\overline{K_n}$ exhibits FR from the state $\frac{1}{\sqrt{2}}\ob{\e_a-\e_b}$ for all time $t$ with respect to adjacency, Laplacian, and signless Laplacian matrices. Therefore, by Theorem \ref{5t10}, we have the cartesian product $(K_n\overline{K_n})^{\square m}$ exhibits fractional revival from $((\frac{1}{\sqrt{2}}\ob{\e_a-\e_b}^T,\mathbf{0})^T)^{\otimes m}$ for all time $t.$ Moreover, if $\mathcal{M}$ is a perfect matching in $K_{2n},$ then the state $\frac{1}{\sqrt{2}}\ob{\e_a-\e_b}$ corresponding to an edge $(a,b)\in {\mathcal{M}}$ is a fixed state in the resulting graph $K_{2n}\backslash{\mathcal{M}}.$ Hence by Theorem \ref{5t10}, the cartesian product $((K_{2n}\backslash{M})\overline{K_{2n}\backslash{\mathcal{M}}})^{\square m}$ exhibits FR from $((\frac{1}{\sqrt{2}}\ob{\e_a-\e_b}^T,\mathbf{0})^T)^{\otimes m}$ for all time $t$ with respect to $M.$
% %  \end{exm}
 
% % % \begin{exm}
% % %   Let $G$ be a graph with fixed state $\u$ with $\bf{1}^T\u=0.$ Then by Theorem \ref{5t7} the complementary prism $G\overline G$ exhibits fractional revival between $(\u^T,\mathbf{0})^T$ and $(\mathbf{0},\u^T)^T$ for all time $t$ with respect to adjacency, Laplacian, and signless Laplacian matrix. Then, by Theorem \ref{5t10}, the cartesian product $G\overline{G}\square G$  exhibits fractional revival between $(\u^T,\mathbf{0})^T\otimes\u$ and $(\mathbf{0},\u^T)^T\otimes \u$ for all time $t.$ Moreover, if $\v$ be any fixed state in $G,$ $G\overline{G}\square G$  exhibits fractional revival between $(\u^T,\mathbf{0})^T\otimes\v$ and $(\mathbf{0},\u^T)^T\otimes \v$ for all time $t$ relative to $M.$
% % % \end{exm}
 
 
 
 \section*{Disclosure statement}
  No potential conflict of interest was reported by the author(s).


 
\section*{Acknowledgements}
 We sincerely thank the reviewers for their insightful comments and valuable suggestions to improve the manuscript. 
% H. Pal is funded by the Science and Engineering Research Board (Project: SRG/2021/000522).
S. Mohapatra is supported by the Department of Science and Technology (INSPIRE: IF210209).

%%%%~~Bibliography~~%%%%%%

\bibliographystyle{abbrv}
\bibliography{references}



\end{document}

